%% Vectors I HW1 (Introduction to Vectors in Three Dimensions): HOUSE FORMAT (house_format/SPEC.md), converted 2026-11-09.
%% Questions, marks, diagrams (house + frustum, unchanged TikZ) and solutions copied word for word from
%% build_hw1.py (pre-house copy: build_hw1.pre_house.py; backups in house_format/phase2/vec_backup/hw1).
%% The same file feeds the worked solutions (Vectors_I_HW1_Solutions.tex).
\documentclass[a4paper,11pt,oneside]{book}
\providecommand\housemode{hw}
% =====================================================================
% house.tex: ONE preamble for every Maximus Pung Jun H2 Math document,
% modelled on Ch10 Definite Integrals + Ch11 Differential Equations (see SPEC.md).
%
% Before % =====================================================================
% house.tex: ONE preamble for every Maximus Pung Jun H2 Math document,
% modelled on Ch10 Definite Integrals + Ch11 Differential Equations (see SPEC.md).
%
% Before % =====================================================================
% house.tex: ONE preamble for every Maximus Pung Jun H2 Math document,
% modelled on Ch10 Definite Integrals + Ch11 Differential Equations (see SPEC.md).
%
% Before \input{house}, set:
%   \def\housemode{notes}    lecture notes: book class, fncychap Lenny chapter page, boxes
%   \def\housemode{hw}       homework / additional / workshop QUESTIONS (+ Answers section)
%   \def\housemode{sol}      worked SOLUTIONS of a hw (RI "Solution | Comments" table)
%   \def\housemode{test}     test question paper (prelim cover, prelim layout)
%   \def\housemode{testsol}  test worked solutions (prelim layout + RI table)
%   \def\housechapterno{10}  (notes: chapter number; optional)
%   \def\housenotesLM{1}     (notes only, optional) keep Latin Modern as in the current
%                            DI/DE notes instead of Times New Roman + TeX Gyre Termes Math
% then \housechapter{Definite Integrals} and \ditle{Homework 1}.
% Document class (put BEFORE \input{house}):
%   notes: \documentclass[a4paper]{book}
%   others: \documentclass[a4paper,11pt,oneside]{book}
% A HW file uses \providecommand\housemode{hw}; its solutions file has NO \documentclass:
%   \def\housemode{sol}\def\diquestionsonly{1}\input{HW1}  (as DI/DE HWn_solutions.tex)
% Compile: TEXINPUTS=.:/workspace/house_format: xelatex (2 runs; solutions until .aux stable)
% =====================================================================
\providecommand\housemode{hw}
\providecommand\housedir{/workspace/house_format}
\newcommand\housechap{}
\newcommand\housechapter[1]{\renewcommand\housechap{#1}}
\def\housexnotes{notes}\def\housexhw{hw}\def\housexsol{sol}\def\housextest{test}\def\housextestsol{testsol}
\edef\housexm{\housemode}
\input{house_base}
\input{house_local}
\ifx\housexm\housexnotes
  \ifdefined\housenotesLM\else\input{house_fonts}\fi
  \input{house_tf}
  \providecommand{\units}[1]{\ \text{units}^{#1}}
  \raggedbottom
\else
  \input{house_fonts}
  \input{house_hw}
  \input{house_solmacros}
  \input{house_tf}
  \ifx\housexm\housexsol \def\prelimkeepcolour{1}\input{house_prelim}\input{house_sol}\tfanswerstrue\fi
  \ifx\housexm\housextest \input{house_prelim}\fi
  \ifx\housexm\housextestsol \def\prelimkeepcolour{1}\input{house_prelim}\input{house_sol}\tfanswerstrue\fi
\fi
, set:
%   \def\housemode{notes}    lecture notes: book class, fncychap Lenny chapter page, boxes
%   \def\housemode{hw}       homework / additional / workshop QUESTIONS (+ Answers section)
%   \def\housemode{sol}      worked SOLUTIONS of a hw (RI "Solution | Comments" table)
%   \def\housemode{test}     test question paper (prelim cover, prelim layout)
%   \def\housemode{testsol}  test worked solutions (prelim layout + RI table)
%   \def\housechapterno{10}  (notes: chapter number; optional)
%   \def\housenotesLM{1}     (notes only, optional) keep Latin Modern as in the current
%                            DI/DE notes instead of Times New Roman + TeX Gyre Termes Math
% then \housechapter{Definite Integrals} and \ditle{Homework 1}.
% Document class (put BEFORE % =====================================================================
% house.tex: ONE preamble for every Maximus Pung Jun H2 Math document,
% modelled on Ch10 Definite Integrals + Ch11 Differential Equations (see SPEC.md).
%
% Before \input{house}, set:
%   \def\housemode{notes}    lecture notes: book class, fncychap Lenny chapter page, boxes
%   \def\housemode{hw}       homework / additional / workshop QUESTIONS (+ Answers section)
%   \def\housemode{sol}      worked SOLUTIONS of a hw (RI "Solution | Comments" table)
%   \def\housemode{test}     test question paper (prelim cover, prelim layout)
%   \def\housemode{testsol}  test worked solutions (prelim layout + RI table)
%   \def\housechapterno{10}  (notes: chapter number; optional)
%   \def\housenotesLM{1}     (notes only, optional) keep Latin Modern as in the current
%                            DI/DE notes instead of Times New Roman + TeX Gyre Termes Math
% then \housechapter{Definite Integrals} and \ditle{Homework 1}.
% Document class (put BEFORE \input{house}):
%   notes: \documentclass[a4paper]{book}
%   others: \documentclass[a4paper,11pt,oneside]{book}
% A HW file uses \providecommand\housemode{hw}; its solutions file has NO \documentclass:
%   \def\housemode{sol}\def\diquestionsonly{1}\input{HW1}  (as DI/DE HWn_solutions.tex)
% Compile: TEXINPUTS=.:/workspace/house_format: xelatex (2 runs; solutions until .aux stable)
% =====================================================================
\providecommand\housemode{hw}
\providecommand\housedir{/workspace/house_format}
\newcommand\housechap{}
\newcommand\housechapter[1]{\renewcommand\housechap{#1}}
\def\housexnotes{notes}\def\housexhw{hw}\def\housexsol{sol}\def\housextest{test}\def\housextestsol{testsol}
\edef\housexm{\housemode}
\input{house_base}
\input{house_local}
\ifx\housexm\housexnotes
  \ifdefined\housenotesLM\else\input{house_fonts}\fi
  \input{house_tf}
  \providecommand{\units}[1]{\ \text{units}^{#1}}
  \raggedbottom
\else
  \input{house_fonts}
  \input{house_hw}
  \input{house_solmacros}
  \input{house_tf}
  \ifx\housexm\housexsol \def\prelimkeepcolour{1}\input{house_prelim}\input{house_sol}\tfanswerstrue\fi
  \ifx\housexm\housextest \input{house_prelim}\fi
  \ifx\housexm\housextestsol \def\prelimkeepcolour{1}\input{house_prelim}\input{house_sol}\tfanswerstrue\fi
\fi
):
%   notes: \documentclass[a4paper]{book}
%   others: \documentclass[a4paper,11pt,oneside]{book}
% A HW file uses \providecommand\housemode{hw}; its solutions file has NO \documentclass:
%   \def\housemode{sol}\def\diquestionsonly{1}\input{HW1}  (as DI/DE HWn_solutions.tex)
% Compile: TEXINPUTS=.:/workspace/house_format: xelatex (2 runs; solutions until .aux stable)
% =====================================================================
\providecommand\housemode{hw}
\providecommand\housedir{/workspace/house_format}
\newcommand\housechap{}
\newcommand\housechapter[1]{\renewcommand\housechap{#1}}
\def\housexnotes{notes}\def\housexhw{hw}\def\housexsol{sol}\def\housextest{test}\def\housextestsol{testsol}
\edef\housexm{\housemode}
% house_base.tex: HOUSE COPY of definite_integrals/common/ws_preamble_build.tex
% (= Maximus's own preamble.tex + the DE v3 / AoI additions). Only change: GC font path. [HOUSE]
% WORKSHEET COPY of ../notes/preamble_build.tex (2026-10-09 version), unchanged.
% preamble_build.tex: copy of Maximus's preamble.tex for the DE v3 build.
% The blocks marked [BUILD COPY CHANGE] are the earlier build fixes.
% The blocks marked [V3 ADDITION] are new in this version (example box,
% spacing). Everything else is unchanged from preamble.tex.

% For margins
\usepackage[a4paper, margin = 20mm]{geometry}

\usepackage[Lenny]{fncychap}
% Rejne, Glenn, Lenny

\usepackage{etoc}

% [V3 ADDITION] cancel: used to show the dx's "cancelling" in Section 11.3
\usepackage{cancel}
\renewcommand{\CancelColor}{\color{red}}

% For enumeration & math symbols
\usepackage{enumerate}
\usepackage{amsmath,amsfonts,amssymb,amsthm,graphicx,lipsum}

% For diagrams, plots and boxes
\usepackage{pgfplots}
\pgfplotsset{compat=1.18}
\usepackage{tikz}
\usepackage[most]{tcolorbox}
\tcbuselibrary{skins}
\usepackage{xcolor}
\pgfplotsset{
    standard/.style={clip = false,
    axis line style = thick,
    trig format=rad,
    axis lines = middle,
    every axis x label/.style={at={(current axis.right of origin)},anchor= west},
    every axis y label/.style={at={(current axis.above origin)},anchor=south}
    }
}

% For GC font
\usepackage{fontspec}
% [BUILD COPY CHANGE 1] The TI-84 key font is loaded only if the file exists
% (checked in ./Fonts/... and ../Fonts/...). Otherwise \ti{} falls back to a
% small boxed sans-serif key. Known letters of the TI font are mapped to their
% key names in the fallback: o = [y=], z = [mode], f = [L3].
\newcommand{\tikey}[1]{\tcbox[on line, enhanced, boxsep=0pt, left=2pt, right=2pt,
  top=1pt, bottom=1pt, boxrule=0.5pt, arc=1.5pt, colback=white,
  colframe=black!70]{\sffamily\footnotesize\strut #1}}
\ExplSyntaxOn
\NewDocumentCommand{\tifallback}{m}{
  \str_case:nnF {#1}
    { {o}{\tikey{y=}} {z}{\tikey{mode}} {f}{\tikey{L3}} }
    { \tikey{#1} }
}
\ExplSyntaxOff
\def\tifontcase{0}
\IfFileExists{Fonts/TI-84 Plus CE Key Fonts/TI84PlusCEKeys.ttf}
  {\def\tifontcase{1}}{}
\IfFileExists{../Fonts/TI-84 Plus CE Key Fonts/TI84PlusCEKeys.ttf}
  {\def\tifontcase{2}}{}
\ifcase\tifontcase\relax
  % 0: font not found, use boxed fallback keys
  \newcommand{\ti}[1]{\tifallback{#1}}
\or
  % 1: original location (as in preamble.tex)
  \newfontfamily\tifont[
    Path = {Fonts/TI-84 Plus CE Key Fonts/},
    UprightFont = {TI84PlusCEKeys.ttf}
  ]{}
  \newcommand{\ti}[1]{{\tifont #1}}
\or
  % 2: one level up (when compiling inside draft-ch1/)
  \newfontfamily\tifont[
    Path = {../Fonts/TI-84 Plus CE Key Fonts/},
    UprightFont = {TI84PlusCEKeys.ttf}
  ]{}
  \newcommand{\ti}[1]{{\tifont #1}}
\fi

% [AoI items 17/18/37] GC DISPLAY FONT. Every GC key, command, menu item and
% screen line goes through \gcfont, so the font can be swapped HERE in one place.
%   \gc{fnInt(}      command / menu item / stored variable in the display font
%   \gckey{2nd}      a calculator key (display font in a small key outline)
%   \gcscreen{...}   a calculator screen (display font on a grey-green LCD panel)
% The font file lives in ../Fonts (shared by notes, slides and worksheet).
\def\gcfontfile{VT323-Regular.ttf}
% [HOUSE] GC font: always the absolute copy in \housedir/Fonts (relative Path= breaks under TEXINPUTS)
\providecommand\housedir{/workspace/house_format}
\IfFileExists{\housedir/Fonts/\gcfontfile}{\newfontfamily\gcfont[Path=\housedir/Fonts/, Scale=1.32]{\gcfontfile}}{%
  \newcommand\gcfont{\ttfamily}}
\newcommand{\gc}[1]{\mbox{\gcfont #1}}
\newcommand{\gckey}[1]{\tcbox[on line, enhanced, boxsep=0pt, left=2pt, right=2pt,
  top=0.6pt, bottom=0.2pt, boxrule=0.5pt, arc=1.5pt, colback=black!4,
  colframe=black!65]{\gcfont\strut #1}}
\definecolor{lcd}{RGB}{214,222,204}
\newcommand{\gcscreen}[1]{\tcbox[on line, enhanced, boxsep=1pt, left=4pt, right=4pt,
  top=2pt, bottom=2pt, boxrule=1.2pt, arc=1pt, colback=lcd, colframe=black!70]{%
  \gcfont\begin{tabular}{@{}l@{}}#1\end{tabular}}}
% integral template as shown on the TI-84 Plus CE (MathPrint): \gcint{0}{2}{e^(-X²)}{X}
\newcommand{\gcint}[4]{{\gcfont\raisebox{-0.35ex}{\Large ∫}\hspace{-0.15em}%
  \raisebox{-0.2ex}{\begin{tabular}{@{}l@{}}\footnotesize #2\\[-1.3ex]\footnotesize #1\end{tabular}}%
  \hspace{0.1em}(#3)\mathrm{d}#4}}
% old key macro kept as an alias so nothing breaks
\renewcommand{\ti}[1]{\gckey{#1}}

% SI units
\usepackage{siunitx}

% setspace
%\usepackage{setspace}
%\onehalfspacing

% For hyperlinks within and out of document
\usepackage{hyperref}
\hypersetup{
    colorlinks = true,
    linkcolor = black,
    urlcolor = blue,
}   

% For captions for pics
\usepackage{capt-of}


% tcolorbox custom environments
% Definition ========================================================
\newtcolorbox[auto counter,number within=section]{definition}[2][drop shadow]{%
enhanced, before skip=5mm, after skip=6mm, % [V3 ADDITION] spacing
colback=blue!5!white,colframe=blue!50!black,fonttitle=\bfseries,
title= Definition~\thetcbcounter. #2,#1}

% Theorem ==========================================================
\newtcolorbox[auto counter,number within=section]{theorem}[2][drop shadow]{%
enhanced, before skip=5mm, after skip=6mm, % [V3 ADDITION] spacing
colback=black!5!white,colframe=black,fonttitle=\bfseries,
title= Theorem~\thetcbcounter. #2,#1}

% Caution box ======================================================
\newtcolorbox{caution}[1][]{enhanced,
before skip=5mm,after skip=6mm, % [V3 ADDITION] was 2mm/3mm
boxrule=1pt,left=11mm,right=2mm,top=3mm,bottom=3mm,
colback=yellow!50,
colframe=yellow!20!black,
sharp corners,rounded corners=southeast,arc is angular,arc=3mm,
underlay={%
\path[fill=tcbcolback!80!black] ([yshift=3mm]interior.south east)--++(-0.4,-0.1)--++(0.1,-0.2);
\path[draw=tcbcolframe,shorten <=-0.05mm,shorten >=-0.05mm] ([yshift=3mm]interior.south east)--++(-0.4,-0.1)--++(0.1,-0.2);
\path[fill=yellow!50!black,draw=none] (interior.south west) rectangle node[white]{\Huge\bfseries !} ([xshift=9mm]interior.north west);
},
drop fuzzy shadow,#1}

% Question box ======================================================
\newtcolorbox{question}[1][]{enhanced,
before skip=5mm,after skip=6mm, % [V3 ADDITION] was 2mm/3mm
boxrule=1pt,left=11mm,right=2mm,top=3mm,bottom=3mm,
colback=blue!20!white,
colframe=blue!20!black,
sharp corners,rounded corners=southeast,arc is angular,arc=3mm,
underlay={%
\path[fill=tcbcolback!80!black] ([yshift=3mm]interior.south east)--++(-0.4,-0.1)--++(0.1,-0.2);
\path[draw=tcbcolframe,shorten <=-0.05mm,shorten >=-0.05mm] ([yshift=3mm]interior.south east)--++(-0.4,-0.1)--++(0.1,-0.2);
\path[fill=blue!50!black,draw=none] (interior.south west) rectangle node[white]{\Huge\bfseries ?} ([xshift=9mm]interior.north west);
},
drop fuzzy shadow,#1}

% Information box ======================================================
\newtcolorbox{info}[1][]{enhanced,
before skip=5mm,after skip=6mm, % [V3 ADDITION] was 2mm/3mm
boxrule=1pt,left=11mm,right=2mm,top=3mm,bottom=3mm,
colback=cyan!20!white,
colframe=cyan!20!black,
sharp corners,rounded corners=southeast,arc is angular,arc=3mm,
underlay={%
\path[fill=tcbcolback!80!black] ([yshift=3mm]interior.south east)--++(-0.4,-0.1)--++(0.1,-0.2);
\path[draw=tcbcolframe,shorten <=-0.05mm,shorten >=-0.05mm] ([yshift=3mm]interior.south east)--++(-0.4,-0.1)--++(0.1,-0.2);
\path[fill=cyan!50!black,draw=none] (interior.south west) rectangle node[white]{\Huge\bfseries i} ([xshift=9mm]interior.north west);
},
drop fuzzy shadow,#1}

% Mistake box ======================================================
\newtcolorbox{mistake}[1][]{enhanced,
before skip=5mm,after skip=6mm, % [V3 ADDITION] was 2mm/3mm
boxrule=1pt,left=11mm,right=2mm,top=3mm,bottom=3mm,
colback=red!20!white,
colframe=red!20!black,
sharp corners,rounded corners=southeast,arc is angular,arc=3mm,
underlay={%
\path[fill=tcbcolback!80!black] ([yshift=3mm]interior.south east)--++(-0.4,-0.1)--++(0.1,-0.2);
\path[draw=tcbcolframe,shorten <=-0.05mm,shorten >=-0.05mm] ([yshift=3mm]interior.south east)--++(-0.4,-0.1)--++(0.1,-0.2);
\path[fill=red!50!black,draw=none] (interior.south west) rectangle node[white]{\Huge\bfseries $\times$} ([xshift=9mm]interior.north west);
},
drop fuzzy shadow,#1}



% [V3 ADDITION] Worked example box ==================================
% Own counter, numbered within the section (Example 11.2.1, 11.2.2, ...).
% For per-subsection numbering (11.2.1.1) change `section' to `subsection'. Blue `question' boxes are kept for thought-provoking questions
% only; worked examples always use this box. Usage:
%   \begin{example} ...question text... \end{example}
%   \solution{...}
\newtcolorbox[auto counter,number within=section]{example}[1][]{%
enhanced, breakable=false, before skip=5mm, after skip=4mm,
colback=teal!5!white, colframe=teal!55!black, boxrule=0.8pt,
left=3mm, right=3mm, top=2mm, bottom=2mm,
sharp corners, rounded corners=northwest, arc=3mm,
fonttitle=\bfseries, coltitle=white,
title=Example~\thetcbcounter, attach boxed title to top left={yshift=-2mm, xshift=3mm},
boxed title style={colback=teal!55!black, sharp corners, rounded corners=northwest, arc=2mm},
drop fuzzy shadow, #1}

% [V3 ADDITION] Spacing between sections, subsections, parts, paragraphs =
% (fncychap's chapter heading is untouched; only \section and below.)
\makeatletter
\renewcommand\section{\@startsection{section}{1}{\z@}%
  {-2ex \@plus -.5ex \@minus -.2ex}{3.2ex \@plus .3ex}%
  {\normalfont\LARGE\bfseries\color{blue!50!black}}}
\renewcommand\subsection{\@startsection{subsection}{2}{\z@}%
  {-1.5ex \@plus -.5ex \@minus -.2ex}{2.6ex \@plus .3ex}%
  {\normalfont\Large\bfseries}}
\renewcommand\subsubsection{\@startsection{subsubsection}{3}{\z@}%
  {-4ex \@plus -1ex \@minus -.5ex}{1.6ex \@plus .2ex}%
  {\normalfont\large\bfseries\color{teal!50!black}}}
\makeatother
\setlength{\parindent}{0pt}
\setlength{\parskip}{0.9ex plus 0.3ex minus 0.2ex}

% Headers and Footers
\usepackage{fancyhdr,lastpage}
% [V3 R2] grey left-margin line on the out-of-syllabus pages (item 13)
\usepackage{eso-pic}
% [AoI item 63] thicker line (2.6pt, was 1.6pt) that starts just BELOW the header rule
% (top 14mm of the page left clear). Flag: copy this line into every chapter's preamble.
\newcommand{\oosmargin}{\AddToShipoutPictureBG{\put(12mm,0){\color{gray!75}\rule{2.6pt}{\dimexpr\paperheight-14mm\relax}}}}
\pagestyle{fancy}
\lhead{Maximus Pung Jun}
% [BUILD COPY CHANGE 2] Header bug fix: \chaptername is the word "Chapter",
% so the original header printed "Chapter 1: Chapter". \chaptermark now stores
% the bare chapter title in \leftmark, and the header uses \leftmark.
\renewcommand{\chaptermark}[1]{\markboth{#1}{}}
\rhead{Chapter \thechapter: \leftmark}
% \cfoot{Page \thepage\  of \pageref{LastPage}} % buggy, does not appear on chapter pages.

% [V3 R2] tick and cross markers for mistake boxes (G3)
\newcommand{\wrong}{\textcolor{red!70!black}{\ensuremath{\boldsymbol{\times}}}}
\newcommand{\cmark}{\textcolor{green!50!black}{\ensuremath{\checkmark}}}

\newcommand{\xmark}{\wrong}   % alias: red cross
% [V3 R2] out-of-syllabus helpers (DE notes item 13)
%   \oosmargin    : grey line in the left margin on every following page
%   \oosmarginoff : stop the margin line
%   \oosnote      : grey banner "Out of syllabus. For your interest only."
%   outofsyllabus : environment = new page + margin line + banner; ends with
%                   \clearpage and switches the margin line off again
\newcommand{\oosmarginoff}{\ClearShipoutPictureBG}
\newcommand{\oosnote}[1][Nothing in this section is examined in H2 Mathematics.]{\par\noindent\colorbox{gray!15}{\parbox{\dimexpr\textwidth-2\fboxsep}{\textbf{Out of syllabus. For your interest only.} #1}}\par}
\newenvironment{outofsyllabus}{\clearpage\oosmargin}{\clearpage\oosmarginoff}

% Notation as in Maximus's notes: italic d, display style everywhere (G1).
% Fractions inside exponents stay small, so never use these macros inside ^{}.
\newcommand{\dydx}{\dfrac{\mathrm{d}y}{\mathrm{d}x}}
\newcommand{\dvdx}{\dfrac{\mathrm{d}v}{\mathrm{d}x}}
\newcommand{\dudx}{\dfrac{\mathrm{d}u}{\mathrm{d}x}}
\newcommand{\ddydx}{\dfrac{\mathrm{d}^{2}y}{\mathrm{d}x^{2}}}
\newcommand{\ddt}[1]{\dfrac{\mathrm{d}#1}{\mathrm{d}t}}
% Small syllabus flag (no longer used in the DE notes, kept for other chapters)
\newcommand{\syllabusflag}[1]{\par\noindent{\footnotesize\color{red!70!black}\textbf{Syllabus note.} #1}\par}


% Teacher's copy toggle ================================================
% [BUILD COPY CHANGE 3] \ifx compares the whole word (the original
% \if\conform test was always false). DE_notes_student.tex predefines
% \conform as HideSolutions.
% A student copy defines \conform as HideSolutions BEFORE loading this preamble.
\providecommand\conform{ShowSolutions}
\newcommand\ShowSolutionsKey{ShowSolutions}
% [V3 ADDITION] In the student copy the hidden solution leaves stretchable
% blank space (\vfill), so the free space on each page is shared out as
% working space below the examples instead of collecting at the bottom.
% [AoI item 10] "Solution." on its own line, so that (a), (b), (c) line up below it.
\newcommand\solution[1]{\ifx\conform\ShowSolutionsKey
  \par\noindent\textbf{Solution.}\par\nobreak\vspace{0.5ex}\noindent #1\par\medskip\else\par\vspace{\workspace}\fi}
% [V3 R2] minimum working space left under each example in the student copy
\newlength\workspace\setlength\workspace{3cm plus 1fill}


\newcommand{\NN}{\mathbb{N}}
\newcommand{\ZZ}{\mathbb{Z}}
\newcommand{\QQ}{\mathbb{Q}}
\newcommand{\RR}{\mathbb{R}}
\newcommand{\CC}{\mathbb{C}}


% house_local.tex: HOUSE COPY of definite_integrals/common/ws_preamble_local.tex ([HOUSE] figstyles renamed)
% WORKSHEET COPY of ../notes/preamble_local.tex (version the worksheet was last built with,
% from ../_backup_pre_limitsum, 2026-10-09 17:31). Local so the worksheet no longer depends
% on ../notes; changes: \input{house_figstyles}; [2026-10-09 follow-up] \mk marks now bold, as in the notes.
% =====================================================================
% preamble_local.tex: extra packages, styles and macros for the
% Definite Integrals chapter ONLY. Loaded after preamble_build.tex.
% Plot styles and the mini 3D macros are in figstyles.tex (shared with the
% slides and worksheets).
% =====================================================================
\usepgfplotslibrary{fillbetween}
\usepackage{adjustbox} % mini 3D pictures (figstyles.tex)
\usetikzlibrary{arrows.meta, decorations.pathreplacing, calc}
\input{house_figstyles}

% Small fractions are allowed only in exponents and in the limits of an
% integral (which are superscripts/subscripts); everywhere else use \dfrac.
\newcommand{\dx}{\,\mathrm{d}x}
\newcommand{\dy}{\,\mathrm{d}y}
\newcommand{\evalat}[3]{\left[\,#1\,\right]_{#2}^{#3}}

% [item 78] right-aligned A-level style marks at the end of a question part
\newcommand{\mk}[1]{{\unskip\nobreak\hfil\penalty50\hskip1.5em\hbox{}\nobreak\hfill\textbf{[#1]}%
  \parfillskip=0pt\finalhyphendemerits=0\par}}
% [item 11] visually emphasised step labels in worked solutions
\newcommand{\step}[1]{\textcolor{teal!55!black}{\textbf{Step #1}}\quad}
% (a), (b), (c) part lists inside questions and solutions
\newenvironment{parts}{\begin{enumerate}[(a)]\setlength{\itemsep}{0.4ex}}{\end{enumerate}}
% SEAB-style solution layout (approved 10.5.8 template): working on the
% left, region sketch on the right (mini 3D pictures go under the sketch)
\newcommand{\solpair}[3][0.56]{\noindent\begin{minipage}[t]{#1\textwidth}\vspace{0pt}#2\end{minipage}\hfill
  \begin{minipage}[t]{\dimexpr0.98\textwidth-#1\textwidth\relax}\vspace{0pt}\centering #3\end{minipage}\par}
% minus-diagram style (drafts/minus_style_template.tex, items 29/31/62)
\pgfplotsset{minus/.style={axis lines=middle, width=5.0cm, height=4.2cm, ticks=none,
  xlabel={$x$}, ylabel={$y$}, every axis x label/.style={at={(ticklabel* cs:1)},anchor=west},
  every axis y label/.style={at={(ticklabel* cs:1)},anchor=south}, clip=false, axis on top,
  axis line style={gray!70}, every axis label/.append style={gray},
  every node/.append style={text=gray!80!black, font=\footnotesize}}}
\tikzset{hlreg/.style={fill=hl!35, draw=hl, thick}}

\ifx\housexm\housexnotes
  \ifdefined\housenotesLM\else% house_fonts.tex: HOUSE COPY of definite_integrals/common/prelim_fonts.tex (unchanged)
% =====================================================================
% prelim_fonts.tex: fonts of the real 9758 prelim papers (Word + MathType; pdffonts on
% ACJC 2026 P1/P2, ACJC 2025 P1, JPJC 2026, RVHS 2026, CJC/DHS/TJC/NYJC/HCI/EJC 2025):
%   text Times New Roman; cover/footer Arial; RI page footer Calibri (here Carlito);
%   maths = MathType: Times New Roman letters/digits (italic variables; upright digits,
%   function names, d, e, f), SymbolMT for the integral, sigma, pi and = + < >,
%   scripts 7 pt on 12 pt (about 58 %), tight operator spacing.
% Candidate C (teacher-approved 13 Oct 2026): the real Times New Roman and Arial
% (ttf-mscorefonts-installer) + URW Standard Symbols PS (fonts-urw-base35, a SymbolMT
% clone with the same Symbol encoding) for the integral, sigma, pi, = + < >;
% TeX Gyre Termes Math only for what MathType draws itself (radicals, stretchy brackets,
% fraction bars, arrows). House style: upright d (\mathrm{d}x), upright e and f.
% Loaded once by di_style / de_core (HW sheets) and prelim_style (Tests, solutions).
% =====================================================================
\ifdefined\prelimfontsloaded\endinput\fi
\def\prelimfontsloaded{1}
\edef\pfrestoreat{\catcode`\noexpand\@=\the\catcode`\@\relax}\makeatletter
% MathType script sizes: 7 pt superscripts on 12 pt text (58 %), 5 pt second level
\DeclareMathSizes{12}{12}{7}{5}
\DeclareMathSizes{11}{11}{6.5}{5}
\DeclareMathSizes{10.95}{10.95}{6.5}{5}
\DeclareMathSizes{10}{10}{6}{5}
\DeclareMathSizes{9}{9}{5.5}{5}
\DeclareMathSizes{14.4}{14.4}{8.5}{6}
\setmainfont{Times New Roman}
\setsansfont{Arial}
\newfontfamily\coverfont{Arial}
\newfontfamily\calibrifont{Carlito}
\usepackage{unicode-math}
\setmathfont{TeX Gyre Termes Math}
\setmathfont[range={up/{latin,Latin,num}}]{Times New Roman}
\setmathfont[range={it/{latin,Latin}}]{Times New Roman Italic}
\setmathfont[range={bfup/{latin,Latin,num}}]{Times New Roman Bold}
\setmathfont[range={"002B,"003D,"003C,"003E}]{Standard Symbols PS}
\newfontfamily\pf@sym{Standard Symbols PS}
\newfontfamily\pf@symi[FakeSlant=0.21]{Standard Symbols PS}
% MathType-like spacing (tighter than TeX around = + -)
\AtBeginDocument{\thickmuskip=2.5mu plus 1mu\medmuskip=1.5mu plus 1mu\thinmuskip=3mu
% Symbol integral and sigma (18 pt at 12 pt text, as MathType), slanted Symbol pi
 \renewcommand{\int}{\mathop{\vcenter{\hbox{\pf@sym\fontsize{\fpeval{1.5*\f@size}}{0pt}\selectfont\char"F2}}}\nolimits}%
 \renewcommand{\sum}{\mathop{\vcenter{\hbox{\pf@sym\fontsize{\fpeval{1.5*\f@size}}{0pt}\selectfont\char"E5}}}\limits}%
 \renewcommand{\pi}{\mathord{\text{\pf@symi\char"70}}}}
\pfrestoreat
\fi
  % =====================================================================
% house_tf.tex: True/False table (standing rules 2026-11-05/06/07):
%   * a T/F set is a NUMBERED QUESTION (no "Check Your Understanding" heading);
%   * table: row label (a),(b),... | Statement (left) | T/F column (right);
%   * row spacing slightly larger than default (extrarowheight 6pt, arraystretch 1.18,
%     as the Vectors Basics Worksheet after the 2026-11-07 change), font size unchanged;
%   * question doc: T/F column blank; solutions doc (\tfanswerstrue): bold T or F,
%     optional short reason under the statement in small grey.
% Usage:
%   \begin{tftable}  \tfrow{statement}{T}{optional reason}  ...  \end{tftable}
% Plain tabular: a question never splits across pages (house rule), so the table is kept whole.
% =====================================================================
\RequirePackage{array}
\makeatletter
\newif\iftfanswers \tfanswersfalse
\newcounter{tfrow}
\newlength{\tf@lab}\setlength{\tf@lab}{0.9cm}
\newlength{\tf@ans}\setlength{\tf@ans}{1.6cm}
\newenvironment{tftable}{\setcounter{tfrow}{0}\par\smallskip
  \setlength{\extrarowheight}{6pt}\renewcommand{\arraystretch}{1.18}%
  \noindent\begin{tabular}{|>{\centering\arraybackslash}p{\tf@lab}|p{\dimexpr\linewidth-\tf@lab-\tf@ans-6\tabcolsep-4\arrayrulewidth\relax}|>{\centering\arraybackslash}p{\tf@ans}|}
  \hline &\textbf{Statement}&\textbf{T/F}\\\hline}
 {\end{tabular}\par\smallskip}
\newcommand{\tfrow}[3]{\stepcounter{tfrow}\textbf{(\alph{tfrow})} & #1%
  \iftfanswers\if\relax\detokenize{#3}\relax\else\newline{\small\color{black!60}#3}\fi\fi
  & \iftfanswers\textbf{#2}\fi\\\hline}
% explicit row label (a set that keeps the Basics Worksheet letters, e.g. (e) omitted)
\newcommand{\tfrowl}[4]{\textbf{#1} & #2%
  \iftfanswers\if\relax\detokenize{#4}\relax\else\newline{\small\color{black!60}#4}\fi\fi
  & \iftfanswers\textbf{#3}\fi\\\hline}
% \tfcompact: allowed exception (Vectors I HW2/HW3) - compressed rows so the whole set
% stays on page 1 of the question sheet (extrarowheight 1.5pt, tabcolsep 4pt, 0.85/1.45 cm)
\newif\iftf@compact
\newcommand{\tfcompact}{\tf@compacttrue\setlength{\tf@lab}{0.85cm}\setlength{\tf@ans}{1.45cm}}
\renewenvironment{tftable}{\setcounter{tfrow}{0}\par\smallskip
  \iftf@compact\setlength{\extrarowheight}{1.5pt}\setlength{\tabcolsep}{4pt}\else\setlength{\extrarowheight}{6pt}\fi
  \renewcommand{\arraystretch}{1.18}%
  \noindent\begin{tabular}{|>{\centering\arraybackslash}p{\tf@lab}|p{\dimexpr\linewidth-\tf@lab-\tf@ans-6\tabcolsep-4\arrayrulewidth\relax}|>{\centering\arraybackslash}p{\tf@ans}|}
  \hline &\textbf{Statement}&\textbf{T/F}\\\hline}
 {\end{tabular}\par\smallskip}
% tflongtable: same table, may break across pages with the header repeated (solutions
% copies whose reasons make the set longer than a page). Needs longtable.
\RequirePackage{longtable}
\newenvironment{tflongtable}{\setcounter{tfrow}{0}\par\smallskip
  \setlength{\extrarowheight}{6pt}\renewcommand{\arraystretch}{1.18}%
  \setlength{\LTleft}{0pt}\setlength{\LTright}{0pt}%
  \begin{longtable}{|>{\centering\arraybackslash}p{\tf@lab}|p{\dimexpr\linewidth-\tf@lab-\tf@ans-6\tabcolsep-4\arrayrulewidth\relax}|>{\centering\arraybackslash}p{\tf@ans}|}
  \hline &\textbf{Statement}&\textbf{T/F}\\\hline\endhead}
 {\end{longtable}\par}
\makeatother

  \providecommand{\units}[1]{\ \text{units}^{#1}}
  \raggedbottom
\else
  % house_fonts.tex: HOUSE COPY of definite_integrals/common/prelim_fonts.tex (unchanged)
% =====================================================================
% prelim_fonts.tex: fonts of the real 9758 prelim papers (Word + MathType; pdffonts on
% ACJC 2026 P1/P2, ACJC 2025 P1, JPJC 2026, RVHS 2026, CJC/DHS/TJC/NYJC/HCI/EJC 2025):
%   text Times New Roman; cover/footer Arial; RI page footer Calibri (here Carlito);
%   maths = MathType: Times New Roman letters/digits (italic variables; upright digits,
%   function names, d, e, f), SymbolMT for the integral, sigma, pi and = + < >,
%   scripts 7 pt on 12 pt (about 58 %), tight operator spacing.
% Candidate C (teacher-approved 13 Oct 2026): the real Times New Roman and Arial
% (ttf-mscorefonts-installer) + URW Standard Symbols PS (fonts-urw-base35, a SymbolMT
% clone with the same Symbol encoding) for the integral, sigma, pi, = + < >;
% TeX Gyre Termes Math only for what MathType draws itself (radicals, stretchy brackets,
% fraction bars, arrows). House style: upright d (\mathrm{d}x), upright e and f.
% Loaded once by di_style / de_core (HW sheets) and prelim_style (Tests, solutions).
% =====================================================================
\ifdefined\prelimfontsloaded\endinput\fi
\def\prelimfontsloaded{1}
\edef\pfrestoreat{\catcode`\noexpand\@=\the\catcode`\@\relax}\makeatletter
% MathType script sizes: 7 pt superscripts on 12 pt text (58 %), 5 pt second level
\DeclareMathSizes{12}{12}{7}{5}
\DeclareMathSizes{11}{11}{6.5}{5}
\DeclareMathSizes{10.95}{10.95}{6.5}{5}
\DeclareMathSizes{10}{10}{6}{5}
\DeclareMathSizes{9}{9}{5.5}{5}
\DeclareMathSizes{14.4}{14.4}{8.5}{6}
\setmainfont{Times New Roman}
\setsansfont{Arial}
\newfontfamily\coverfont{Arial}
\newfontfamily\calibrifont{Carlito}
\usepackage{unicode-math}
\setmathfont{TeX Gyre Termes Math}
\setmathfont[range={up/{latin,Latin,num}}]{Times New Roman}
\setmathfont[range={it/{latin,Latin}}]{Times New Roman Italic}
\setmathfont[range={bfup/{latin,Latin,num}}]{Times New Roman Bold}
\setmathfont[range={"002B,"003D,"003C,"003E}]{Standard Symbols PS}
\newfontfamily\pf@sym{Standard Symbols PS}
\newfontfamily\pf@symi[FakeSlant=0.21]{Standard Symbols PS}
% MathType-like spacing (tighter than TeX around = + -)
\AtBeginDocument{\thickmuskip=2.5mu plus 1mu\medmuskip=1.5mu plus 1mu\thinmuskip=3mu
% Symbol integral and sigma (18 pt at 12 pt text, as MathType), slanted Symbol pi
 \renewcommand{\int}{\mathop{\vcenter{\hbox{\pf@sym\fontsize{\fpeval{1.5*\f@size}}{0pt}\selectfont\char"F2}}}\nolimits}%
 \renewcommand{\sum}{\mathop{\vcenter{\hbox{\pf@sym\fontsize{\fpeval{1.5*\f@size}}{0pt}\selectfont\char"E5}}}\limits}%
 \renewcommand{\pi}{\mathord{\text{\pf@symi\char"70}}}}
\pfrestoreat

  % house_hw.tex: HOUSE COPY of definite_integrals/common/di_style.tex; [HOUSE] chapter name from \housechapter{...}
% =====================================================================
% di_style.tex: Definite Integrals homework / test layout.
% Loads the house preamble (local copies of the AoI worksheet files:
% ws_preamble_build.tex, ws_preamble_local.tex -> ws_figstyles.tex),
% then adds the homework layout used by HW1-HW3 and the Test.
% Usage (before \begin{document}):  \ditle{Homework 1}
% =====================================================================
% [HOUSE] base, local and fonts are loaded by house.tex
\usepackage{needspace}
\usepackage{zref-savepos,zref-abspage}\makeatletter\zref@addprop{savepos}{abspage}\makeatother
\usepackage{array}

% ---- shortlist toggle ------------------------------------------------
% \ifshortlist true  -> only the questions in SHORTLIST.md (+ challenges)
% A wrapper file sets it with \def\setshortlist{1} before \input-ing the
% main file; the main file can also be edited to say \shortlisttrue.
\newif\ifshortlist
% Task 5: the FINAL homework (= former shortlist) is the default build. HWn.tex -> HWn.pdf.
% A wrapper that sets \def\setfull{1} builds the archive with all original questions.
\ifdefined\setfull\else\shortlisttrue\fi

\makeatletter
% ---- REVIEW build: highlight amended content ---------------------------
% A wrapper sets \def\setreview{1}. In review mode amended text is shaded
% yellow and carries a hover note (PDF tooltip / sticky note); otherwise the
% macros print the content plainly.
%   \chg{note}{inline content}        short inline change (no line breaks), text mode
%   \chgm{note}{math}                 the same inside math mode
%   \begin{chgblock}{note} ... \end{chgblock}   paragraph/display-level change
\newif\ifreview
\ifdefined\setreview \reviewtrue \fi
% hover notes are raw PDF annotations (xdvipdfmx \special; no extra packages).
% Keep notes plain ASCII; balanced parentheses are fine.
\newsavebox{\chg@box}
\newcommand{\chg@ann}[4]{% width, height, depth, note -> invisible markup annotation with a hover note
  \special{pdf:ann width #1 height #2 depth #3 << /Type /Annot /Subtype /Square /Border [0 0 0] /BS << /W 0 >> /C [1 0.85 0] /CA 0.01
    /T (Review) /Subj (Amendment) /Contents (#4) >>}}
\newcommand{\chg@sticky}[1]{\special{pdf:ann width 12pt height 12pt depth 0pt << /Type /Annot /Subtype /Text /Name /Comment
    /C [1 0.85 0] /T (Review) /Subj (Amendment) /Contents (#1) >>}}
\newcommand{\chg}[2]{\ifreview\sbox\chg@box{\colorbox{yellow!60}{#2}}%
  \leavevmode\rlap{\chg@ann{\the\wd\chg@box}{\the\ht\chg@box}{\the\dp\chg@box}{#1}}\usebox\chg@box\else#2\fi}
% \chgm{note}{math}: the same, for use INSIDE math mode (display style kept)
\newcommand{\chgm}[2]{\ifreview\mathchoice{\chg{#1}{$\displaystyle#2$}}{\chg{#1}{$\textstyle#2$}}{\chg{#1}{$\scriptstyle#2$}}{\chg{#1}{$\scriptscriptstyle#2$}}\else#2\fi}
\newtcolorbox{chgbox}[1]{enhanced, breakable, colback=yellow!35, colframe=yellow!35, boxrule=0pt, sharp corners,
  left=2pt, right=14pt, top=1pt, bottom=1pt, before skip=2pt, after skip=2pt,
  overlay first and middle={}, overlay last={}, overlay unbroken and first={\node[anchor=north east, inner sep=1pt] at (frame.north east) {\chg@sticky{#1}\hbox to 12pt{\vrule width 0pt height 12pt}};}}
% \chgnote{note}: just the sticky-note icon (e.g. inside a figure); nothing outside review mode
\newcommand{\chgnote}[1]{\ifreview\chg@sticky{#1}\hbox to 12pt{\vrule width 0pt height 12pt}\fi}
\newenvironment{chgblock}[1]{\ifreview\begin{chgbox}{#1}\fi}{\ifreview\end{chgbox}\fi}

% ---- header ---------------------------------------------------------
\newcommand\di@title{}
\makeatletter
\newcommand{\ditle}[1]{\gdef\di@title{#1}}
\fancypagestyle{di}{\fancyhf{}%
  \lhead{\small \housechap: \di@title}%
  \rhead{\small Compiled by Maximus Pung Jun}%
  \cfoot{\small\thepage}%
  \renewcommand{\headrulewidth}{0.4pt}}
\pagestyle{di}
% title block on page 1, as in his Typst sheets
\newcommand{\dititleblock}[1][]{\thispagestyle{di}%
  {\LARGE\scshape \housechap: \di@title}\par\smallskip
  {Compiled by Maximus Pung Jun}\par\ifx&#1&\else\smallskip{\small\textcolor{black!60}{#1}}\par\fi\medskip}
\makeatother
\newcommand{\disection}[1]{\par\bigskip\noindent{\Large\scshape #1}\par\medskip}

% ---- questions --------------------------------------------------------
% \Qhead{label}{source tag}: question heading; grey [tag] straight after the number, on the left (rule 3B, 9 Nov 2026)
\newcommand{\Qhead}[2]{\par\needspace{8\baselineskip}\bigskip\noindent
  {\large\bfseries #1}\setbox0=\hbox{#2}\ifdim\wd0=0pt\else\enspace{\small\textcolor{black!60}{[#2]}}\fi\par\smallskip}
% part lists with bold labels as he writes them: \pt{(a)}, \pt{(i)}, \pt{(b) (i)}
\newenvironment{qparts}[1][2.4em]{\begin{list}{}{\setlength\leftmargin{#1}%
  \setlength\labelwidth{\dimexpr#1-0.5em\relax}\setlength\labelsep{0.5em}%
  \setlength\itemsep{0.8ex}\setlength\topsep{0.6ex}\setlength\parsep{0.6ex}%
  \renewcommand\makelabel[1]{\textbf{##1}\hfil}}}{\end{list}}
\newcommand{\pt}[1]{\item[#1]}
% thin horizontal label for a challenge question
\newcommand{\challengebanner}{\par\bigskip\noindent\textcolor{black!50}{\rule[0.5ex]{\linewidth}{0.4pt}}\par
  \noindent{\large\scshape Challenge (optional)}\par}
% upright "units"
\newcommand{\units}[1]{\ \text{units}^{#1}}
% END line
\newcommand{\diend}[1][END OF WORKSHEET]{\par\vspace{2\baselineskip}\begin{center}\textbf{#1}\end{center}}

% ---- question registry (full / shortlist from one source) -------------
% \defQ{n}{tag}{body}   question n of the original sheet
% \defA{n}{answer body} its final answer, as printed in his answer section
% \putQ{n} / \putC{n}   print question n as the next (challenge) question
% \putA{n}              print its answer under the same new label
\newcounter{diq}
\makeatletter
\newcommand{\defQ}[3]{\expandafter\gdef\csname diQtag@#1\endcsname{#2}%
  \expandafter\long\expandafter\gdef\csname diQbody@#1\endcsname{#3}}
\newcommand{\defA}[2]{\expandafter\long\expandafter\gdef\csname diA@#1\endcsname{#2}}
\newcommand{\di@tag}[1]{\csname diQtag@#1\endcsname}% source tag only (no original Q number)
% A question never breaks across pages: it is measured and starts on a new page unless it
% fits whole on the rest of this one (only a question taller than a page may break, and then
% only between parts). Page positions of its start/end (zref labels q-s<n>/q-e<n> in the .aux)
% let the checker verify this.
\newsavebox{\di@qmeas}
\newcommand{\di@placeQ}[1]{\par\setbox\di@qmeas=\vbox{#1}%
  \ifdim\dimexpr\ht\di@qmeas+\dp\di@qmeas\relax<0.95\textheight
    \needspace{\dimexpr\ht\di@qmeas+\dp\di@qmeas+\bigskipamount\relax}\else\needspace{14\baselineskip}\fi
  \zsavepos{q-s\thediq}#1\par\zsavepos{q-e\thediq}}
\newcommand{\putQ}[1]{\stepcounter{diq}%
  \expandafter\xdef\csname diLab@#1\endcsname{\thediq}%
  \di@placeQ{\Qhead{\thediq}{\di@tag{#1}}\csname diQbody@#1\endcsname\par}}
\newcommand{\putC}[1]{\stepcounter{diq}%
  \expandafter\xdef\csname diLab@#1\endcsname{\thediq}%
  \expandafter\gdef\csname diChal@#1\endcsname{1}%
  \di@placeQ{\Qhead{\thediq\quad\normalsize\mdseries\scshape Challenge}{\di@tag{#1}}\csname diQbody@#1\endcsname\par}}
\newcommand{\putA}[1]{\par\needspace{4\baselineskip}%
  \begin{qparts}[2.4em]\item[{\csname diLab@#1\endcsname}]\csname diA@#1\endcsname\end{qparts}\medskip}
% ---- worked solutions: \defS{n}{body}, \putQS{n} prints question n + its solution ----
\newcommand{\defS}[2]{\expandafter\long\expandafter\gdef\csname diS@#1\endcsname{#2}}
\newcommand{\disolhead}{\par\medskip\noindent{\bfseries Solution}\par\nobreak\vspace{-0.6ex}\noindent\textcolor{black!35}{\rule{\linewidth}{0.4pt}}\par\nobreak}
\newcommand{\putQS}[1]{\putQ{#1}\disolhead\csname diS@#1\endcsname\par}
% source of a worked solution; visible only in REVIEW builds. #1: his / pub / pip
\newcommand{\solsrc}[2]{\ifreview\par\noindent\def\di@a{#1}\def\di@b{pip}%
  \ifx\di@a\di@b\colorbox{yellow!60}{\parbox{\dimexpr\linewidth-2\fboxsep-14pt}{\small\textbf{Pip-worked:} #2}}\chgnote{Pip-worked solution. No worked solution by the teacher and no published solution was found. Please check it.}%
  \else{\small\textcolor{black!55}{Source of this solution: #2}}\fi\par\smallskip\fi}
\makeatother
% compact sketch for worked solutions: \solfig[axis options]{plot commands}
\newcommand{\solfig}[2][]{\par\begin{center}\begin{tikzpicture}\begin{axis}[qfig, width=6.6cm, height=5cm, #1]#2\end{axis}\end{tikzpicture}\end{center}}
% a sketch added to an existing worked solution (teacher rule), shaded in REVIEW
\newenvironment{addedsketch}{\begin{chgblock}{Sketch added: teacher rule - every area or volume solution has a sketch.}}{\end{chgblock}}
% answer part list (compact)
\newenvironment{aparts}[1][2.4em]{\begin{qparts}[#1]\setlength\itemsep{0.3ex}}{\end{qparts}}

% ---- figure helpers ---------------------------------------------------------
\tikzset{rect/.style={fill=blue!10, draw=blue!55!black, thin}}
\pgfplotsset{qfig/.style={sketch, every node/.append style={font=\scriptsize}}}
% aligned working that starts on the same line as a part label: \pt{(iii)} \begin{topalign} ... \end{topalign}
\newenvironment{topalign}{\hfill$\displaystyle\begin{aligned}[t]}{\end{aligned}$\hfill\null\par\smallskip}
% one display that starts on the same line as a part label
\newcommand{\topdisplay}[1]{\hfill$\displaystyle #1$\hfill\null\par\smallskip}

  % house_solmacros.tex: HOUSE COPY of de-docs/common/de_style.tex minus \input{de_core} and the DE test titles.
% Adds the Ch11 solution helpers: soln/spart tables, sbs/sbsplain side-by-side, \note, \award, \tagm, \dstart.
% =====================================================================
% de_style.tex: Differential Equations (Ch11) homework / test layout.
% = Ch10 di_style (copied as de_core.tex; header text changed) + the DE
% notation and the framed solution boxes of the v1 DE documents
% (soln / spart, used by the Test worked solutions).
% =====================================================================
\providecommand{\e}{\mathrm{e}}
\providecommand{\cosec}{\operatorname{cosec}}
\providecommand{\dt}{\,\mathrm{d}t}
\providecommand{\abs}[1]{\left|#1\right|}
\providecommand{\dd}[2]{\dfrac{\mathrm{d}#1}{\mathrm{d}#2}}
\providecommand{\dvdx}{\dfrac{\mathrm{d}v}{\mathrm{d}x}}
\providecommand{\dudx}{\dfrac{\mathrm{d}u}{\mathrm{d}x}}
\colorlet{notegrey}{black!60}
\makeatletter
\protected\def\note#1{\ifmmode\text{\small\color{notegrey}\itshape #1}\else{\small\color{notegrey}\itshape #1}\fi}
\newcommand{\award}[1]{\unskip\nobreak\hfill\textbf{\small #1}}
\newcommand{\tagm}[1]{\tag*{\textbf{\small #1}}}
\newcommand{\mfAminus}{\[\begin{aligned}&\int\dfrac{1}{a^{2}-x^{2}}\dx\\&\quad=\dfrac{1}{2a}\ln\abs{\dfrac{a+x}{a-x}}+c\end{aligned}\]}
\newcommand{\mfXminus}{\[\begin{aligned}&\int\dfrac{1}{x^{2}-a^{2}}\dx\\&\quad=\dfrac{1}{2a}\ln\abs{\dfrac{x-a}{x+a}}+c\end{aligned}\]}
\newcommand{\sbs@start}{\if@inlabel\def\sbs@shift{\vspace{-\ht\strutbox}}\else
  \ifvmode\else\par\fi\addvspace{0.4ex}\def\sbs@shift{}\fi\noindent}
\newsavebox{\sbs@box}
\newenvironment{sbs}[1]{%
  \sbs@start\begin{minipage}[t]{#1\linewidth}\vspace{0pt}\sbs@shift\raggedright%
  \def\sbsnext{\end{minipage}\hfill\begin{minipage}[t]{\dimexpr0.97\linewidth-#1\linewidth\relax}\vspace{0pt}\sbs@shift\setlength{\fboxsep}{1.5mm}\begin{lrbox}{\sbs@box}\begin{minipage}[t]{\dimexpr\linewidth-2\fboxsep-1pt\relax}\small\raggedright}%
}{\end{minipage}\end{lrbox}\setlength{\fboxrule}{0.5pt}\fbox{\usebox{\sbs@box}}\end{minipage}\par}
\newenvironment{sbsplain}[1]{%
  \sbs@start\begin{minipage}[t]{#1\linewidth}\vspace{0pt}\sbs@shift\raggedright%
  \def\sbsnext{\end{minipage}\hfill\begin{minipage}[t]{\dimexpr0.97\linewidth-#1\linewidth\relax}\vspace{0pt}\sbs@shift\centering}%
}{\end{minipage}\par}
% ---- framed worked-solution box (v1 DE mark-scheme layout, kept for the Test) ----
\newlength{\solcol}\setlength{\solcol}{15mm}
\newcounter{spartno}
\newcommand{\solskips}{\abovedisplayskip=4pt plus 2pt\belowdisplayskip=4pt plus 2pt\abovedisplayshortskip=2pt\belowdisplayshortskip=3pt}
\newtcolorbox{soln}[1][]{enhanced, breakable, sharp corners,
  colback=white, colframe=black, boxrule=0.6pt,
  boxsep=0pt, left=0pt, right=3mm, top=1.5mm, bottom=2mm,
  before skip=1.2ex, after skip=1.5ex,
  before upper={\setcounter{spartno}{0}\parskip=0.5ex\solskips\@fleqntrue\@mathmargin=0pt\raggedright},
  overlay={\draw[line width=0.4pt] ([xshift=\solcol]frame.north west) -- ([xshift=\solcol]frame.south west);},
  #1}
\newenvironment{spart}[1]{%
  \stepcounter{spartno}%
  \ifnum\value{spartno}>1 \par\vspace{0.2ex}\noindent
    \makebox[\linewidth][l]{\rule{\dimexpr\linewidth+3mm\relax}{0.4pt}}\par\nobreak\vspace{0.3ex}\fi
  \begin{list}{}{\leftmargin=\dimexpr\solcol+2.5mm\relax
    \labelwidth=\solcol \labelsep=2.5mm \itemindent=0pt
    \topsep=0.3ex \parsep=0.5ex \rightmargin=0pt \@beginparpenalty=\@M}%
  \item[{\makebox[\solcol][l]{\strut\hspace{1.5mm}\textbf{#1}}}]%
  \ignorespaces}{\end{list}}
\makeatother
% ---- DE sketches: thin light axes (teacher rule), as the Ch10 'sketch' style ----
\pgfplotsset{dsketch/.style={sketch, axis lines=middle, axis line style={thin, black!45, -{Stealth[length=1.6mm]}},
    ticks=none, clip=false, width=6.5cm, height=4.6cm,
    every axis x label/.style={at={(current axis.right of origin)},anchor=west},
    every axis y label/.style={at={(current axis.above origin)},anchor=south},
    every node/.append style={font=\small}},
  dcurve/.style={black, semithick, smooth, samples=120}}
\tikzset{asym/.style={black, dashed, thin}}
\makeatletter
\newcommand{\dstart}{\leavevmode\vadjust{\nobreak\vskip\dimexpr-\baselineskip-\abovedisplayskip+1pt\relax}}
\newcommand{\sdstart}{\leavevmode\vadjust{\nobreak\vskip\dimexpr-\baselineskip-\abovedisplayskip+4pt\relax}}
\makeatother

  % =====================================================================
% house_tf.tex: True/False table (standing rules 2026-11-05/06/07):
%   * a T/F set is a NUMBERED QUESTION (no "Check Your Understanding" heading);
%   * table: row label (a),(b),... | Statement (left) | T/F column (right);
%   * row spacing slightly larger than default (extrarowheight 6pt, arraystretch 1.18,
%     as the Vectors Basics Worksheet after the 2026-11-07 change), font size unchanged;
%   * question doc: T/F column blank; solutions doc (\tfanswerstrue): bold T or F,
%     optional short reason under the statement in small grey.
% Usage:
%   \begin{tftable}  \tfrow{statement}{T}{optional reason}  ...  \end{tftable}
% Plain tabular: a question never splits across pages (house rule), so the table is kept whole.
% =====================================================================
\RequirePackage{array}
\makeatletter
\newif\iftfanswers \tfanswersfalse
\newcounter{tfrow}
\newlength{\tf@lab}\setlength{\tf@lab}{0.9cm}
\newlength{\tf@ans}\setlength{\tf@ans}{1.6cm}
\newenvironment{tftable}{\setcounter{tfrow}{0}\par\smallskip
  \setlength{\extrarowheight}{6pt}\renewcommand{\arraystretch}{1.18}%
  \noindent\begin{tabular}{|>{\centering\arraybackslash}p{\tf@lab}|p{\dimexpr\linewidth-\tf@lab-\tf@ans-6\tabcolsep-4\arrayrulewidth\relax}|>{\centering\arraybackslash}p{\tf@ans}|}
  \hline &\textbf{Statement}&\textbf{T/F}\\\hline}
 {\end{tabular}\par\smallskip}
\newcommand{\tfrow}[3]{\stepcounter{tfrow}\textbf{(\alph{tfrow})} & #1%
  \iftfanswers\if\relax\detokenize{#3}\relax\else\newline{\small\color{black!60}#3}\fi\fi
  & \iftfanswers\textbf{#2}\fi\\\hline}
% explicit row label (a set that keeps the Basics Worksheet letters, e.g. (e) omitted)
\newcommand{\tfrowl}[4]{\textbf{#1} & #2%
  \iftfanswers\if\relax\detokenize{#4}\relax\else\newline{\small\color{black!60}#4}\fi\fi
  & \iftfanswers\textbf{#3}\fi\\\hline}
% \tfcompact: allowed exception (Vectors I HW2/HW3) - compressed rows so the whole set
% stays on page 1 of the question sheet (extrarowheight 1.5pt, tabcolsep 4pt, 0.85/1.45 cm)
\newif\iftf@compact
\newcommand{\tfcompact}{\tf@compacttrue\setlength{\tf@lab}{0.85cm}\setlength{\tf@ans}{1.45cm}}
\renewenvironment{tftable}{\setcounter{tfrow}{0}\par\smallskip
  \iftf@compact\setlength{\extrarowheight}{1.5pt}\setlength{\tabcolsep}{4pt}\else\setlength{\extrarowheight}{6pt}\fi
  \renewcommand{\arraystretch}{1.18}%
  \noindent\begin{tabular}{|>{\centering\arraybackslash}p{\tf@lab}|p{\dimexpr\linewidth-\tf@lab-\tf@ans-6\tabcolsep-4\arrayrulewidth\relax}|>{\centering\arraybackslash}p{\tf@ans}|}
  \hline &\textbf{Statement}&\textbf{T/F}\\\hline}
 {\end{tabular}\par\smallskip}
% tflongtable: same table, may break across pages with the header repeated (solutions
% copies whose reasons make the set longer than a page). Needs longtable.
\RequirePackage{longtable}
\newenvironment{tflongtable}{\setcounter{tfrow}{0}\par\smallskip
  \setlength{\extrarowheight}{6pt}\renewcommand{\arraystretch}{1.18}%
  \setlength{\LTleft}{0pt}\setlength{\LTright}{0pt}%
  \begin{longtable}{|>{\centering\arraybackslash}p{\tf@lab}|p{\dimexpr\linewidth-\tf@lab-\tf@ans-6\tabcolsep-4\arrayrulewidth\relax}|>{\centering\arraybackslash}p{\tf@ans}|}
  \hline &\textbf{Statement}&\textbf{T/F}\\\hline\endhead}
 {\end{longtable}\par}
\makeatother

  \ifx\housexm\housexsol \def\prelimkeepcolour{1}% house_prelim.tex: HOUSE COPY of definite_integrals/common/prelim_style.tex; [HOUSE] titles from \housechapter
% =====================================================================
% prelim_style.tex: makes the Test look like a Singapore JC H2 Mathematics
% 9758 prelim paper (studied from ACJC/RI 2026 P1/P2 on the box:
% body Times New Roman + Times-style maths, cover in Arial).
% Fonts: prelim_fonts.tex (real Times New Roman/Arial + Times-style maths).
% Shared by Ch10 (di_style) and Ch11 (de_style; sets \prelimtitle/\prelimsubject). Loaded AFTER the chapter style (overrides its layout for the test only).
% =====================================================================
% [HOUSE] fonts loaded by house.tex
\providecommand{\prelimtitle}{\MakeUppercase{\housechap} TEST}
\providecommand{\prelimsubject}{\housechap}
\geometry{a4paper, top=22mm, bottom=24mm, left=25mm, right=25mm, headsep=6mm, footskip=12mm}
\setlength{\parindent}{0pt}
\setlength{\parskip}{0.9ex plus 0.3ex minus 0.2ex}
\linespread{1.3}
\setlength{\emergencystretch}{3em}
\clubpenalty=10000 \widowpenalty=10000 \interlinepenalty=10000
\AtBeginDocument{\fontsize{12}{14.4}\selectfont}
\renewcommand{\normalsize}{\fontsize{12}{14.4}\selectfont}

% ---- figures: plain black/grey like a printed paper ----------------------
\ifdefined\prelimkeepcolour\else % question paper: black/grey figures; solutions keep colour
\pgfplotsset{curve/.style={black, semithick, smooth, samples=150},
  line2/.style={black, semithick, samples=2},
  sketch/.append style={axis line style={thin, black}, every axis label/.append style={black}}}
\tikzset{region/.style={fill=black!15, draw=none}, regionB/.style={fill=black!30, draw=none},
  rect/.style={fill=black!8, draw=black, thin}, ax3/.style={-{Stealth[length=1.5mm]}, thin, black}}
\colorlet{blue}{black}\colorlet{red}{black}\colorlet{orange}{black}
\fi

% ---- header / footer: page number centred top (bold), [Turn over at bottom right
\makeatletter
\providecommand{\prelimfootlabel}{\prelimtitle\ 2026}
\newcommand{\pp@footer}{{\coverfont\scriptsize \prelimfootlabel\hspace{3em}H2 MATHEMATICS 9758}}
\fancypagestyle{prelim}{\fancyhf{}%
  \chead{\bfseries\thepage}%
  \lfoot{\pp@footer}%
  \rfoot{\ifnum\value{page}<\getpagerefnumber{LastPage}\relax\textbf{[Turn over}\fi}%
  \renewcommand{\headrulewidth}{0pt}\renewcommand{\footrulewidth}{0pt}}
\fancypagestyle{prelimcover}{\fancyhf{}%
  \rfoot{\textbf{[Turn over}}%
  \renewcommand{\headrulewidth}{0pt}\renewcommand{\footrulewidth}{0pt}}
\pagestyle{prelim}

% ---- marks: [n] regular weight, right-aligned at the end of the part ----------
\renewcommand{\mk}[1]{{\unskip\nobreak\hfil\penalty50\hskip1.5em\hbox{}\nobreak\hfill[#1]%
  \parfillskip=0pt\finalhyphendemerits=0\par}}

% ---- questions: bold number at the left margin, text indented; grey [tag] straight after the number (3B)
\newsavebox{\pp@qmeas}
\newcommand{\pp@qlist}[1]{%
  \begin{list}{}{\setlength\leftmargin{2.6em}\setlength\labelwidth{2.1em}\setlength\labelsep{0.5em}%
     \setlength\topsep{0pt}\setlength\itemsep{0pt}\setlength\parsep{0.9ex}%
     \renewcommand\makelabel[1]{\textbf{##1}\hfil}}%
  \item[\thediq]\@endparpenalty\@M\setbox0=\hbox{\di@tag{#1}}\ifdim\wd0=0pt\else{\small\textcolor{black!60}{[\di@tag{#1}]}}\par\nobreak\fi\csname diQbody@#1\endcsname\par\penalty\@M
  \end{list}}
% a question starts on a new page unless it fits whole on the current one
% (questions taller than a page may break, but only between parts)
\renewcommand{\putQ}[1]{\stepcounter{diq}%
  \expandafter\xdef\csname diLab@#1\endcsname{\thediq}%
  \par\setbox\pp@qmeas=\vbox{\pp@qlist{#1}}%
  \ifdim\dimexpr\ht\pp@qmeas+\dp\pp@qmeas\relax<0.95\textheight
    \needspace{\dimexpr\ht\pp@qmeas+\dp\pp@qmeas\relax}\else\needspace{14\baselineskip}\fi
  \zsavepos{q-s\thediq}\pp@qlist{#1}\par\zsavepos{q-e\thediq}\vspace{5ex plus 2ex minus 1ex}}
% part labels (a)/(i) bold, indented under the question text
\renewenvironment{qparts}[1][2.6em]{\begin{list}{}{\setlength\leftmargin{#1}%
  \setlength\labelwidth{\dimexpr#1-0.5em\relax}\setlength\labelsep{0.5em}%
  \setlength\itemsep{1.4ex}\setlength\topsep{1.2ex}\setlength\parsep{0.8ex}%
  \renewcommand\makelabel[1]{\textbf{##1}\hfil}}}{\end{list}}
\makeatother

% ---- cover page --------------------------------------------------------------
\newcommand{\prelimcover}[3]{% #1 total marks, #2 duration, #3 date line
\thispagestyle{prelimcover}%
{\coverfont\fontsize{10}{12}\selectfont\linespread{1}\selectfont
\noindent\begin{minipage}[t]{\linewidth}
  {\bfseries \prelimtitle}\\[0.2ex]
  {\bfseries Compiled by Maximus Pung Jun}\\[1.6ex]
  Higher 2
\end{minipage}\par\vspace{2ex}
\noindent\rule{\linewidth}{1.2pt}\par\vspace{0.4ex}
\noindent{\Large\bfseries MATHEMATICS}\hfill{\Large\bfseries 9758}\par\vspace{1.2ex}
\noindent \prelimsubject\hfill\textbf{#3}\\
QUESTION PAPER\hfill\textbf{#2}\par\vspace{2.5ex}
\noindent\begin{tabular}{@{}p{46mm}l@{}}Additional Materials: & Writing paper\\ & List of Formulae and Results (MF27)\end{tabular}\par\vspace{2ex}
\noindent\rule{\linewidth}{0.6pt}\par\vspace{0.6ex}
{\bfseries READ THESE INSTRUCTIONS FIRST}\par\vspace{2ex}
\setlength{\parskip}{0pt}
Write your name and class on all the work you hand in.\\
Write in dark blue or black pen.\\
You may use an HB pencil for any diagrams or graphs.\\
Do not use staples, paper clips, glue or correction fluid.\par\vspace{2ex}
Answer \textbf{all} the questions.\\
Give non-exact numerical answers correct to 3 significant figures, or 1 decimal place in the case of angles in degrees, unless a different level of accuracy is specified in the question.\\
You are expected to use an approved graphing calculator.\\
Unsupported answers from a graphing calculator are allowed unless a question specifically states otherwise.\\
Where unsupported answers from a graphing calculator are \textbf{not} allowed in a question, you are required to present the mathematical steps using mathematical notations and not calculator commands.\\
You are reminded of the need for clear presentation in your answers.\par\vspace{2ex}
The number of marks is given in brackets [\ ] at the end of each question or part question.\\
The total number of marks for this paper is #1.
\vfill
\noindent\rule{\linewidth}{0.6pt}\par\vspace{0.4ex}
\centerline{This document consists of \textbf{\pageref{LastPage}} printed pages.}\par\vspace{1ex}
}\clearpage}

% ---- solutions header block (page 1 of the worked solutions; same cover style) ----
\newcommand{\prelimsolhead}[1]{% #1 total marks
{\coverfont\fontsize{10}{12}\selectfont\linespread{1}\selectfont
\noindent{\bfseries \prelimtitle}\\[0.2ex]{\bfseries Compiled by Maximus Pung Jun}\\[1.6ex]Higher 2\par\vspace{2ex}
\noindent\rule{\linewidth}{1.2pt}\par\vspace{0.4ex}
\noindent{\Large\bfseries MATHEMATICS}\hfill{\Large\bfseries 9758}\par\vspace{1.2ex}
\noindent \prelimsubject\hfill\textbf{2026}\\
WORKED SOLUTIONS\hfill\textbf{Total marks: #1}\par\vspace{1ex}
\noindent\rule{\linewidth}{0.6pt}\par}\vspace{4ex}}
% house_sol.tex: HOUSE COPY of definite_integrals/common/ri_solstyle.tex (unchanged)
% =====================================================================
% ri_solstyle.tex: worked solutions in the layout of the RI 2026 Year 6 Prelim
% "Solutions with comments" (grail.moe 26612/26613), loaded AFTER prelim_style.
%   * page 1: bold school-style block + bold underlined title; later pages: running
%     title top-left with a rule under it; footer "Page N of M" bottom right (Calibri).
%   * each question: prelim-format question text (bold number, marks right, source tag
%     bottom right), then a bordered table:  Solution | Comments  (bold header row);
%     body rows  label + [marks] (bold, narrow left column) | working | comments,
%     a full-width rule between parts, working left-aligned (fleqn).
% Rows come from the existing sources: \pt in qparts (Ch10 HW), \slabel (Ch10 Test),
% spart (Ch11). Marks under each label are taken from the question's \mk.
% =====================================================================
\makeatletter
\usepackage{lastpage}
\usepackage{adjustbox}
\usepackage{zref-savepos,zref-abspage}
\zref@addprop{savepos}{abspage}
\geometry{a4paper, top=22mm, bottom=24mm, left=22mm, right=22mm, headsep=8mm, footskip=12mm, headheight=24pt}
\providecommand{\risoltitle}{\prelimtitle: Solutions with comments}
\providecommand{\risolschool}{Compiled by Maximus Pung Jun}
\fancypagestyle{risol}{\fancyhf{}%
  \lhead{\fontsize{11}{12}\selectfont 2026 \risoltitle\\[-1.2ex]\rule{\textwidth}{0.6pt}}%
  \rfoot{\calibrifont\fontsize{11}{13}\selectfont Page \textbf{\thepage} of \textbf{\pageref*{LastPage}}}%
  \renewcommand{\headrulewidth}{0pt}\renewcommand{\footrulewidth}{0pt}}
\fancypagestyle{risolfirst}{\fancyhf{}%
  \rfoot{\calibrifont\fontsize{11}{13}\selectfont Page \textbf{\thepage} of \textbf{\pageref*{LastPage}}}%
  \renewcommand{\headrulewidth}{0pt}\renewcommand{\footrulewidth}{0pt}}
\pagestyle{risol}
\newif\ifri@fresh
\newcommand{\risolfront}{\thispagestyle{risolfirst}\global\ri@freshtrue%
  {\noindent\bfseries\linespread{1}\selectfont \risolschool\\ H2 Mathematics (9758)\\ 2026\par}%
  \vspace{4ex}{\noindent\bfseries\underline{2026 \risoltitle}\par}\vspace{3ex}}
% ---- table geometry (RI: label 42pt | working 284pt | comments 142pt of 468pt) ----
\newlength{\ri@L}\setlength{\ri@L}{15mm}
\newlength{\ri@C}\setlength{\ri@C}{40mm}
\newlength{\ri@pad}\setlength{\ri@pad}{1.8mm}
\newif\ifri@sol \newif\ifri@first
\newcount\ri@depth
\def\ri@outer{}
\newtcolorbox{risol}{enhanced, breakable, sharp corners, frame hidden, colback=white,
  boxrule=0pt, boxsep=0pt, left=\dimexpr\ri@L+\ri@pad\relax, right=\dimexpr\ri@C+\ri@pad\relax,
  top=1.2mm, bottom=2.5mm, toptitle=0.6mm, bottomtitle=0.6mm, lefttitle=\ri@pad, righttitle=0pt,
  colbacktitle=white, coltitle=black, fonttitle=\bfseries,
  title={Solution\hfill\makebox[\ri@C][l]{\hspace{\ri@pad}Comments}},
  before skip=2ex, after skip=2ex,
  overlay unbroken and first={\draw[line width=0.5pt] (frame.south west) rectangle (frame.north east);
     \draw[line width=0.5pt] (title.south west) -- (title.south east);
     \draw[line width=0.5pt] ([xshift=\ri@L]interior.north west) -- ([xshift=\ri@L]frame.south west);
     \draw[line width=0.5pt] ([xshift=-\ri@C]frame.north east) -- ([xshift=-\ri@C]frame.south east);},
  overlay middle and last={\draw[line width=0.5pt] (frame.south west) rectangle (frame.north east);
     \draw[line width=0.5pt] ([xshift=\ri@L]frame.north west) -- ([xshift=\ri@L]frame.south west);
     \draw[line width=0.5pt] ([xshift=-\ri@C]frame.north east) -- ([xshift=-\ri@C]frame.south east);},
  after upper={\ri@endrow\protected@write\@auxout{}{\string\ri@settotal{\ri@qid}{\the\ri@sumH}}},
  before upper={\global\ri@sumH=14mm\ri@soltrue\ri@firsttrue\ri@depth=0 \parindent=0pt \parskip=0.6ex
     \predisplaypenalty=100 
     \linespread{1.12}\selectfont\raggedright\@fleqntrue\@mathmargin=0pt
     \abovedisplayskip=5pt plus 2pt\belowdisplayskip=5pt plus 2pt
     \abovedisplayshortskip=3pt plus 1pt\belowdisplayshortskip=4pt plus 2pt
     \ifcsname ri@hp@\ri@qid\endcsname\else\ri@row{}\fi}}
\def\ri@qid{}
% one row: full-width rule (glue leaders, so it vanishes at a page break where the frame
% edge takes its place; not before the first row), then bold label + bold [marks]
% in the left column, on the baseline of the first line of the working
\ExplSyntaxOn
\tl_new:N \g__ri_tmp_tl
\cs_new_protected:Npn \ri_setkey:Nn #1#2
  { \group_begin: \cs_set_eq:NN \chg \use_ii:nn \cs_set_eq:NN \chgm \use_ii:nn
    \cs_set_eq:NN \textbf \use:n \cs_set_eq:NN \mathrm \use:n
    \tl_gset:Nx \g__ri_tmp_tl {#2} \group_end:
    \tl_set_eq:NN #1 \g__ri_tmp_tl \tl_remove_all:Nn #1 {~} }
\cs_set_eq:NN \ri@setkey \ri_setkey:Nn
\ExplSyntaxOff
% ---- page-break rules (positions from the previous run, zref labels in the .aux) ----
% * a row (one part) never breaks across pages: if it does not fit on the rest of the
%   page it starts the next page (forced break between rows); only a row taller than a
%   whole page may break inside;
% * a whole solution that fits on one page is kept together (needspace before the table);
%   a longer one starts where at least its first row fits;
% * a row is at least as tall as its label + [marks].
\newcount\ri@n \newbox\ri@labbox \def\ri@open{0}
\newdimen\ri@ybot \newdimen\ri@ytop \newdimen\ri@h \newdimen\ri@sumH
\AtBeginDocument{\ri@ybot=\dimexpr\paperheight-1in-\voffset-\topmargin-\headheight-\headsep-\textheight\relax
  \ri@ytop=\dimexpr\ri@ybot+\textheight\relax}
\def\ri@pg#1{\zref@extractdefault{#1}{abspage}{-1}}
% \ri@rowh{n}: height of row n in the previous run -> \ri@h (0pt if unknown)
\newcommand{\ri@rowh}[1]{\ri@h=0pt
  \ifnum\ri@pg{ri-s#1}>0 \ifnum\ri@pg{ri-e#1}>0
    \ifnum\ri@pg{ri-s#1}=\ri@pg{ri-e#1}\relax \ri@h=\dimexpr\zposy{ri-s#1}sp-\zposy{ri-e#1}sp\relax
    \else \ri@h=\dimexpr\zposy{ri-s#1}sp-\ri@ybot+\ri@ytop-\zposy{ri-e#1}sp-6mm\relax\fi\fi\fi}
\newcommand{\ri@setfirst}[2]{\expandafter\gdef\csname ri@first@#1\endcsname{#2}}
\newcommand{\ri@setkept}[1]{\expandafter\gdef\csname ri@kept@#1\endcsname{}}
\newcommand{\ri@settotal}[2]{\expandafter\gdef\csname ri@total@#1\endcsname{#2}}
\newcommand{\ri@endrow}{\ifnum\ri@open=1 \par\zsavepos{ri-e\the\ri@n}%
  \ifnum\ri@pg{ri-s\the\ri@n}=\ri@pg{ri-e\the\ri@n}\relax
    \dimen@=\dimexpr\ri@curH sp+1.6mm-\zposy{ri-s\the\ri@n}sp+\zposy{ri-e\the\ri@n}sp\relax
    \ifdim\dimen@>0pt \nobreak\kern\dimen@\penalty\z@\fi\fi
  \ri@rowh{\the\ri@n}\global\advance\ri@sumH\dimexpr\ri@h+3.5mm\relax
  \gdef\ri@open{0}\fi}
\newcommand{\ri@row}[1]{\par\ri@endrow\global\advance\ri@n\@ne\gdef\ri@open{1}\@beginparpenalty\@M
  \ifri@first\global\ri@firstfalse\ri@firstfalse
    \protected@write\@auxout{}{\string\ri@setfirst{\ri@qid}{\the\ri@n}}%
  \else
    \zsavepos{ri-p\the\ri@n}%
    % push the row to the next page if it would not fit on the rest of this one
    \ri@rowh{\the\ri@n}%
    \ifnum\ri@pg{ri-p\the\ri@n}>0 \ifdim\ri@h>0pt
      \ifdim\dimexpr\ri@h+4mm\relax<\dimexpr\textheight-10mm\relax
        \ifdim\dimexpr\ri@h+7mm\relax>\dimexpr\zposy{ri-p\the\ri@n}sp-\ri@ybot\relax
          \penalty-\@M\fi\fi\fi\fi
    \vskip0.6ex\nointerlineskip
    \leaders\hbox{\hskip-\dimexpr\ri@L+\ri@pad\relax\vrule width\dimexpr\linewidth+\ri@L+\ri@C+2\ri@pad\relax height0.5pt depth0pt}\vskip0.5pt
    \nobreak\vskip1.4mm\fi
  \ri@setkey\ri@key{\ri@qid#1}\def\ri@lab{#1}%
  \setbox\ri@labbox=\vtop{\hsize=\dimexpr\ri@L-1.6mm\relax\bfseries\raggedright\parskip=0pt
      \strut\ri@lab\ifcsname ri@mk@\ri@key\endcsname\ifx\ri@lab\@empty\else\\\fi
      [\csname ri@mk@\ri@key\endcsname]\fi\strut\par}%
  \xdef\ri@curH{\number\dimexpr\ht\ri@labbox+\dp\ri@labbox\relax}%
  \nointerlineskip\zsavepos{ri-s\the\ri@n}\nointerlineskip\nobreak
  \vbox to 0pt{\hbox{\hskip-\dimexpr\ri@L+\ri@pad-1.2mm\relax\box\ri@labbox}\vss}\nobreak\prevdepth=\dp\strutbox
  \vskip-\parskip\ignorespaces}
% ---- qparts / \pt: question mode records labels (for the marks); solution mode makes rows
\NewCommandCopy\ri@qparts@begin\qparts \let\ri@qparts@end\endqparts
\renewenvironment{qparts}[1][2.6em]{\advance\ri@depth\@ne
  \ifri@sol\else\ri@qparts@begin[#1]\fi}{\ifri@sol\par\else\ri@qparts@end\fi}
\def\ri@qpartsname{qparts}
\renewcommand{\pt}[1]{\ifnum\ri@depth<2 \gdef\ri@outer{#1}\gdef\ri@cur{#1}\else\expandafter\gdef\expandafter\ri@cur\expandafter{\ri@outer#1}\fi
  \ifri@sol\expandafter\ri@solpt\else\expandafter\ri@qpt\fi{#1}}
\def\ri@qpt#1{\expandafter\gdef\csname ri@hp@\ri@qid\endcsname{}\item[#1]}
\def\ri@solpt#1{\@ifnextchar\begin\ri@ptpeek{\ri@row{\ri@cur}}}
\def\ri@ptpeek\begin#1{\def\ri@tmp{#1}\ifx\ri@tmp\ri@qpartsname
   \expandafter\@firstoftwo\else\expandafter\@secondoftwo\fi
  {\begin{qparts}}{\ri@row{\ri@cur}\begin{#1}}}
% marks: record [n] against the current question/part
\let\ri@mk\mk
\renewcommand{\mk}[1]{\ri@setkey\ri@k{\ri@qid\ri@cur}\ri@setkey\ri@mv{#1}\expandafter\xdef\csname ri@mk@\ri@k\endcsname{\ri@mv}\ri@mk{#1}}
% \slabel (Ch10 Test answers): letter labels set the outer part, roman ones are appended;
% consecutive \slabel calls merge into one row label
\def\ri@romans{(i)(ii)(iii)(iv)(v)(vi)(vii)}
\def\ri@firstpart(#1)#2\@nil{(#1)}
\def\slabel#1{\ifri@sol\expandafter\ri@solslabel\else\expandafter\ri@qslabel\fi{#1}}
\def\ri@qslabel#1{\par\smallskip\noindent\textbf{#1}\par\nobreak}
\def\ri@solslabel#1{\@expandtwoargs\in@{#1}{\ri@romans}%
  \ifin@\expandafter\gdef\expandafter\ri@cur\expandafter{\ri@outer#1}\else\gdef\ri@cur{#1}\xdef\ri@outer{\ri@firstpart#1\@nil}\fi
  \@ifnextchar\slabel{}{\ri@row{\ri@cur}}}
% Ch11 spart{label}: a row; soln: the table
\def\spart#1{\gdef\ri@cur{#1}\ri@row{#1}}\def\endspart{\par}
\def\soln{\@ifnextchar[{\ri@solnopt}{\risol}}\def\ri@solnopt[#1]{\risol}\def\endsoln{\endrisol}
% ---- working that started on the label line is left-aligned in the working column
\AtBeginDocument{\@ifundefined{topalign}{}{%
  \renewenvironment{topalign}{\par\noindent$\displaystyle\begin{aligned}[t]}{\end{aligned}$\par\smallskip}%
  \renewcommand{\topdisplay}[1]{\par\noindent$\displaystyle #1$\par\smallskip}}}
% ---- Ch11 side-by-side (text | sketch): the sketch shrinks to fit its half of the working column
\AtBeginDocument{\@ifundefined{sbsplain}{}{%
  \renewenvironment{sbsplain}[1]{%
    \sbs@start\begin{minipage}[t]{#1\linewidth}\vspace{0pt}\sbs@shift\raggedright%
    \def\sbsnext{\end{minipage}\hfill\begin{minipage}[t]{\dimexpr0.97\linewidth-#1\linewidth\relax}\vspace{0pt}\sbs@shift\centering
      \begin{adjustbox}{max width=\linewidth}}%
  }{\end{adjustbox}\end{minipage}\par}}}
% ---- figures in the working column shrink to fit it (only if wider), colours kept
\newcommand{\rifitfig}[1]{\expandafter\let\csname ri@orig@\string#1\endcsname#1%
  \def#1{\@ifnextchar\bgroup{\ri@fitarg#1}{\ri@fitnoarg#1}}}
\def\ri@fitarg#1#2{\ifri@sol\adjustbox{max width=\linewidth}{\csname ri@orig@\string#1\endcsname{#2}}\else\csname ri@orig@\string#1\endcsname{#2}\fi}
\def\ri@fitnoarg#1{\ifri@sol\adjustbox{max width=\linewidth}{\csname ri@orig@\string#1\endcsname}\else\csname ri@orig@\string#1\endcsname\fi}
\@ifundefined{solfig}{}{\renewcommand{\solfig}[2][]{\par\begin{center}\adjustbox{max width=\linewidth}{\begin{tikzpicture}\begin{axis}[qfig, width=6.6cm, height=5cm, #1]#2\end{axis}\end{tikzpicture}}\end{center}}}
% \ricomment{text}: text in the Comments column of the current row (call right after \pt,
% \slabel or \spart). The row is made at least as tall as the comment.
\newbox\ri@cmtbox
\newcommand{\ricomment}[1]{\par\setbox\ri@cmtbox=\vtop{\hsize=\dimexpr\ri@C-2\ri@pad\relax
    \parskip=0pt\raggedright\small{\color{black!60}\strut#1\strut\par}}%
  \ifdim\dimexpr\ht\ri@cmtbox+\dp\ri@cmtbox\relax>\ri@curH sp\xdef\ri@curH{\number\dimexpr\ht\ri@cmtbox+\dp\ri@cmtbox\relax}\fi
  \nointerlineskip\vbox to 0pt{\hbox to\linewidth{\hskip\dimexpr\linewidth+2\ri@pad\relax\box\ri@cmtbox\hss}\vss}\nobreak\prevdepth=\dp\strutbox
  \vskip-\parskip\ignorespaces}
% \putTFlong{n}: the True/False question n with its answer table allowed to break across
% pages (the body of \defQ{n} must be: intro text, then \begin{tftable}...\end{tftable});
% the solutions copy prints it with tflongtable.
\renewcommand{\putTFlong}[1]{\ri@newq\begingroup\let\tftable\tflongtable\let\endtftable\endtflongtable
  \stepcounter{diq}\expandafter\xdef\csname diLab@#1\endcsname{\thediq}%
  \par\noindent\hbox to 2.6em{\textbf{\thediq}\hfil}\setbox0=\hbox{\di@tag{#1}}\ifdim\wd0=0pt\else{\small\textcolor{black!60}{[\di@tag{#1}]}}\par\noindent\hskip2.6em\fi\ignorespaces\csname diQbody@#1\endcsname\par
  \endgroup\vspace{3ex}}
% ---- question + solution; each question starts on a new page (as RI)
\def\ri@qid{}
\let\ri@putQ\putQ
\renewcommand{\putQ}[1]{\gdef\ri@qid{#1}\global\ri@depth=0 \gdef\ri@cur{}\ri@putQ{#1}}
% \risolQ{n}{body macro name}: question n + its solution body in the table;
% \risolQraw: body already contains its own soln environment (Ch11 Test)
\newcommand{\ri@newq}{\ifri@fresh\global\ri@freshfalse\else\clearpage\fi}
% keep the whole table together if it fits on a page; otherwise make sure at least its
% header and first row fit before starting it
\newcommand{\ri@keep}[1]{\par
  \ifcsname ri@total@#1\endcsname
    % hysteresis: a table kept whole in the previous run stays whole up to textheight-4.5mm,
    % so a total near the limit cannot flip between runs
    \ifdim\csname ri@total@#1\endcsname<\dimexpr\textheight-\ifcsname ri@kept@#1\endcsname4.5\else8\fi mm\relax
      \protected@write\@auxout{}{\string\ri@setkept{#1}}%
      \needspace{\dimexpr\csname ri@total@#1\endcsname+4mm\relax}%
    \else\ifcsname ri@first@#1\endcsname
      \ri@rowh{\csname ri@first@#1\endcsname}%
      \ifdim\ri@h<\dimexpr\textheight-20mm\relax\needspace{\dimexpr\ri@h+18mm\relax}\fi\fi\fi\fi}
\newcommand{\risolQ}[2]{\ri@newq\putQ{#1}\vspace{-3ex}\ri@keep{#1}\begin{risol}\csname #2@#1\endcsname\par\end{risol}}
\newcommand{\risolQraw}[2]{\ri@newq\putQ{#1}\vspace{-3ex}\ri@keep{#1}\csname #2@#1\endcsname\par}
\renewcommand{\putQS}[1]{\risolQ{#1}{diS}}
% end marker only if it fits on the last page (no page holding just the marker)
\renewcommand{\diend}[1][END OF SOLUTIONS]{\par\ifdim\dimexpr\pagegoal-\pagetotal\relax<4\baselineskip\else\vspace{1.5ex}\begin{center}\textbf{#1}\end{center}\fi}
\makeatother
\tfanswerstrue\fi
  \ifx\housexm\housextest % house_prelim.tex: HOUSE COPY of definite_integrals/common/prelim_style.tex; [HOUSE] titles from \housechapter
% =====================================================================
% prelim_style.tex: makes the Test look like a Singapore JC H2 Mathematics
% 9758 prelim paper (studied from ACJC/RI 2026 P1/P2 on the box:
% body Times New Roman + Times-style maths, cover in Arial).
% Fonts: prelim_fonts.tex (real Times New Roman/Arial + Times-style maths).
% Shared by Ch10 (di_style) and Ch11 (de_style; sets \prelimtitle/\prelimsubject). Loaded AFTER the chapter style (overrides its layout for the test only).
% =====================================================================
% [HOUSE] fonts loaded by house.tex
\providecommand{\prelimtitle}{\MakeUppercase{\housechap} TEST}
\providecommand{\prelimsubject}{\housechap}
\geometry{a4paper, top=22mm, bottom=24mm, left=25mm, right=25mm, headsep=6mm, footskip=12mm}
\setlength{\parindent}{0pt}
\setlength{\parskip}{0.9ex plus 0.3ex minus 0.2ex}
\linespread{1.3}
\setlength{\emergencystretch}{3em}
\clubpenalty=10000 \widowpenalty=10000 \interlinepenalty=10000
\AtBeginDocument{\fontsize{12}{14.4}\selectfont}
\renewcommand{\normalsize}{\fontsize{12}{14.4}\selectfont}

% ---- figures: plain black/grey like a printed paper ----------------------
\ifdefined\prelimkeepcolour\else % question paper: black/grey figures; solutions keep colour
\pgfplotsset{curve/.style={black, semithick, smooth, samples=150},
  line2/.style={black, semithick, samples=2},
  sketch/.append style={axis line style={thin, black}, every axis label/.append style={black}}}
\tikzset{region/.style={fill=black!15, draw=none}, regionB/.style={fill=black!30, draw=none},
  rect/.style={fill=black!8, draw=black, thin}, ax3/.style={-{Stealth[length=1.5mm]}, thin, black}}
\colorlet{blue}{black}\colorlet{red}{black}\colorlet{orange}{black}
\fi

% ---- header / footer: page number centred top (bold), [Turn over at bottom right
\makeatletter
\providecommand{\prelimfootlabel}{\prelimtitle\ 2026}
\newcommand{\pp@footer}{{\coverfont\scriptsize \prelimfootlabel\hspace{3em}H2 MATHEMATICS 9758}}
\fancypagestyle{prelim}{\fancyhf{}%
  \chead{\bfseries\thepage}%
  \lfoot{\pp@footer}%
  \rfoot{\ifnum\value{page}<\getpagerefnumber{LastPage}\relax\textbf{[Turn over}\fi}%
  \renewcommand{\headrulewidth}{0pt}\renewcommand{\footrulewidth}{0pt}}
\fancypagestyle{prelimcover}{\fancyhf{}%
  \rfoot{\textbf{[Turn over}}%
  \renewcommand{\headrulewidth}{0pt}\renewcommand{\footrulewidth}{0pt}}
\pagestyle{prelim}

% ---- marks: [n] regular weight, right-aligned at the end of the part ----------
\renewcommand{\mk}[1]{{\unskip\nobreak\hfil\penalty50\hskip1.5em\hbox{}\nobreak\hfill[#1]%
  \parfillskip=0pt\finalhyphendemerits=0\par}}

% ---- questions: bold number at the left margin, text indented; grey [tag] straight after the number (3B)
\newsavebox{\pp@qmeas}
\newcommand{\pp@qlist}[1]{%
  \begin{list}{}{\setlength\leftmargin{2.6em}\setlength\labelwidth{2.1em}\setlength\labelsep{0.5em}%
     \setlength\topsep{0pt}\setlength\itemsep{0pt}\setlength\parsep{0.9ex}%
     \renewcommand\makelabel[1]{\textbf{##1}\hfil}}%
  \item[\thediq]\@endparpenalty\@M\setbox0=\hbox{\di@tag{#1}}\ifdim\wd0=0pt\else{\small\textcolor{black!60}{[\di@tag{#1}]}}\par\nobreak\fi\csname diQbody@#1\endcsname\par\penalty\@M
  \end{list}}
% a question starts on a new page unless it fits whole on the current one
% (questions taller than a page may break, but only between parts)
\renewcommand{\putQ}[1]{\stepcounter{diq}%
  \expandafter\xdef\csname diLab@#1\endcsname{\thediq}%
  \par\setbox\pp@qmeas=\vbox{\pp@qlist{#1}}%
  \ifdim\dimexpr\ht\pp@qmeas+\dp\pp@qmeas\relax<0.95\textheight
    \needspace{\dimexpr\ht\pp@qmeas+\dp\pp@qmeas\relax}\else\needspace{14\baselineskip}\fi
  \zsavepos{q-s\thediq}\pp@qlist{#1}\par\zsavepos{q-e\thediq}\vspace{5ex plus 2ex minus 1ex}}
% part labels (a)/(i) bold, indented under the question text
\renewenvironment{qparts}[1][2.6em]{\begin{list}{}{\setlength\leftmargin{#1}%
  \setlength\labelwidth{\dimexpr#1-0.5em\relax}\setlength\labelsep{0.5em}%
  \setlength\itemsep{1.4ex}\setlength\topsep{1.2ex}\setlength\parsep{0.8ex}%
  \renewcommand\makelabel[1]{\textbf{##1}\hfil}}}{\end{list}}
\makeatother

% ---- cover page --------------------------------------------------------------
\newcommand{\prelimcover}[3]{% #1 total marks, #2 duration, #3 date line
\thispagestyle{prelimcover}%
{\coverfont\fontsize{10}{12}\selectfont\linespread{1}\selectfont
\noindent\begin{minipage}[t]{\linewidth}
  {\bfseries \prelimtitle}\\[0.2ex]
  {\bfseries Compiled by Maximus Pung Jun}\\[1.6ex]
  Higher 2
\end{minipage}\par\vspace{2ex}
\noindent\rule{\linewidth}{1.2pt}\par\vspace{0.4ex}
\noindent{\Large\bfseries MATHEMATICS}\hfill{\Large\bfseries 9758}\par\vspace{1.2ex}
\noindent \prelimsubject\hfill\textbf{#3}\\
QUESTION PAPER\hfill\textbf{#2}\par\vspace{2.5ex}
\noindent\begin{tabular}{@{}p{46mm}l@{}}Additional Materials: & Writing paper\\ & List of Formulae and Results (MF27)\end{tabular}\par\vspace{2ex}
\noindent\rule{\linewidth}{0.6pt}\par\vspace{0.6ex}
{\bfseries READ THESE INSTRUCTIONS FIRST}\par\vspace{2ex}
\setlength{\parskip}{0pt}
Write your name and class on all the work you hand in.\\
Write in dark blue or black pen.\\
You may use an HB pencil for any diagrams or graphs.\\
Do not use staples, paper clips, glue or correction fluid.\par\vspace{2ex}
Answer \textbf{all} the questions.\\
Give non-exact numerical answers correct to 3 significant figures, or 1 decimal place in the case of angles in degrees, unless a different level of accuracy is specified in the question.\\
You are expected to use an approved graphing calculator.\\
Unsupported answers from a graphing calculator are allowed unless a question specifically states otherwise.\\
Where unsupported answers from a graphing calculator are \textbf{not} allowed in a question, you are required to present the mathematical steps using mathematical notations and not calculator commands.\\
You are reminded of the need for clear presentation in your answers.\par\vspace{2ex}
The number of marks is given in brackets [\ ] at the end of each question or part question.\\
The total number of marks for this paper is #1.
\vfill
\noindent\rule{\linewidth}{0.6pt}\par\vspace{0.4ex}
\centerline{This document consists of \textbf{\pageref{LastPage}} printed pages.}\par\vspace{1ex}
}\clearpage}

% ---- solutions header block (page 1 of the worked solutions; same cover style) ----
\newcommand{\prelimsolhead}[1]{% #1 total marks
{\coverfont\fontsize{10}{12}\selectfont\linespread{1}\selectfont
\noindent{\bfseries \prelimtitle}\\[0.2ex]{\bfseries Compiled by Maximus Pung Jun}\\[1.6ex]Higher 2\par\vspace{2ex}
\noindent\rule{\linewidth}{1.2pt}\par\vspace{0.4ex}
\noindent{\Large\bfseries MATHEMATICS}\hfill{\Large\bfseries 9758}\par\vspace{1.2ex}
\noindent \prelimsubject\hfill\textbf{2026}\\
WORKED SOLUTIONS\hfill\textbf{Total marks: #1}\par\vspace{1ex}
\noindent\rule{\linewidth}{0.6pt}\par}\vspace{4ex}}
\fi
  \ifx\housexm\housextestsol \def\prelimkeepcolour{1}% house_prelim.tex: HOUSE COPY of definite_integrals/common/prelim_style.tex; [HOUSE] titles from \housechapter
% =====================================================================
% prelim_style.tex: makes the Test look like a Singapore JC H2 Mathematics
% 9758 prelim paper (studied from ACJC/RI 2026 P1/P2 on the box:
% body Times New Roman + Times-style maths, cover in Arial).
% Fonts: prelim_fonts.tex (real Times New Roman/Arial + Times-style maths).
% Shared by Ch10 (di_style) and Ch11 (de_style; sets \prelimtitle/\prelimsubject). Loaded AFTER the chapter style (overrides its layout for the test only).
% =====================================================================
% [HOUSE] fonts loaded by house.tex
\providecommand{\prelimtitle}{\MakeUppercase{\housechap} TEST}
\providecommand{\prelimsubject}{\housechap}
\geometry{a4paper, top=22mm, bottom=24mm, left=25mm, right=25mm, headsep=6mm, footskip=12mm}
\setlength{\parindent}{0pt}
\setlength{\parskip}{0.9ex plus 0.3ex minus 0.2ex}
\linespread{1.3}
\setlength{\emergencystretch}{3em}
\clubpenalty=10000 \widowpenalty=10000 \interlinepenalty=10000
\AtBeginDocument{\fontsize{12}{14.4}\selectfont}
\renewcommand{\normalsize}{\fontsize{12}{14.4}\selectfont}

% ---- figures: plain black/grey like a printed paper ----------------------
\ifdefined\prelimkeepcolour\else % question paper: black/grey figures; solutions keep colour
\pgfplotsset{curve/.style={black, semithick, smooth, samples=150},
  line2/.style={black, semithick, samples=2},
  sketch/.append style={axis line style={thin, black}, every axis label/.append style={black}}}
\tikzset{region/.style={fill=black!15, draw=none}, regionB/.style={fill=black!30, draw=none},
  rect/.style={fill=black!8, draw=black, thin}, ax3/.style={-{Stealth[length=1.5mm]}, thin, black}}
\colorlet{blue}{black}\colorlet{red}{black}\colorlet{orange}{black}
\fi

% ---- header / footer: page number centred top (bold), [Turn over at bottom right
\makeatletter
\providecommand{\prelimfootlabel}{\prelimtitle\ 2026}
\newcommand{\pp@footer}{{\coverfont\scriptsize \prelimfootlabel\hspace{3em}H2 MATHEMATICS 9758}}
\fancypagestyle{prelim}{\fancyhf{}%
  \chead{\bfseries\thepage}%
  \lfoot{\pp@footer}%
  \rfoot{\ifnum\value{page}<\getpagerefnumber{LastPage}\relax\textbf{[Turn over}\fi}%
  \renewcommand{\headrulewidth}{0pt}\renewcommand{\footrulewidth}{0pt}}
\fancypagestyle{prelimcover}{\fancyhf{}%
  \rfoot{\textbf{[Turn over}}%
  \renewcommand{\headrulewidth}{0pt}\renewcommand{\footrulewidth}{0pt}}
\pagestyle{prelim}

% ---- marks: [n] regular weight, right-aligned at the end of the part ----------
\renewcommand{\mk}[1]{{\unskip\nobreak\hfil\penalty50\hskip1.5em\hbox{}\nobreak\hfill[#1]%
  \parfillskip=0pt\finalhyphendemerits=0\par}}

% ---- questions: bold number at the left margin, text indented; grey [tag] straight after the number (3B)
\newsavebox{\pp@qmeas}
\newcommand{\pp@qlist}[1]{%
  \begin{list}{}{\setlength\leftmargin{2.6em}\setlength\labelwidth{2.1em}\setlength\labelsep{0.5em}%
     \setlength\topsep{0pt}\setlength\itemsep{0pt}\setlength\parsep{0.9ex}%
     \renewcommand\makelabel[1]{\textbf{##1}\hfil}}%
  \item[\thediq]\@endparpenalty\@M\setbox0=\hbox{\di@tag{#1}}\ifdim\wd0=0pt\else{\small\textcolor{black!60}{[\di@tag{#1}]}}\par\nobreak\fi\csname diQbody@#1\endcsname\par\penalty\@M
  \end{list}}
% a question starts on a new page unless it fits whole on the current one
% (questions taller than a page may break, but only between parts)
\renewcommand{\putQ}[1]{\stepcounter{diq}%
  \expandafter\xdef\csname diLab@#1\endcsname{\thediq}%
  \par\setbox\pp@qmeas=\vbox{\pp@qlist{#1}}%
  \ifdim\dimexpr\ht\pp@qmeas+\dp\pp@qmeas\relax<0.95\textheight
    \needspace{\dimexpr\ht\pp@qmeas+\dp\pp@qmeas\relax}\else\needspace{14\baselineskip}\fi
  \zsavepos{q-s\thediq}\pp@qlist{#1}\par\zsavepos{q-e\thediq}\vspace{5ex plus 2ex minus 1ex}}
% part labels (a)/(i) bold, indented under the question text
\renewenvironment{qparts}[1][2.6em]{\begin{list}{}{\setlength\leftmargin{#1}%
  \setlength\labelwidth{\dimexpr#1-0.5em\relax}\setlength\labelsep{0.5em}%
  \setlength\itemsep{1.4ex}\setlength\topsep{1.2ex}\setlength\parsep{0.8ex}%
  \renewcommand\makelabel[1]{\textbf{##1}\hfil}}}{\end{list}}
\makeatother

% ---- cover page --------------------------------------------------------------
\newcommand{\prelimcover}[3]{% #1 total marks, #2 duration, #3 date line
\thispagestyle{prelimcover}%
{\coverfont\fontsize{10}{12}\selectfont\linespread{1}\selectfont
\noindent\begin{minipage}[t]{\linewidth}
  {\bfseries \prelimtitle}\\[0.2ex]
  {\bfseries Compiled by Maximus Pung Jun}\\[1.6ex]
  Higher 2
\end{minipage}\par\vspace{2ex}
\noindent\rule{\linewidth}{1.2pt}\par\vspace{0.4ex}
\noindent{\Large\bfseries MATHEMATICS}\hfill{\Large\bfseries 9758}\par\vspace{1.2ex}
\noindent \prelimsubject\hfill\textbf{#3}\\
QUESTION PAPER\hfill\textbf{#2}\par\vspace{2.5ex}
\noindent\begin{tabular}{@{}p{46mm}l@{}}Additional Materials: & Writing paper\\ & List of Formulae and Results (MF27)\end{tabular}\par\vspace{2ex}
\noindent\rule{\linewidth}{0.6pt}\par\vspace{0.6ex}
{\bfseries READ THESE INSTRUCTIONS FIRST}\par\vspace{2ex}
\setlength{\parskip}{0pt}
Write your name and class on all the work you hand in.\\
Write in dark blue or black pen.\\
You may use an HB pencil for any diagrams or graphs.\\
Do not use staples, paper clips, glue or correction fluid.\par\vspace{2ex}
Answer \textbf{all} the questions.\\
Give non-exact numerical answers correct to 3 significant figures, or 1 decimal place in the case of angles in degrees, unless a different level of accuracy is specified in the question.\\
You are expected to use an approved graphing calculator.\\
Unsupported answers from a graphing calculator are allowed unless a question specifically states otherwise.\\
Where unsupported answers from a graphing calculator are \textbf{not} allowed in a question, you are required to present the mathematical steps using mathematical notations and not calculator commands.\\
You are reminded of the need for clear presentation in your answers.\par\vspace{2ex}
The number of marks is given in brackets [\ ] at the end of each question or part question.\\
The total number of marks for this paper is #1.
\vfill
\noindent\rule{\linewidth}{0.6pt}\par\vspace{0.4ex}
\centerline{This document consists of \textbf{\pageref{LastPage}} printed pages.}\par\vspace{1ex}
}\clearpage}

% ---- solutions header block (page 1 of the worked solutions; same cover style) ----
\newcommand{\prelimsolhead}[1]{% #1 total marks
{\coverfont\fontsize{10}{12}\selectfont\linespread{1}\selectfont
\noindent{\bfseries \prelimtitle}\\[0.2ex]{\bfseries Compiled by Maximus Pung Jun}\\[1.6ex]Higher 2\par\vspace{2ex}
\noindent\rule{\linewidth}{1.2pt}\par\vspace{0.4ex}
\noindent{\Large\bfseries MATHEMATICS}\hfill{\Large\bfseries 9758}\par\vspace{1.2ex}
\noindent \prelimsubject\hfill\textbf{2026}\\
WORKED SOLUTIONS\hfill\textbf{Total marks: #1}\par\vspace{1ex}
\noindent\rule{\linewidth}{0.6pt}\par}\vspace{4ex}}
% house_sol.tex: HOUSE COPY of definite_integrals/common/ri_solstyle.tex (unchanged)
% =====================================================================
% ri_solstyle.tex: worked solutions in the layout of the RI 2026 Year 6 Prelim
% "Solutions with comments" (grail.moe 26612/26613), loaded AFTER prelim_style.
%   * page 1: bold school-style block + bold underlined title; later pages: running
%     title top-left with a rule under it; footer "Page N of M" bottom right (Calibri).
%   * each question: prelim-format question text (bold number, marks right, source tag
%     bottom right), then a bordered table:  Solution | Comments  (bold header row);
%     body rows  label + [marks] (bold, narrow left column) | working | comments,
%     a full-width rule between parts, working left-aligned (fleqn).
% Rows come from the existing sources: \pt in qparts (Ch10 HW), \slabel (Ch10 Test),
% spart (Ch11). Marks under each label are taken from the question's \mk.
% =====================================================================
\makeatletter
\usepackage{lastpage}
\usepackage{adjustbox}
\usepackage{zref-savepos,zref-abspage}
\zref@addprop{savepos}{abspage}
\geometry{a4paper, top=22mm, bottom=24mm, left=22mm, right=22mm, headsep=8mm, footskip=12mm, headheight=24pt}
\providecommand{\risoltitle}{\prelimtitle: Solutions with comments}
\providecommand{\risolschool}{Compiled by Maximus Pung Jun}
\fancypagestyle{risol}{\fancyhf{}%
  \lhead{\fontsize{11}{12}\selectfont 2026 \risoltitle\\[-1.2ex]\rule{\textwidth}{0.6pt}}%
  \rfoot{\calibrifont\fontsize{11}{13}\selectfont Page \textbf{\thepage} of \textbf{\pageref*{LastPage}}}%
  \renewcommand{\headrulewidth}{0pt}\renewcommand{\footrulewidth}{0pt}}
\fancypagestyle{risolfirst}{\fancyhf{}%
  \rfoot{\calibrifont\fontsize{11}{13}\selectfont Page \textbf{\thepage} of \textbf{\pageref*{LastPage}}}%
  \renewcommand{\headrulewidth}{0pt}\renewcommand{\footrulewidth}{0pt}}
\pagestyle{risol}
\newif\ifri@fresh
\newcommand{\risolfront}{\thispagestyle{risolfirst}\global\ri@freshtrue%
  {\noindent\bfseries\linespread{1}\selectfont \risolschool\\ H2 Mathematics (9758)\\ 2026\par}%
  \vspace{4ex}{\noindent\bfseries\underline{2026 \risoltitle}\par}\vspace{3ex}}
% ---- table geometry (RI: label 42pt | working 284pt | comments 142pt of 468pt) ----
\newlength{\ri@L}\setlength{\ri@L}{15mm}
\newlength{\ri@C}\setlength{\ri@C}{40mm}
\newlength{\ri@pad}\setlength{\ri@pad}{1.8mm}
\newif\ifri@sol \newif\ifri@first
\newcount\ri@depth
\def\ri@outer{}
\newtcolorbox{risol}{enhanced, breakable, sharp corners, frame hidden, colback=white,
  boxrule=0pt, boxsep=0pt, left=\dimexpr\ri@L+\ri@pad\relax, right=\dimexpr\ri@C+\ri@pad\relax,
  top=1.2mm, bottom=2.5mm, toptitle=0.6mm, bottomtitle=0.6mm, lefttitle=\ri@pad, righttitle=0pt,
  colbacktitle=white, coltitle=black, fonttitle=\bfseries,
  title={Solution\hfill\makebox[\ri@C][l]{\hspace{\ri@pad}Comments}},
  before skip=2ex, after skip=2ex,
  overlay unbroken and first={\draw[line width=0.5pt] (frame.south west) rectangle (frame.north east);
     \draw[line width=0.5pt] (title.south west) -- (title.south east);
     \draw[line width=0.5pt] ([xshift=\ri@L]interior.north west) -- ([xshift=\ri@L]frame.south west);
     \draw[line width=0.5pt] ([xshift=-\ri@C]frame.north east) -- ([xshift=-\ri@C]frame.south east);},
  overlay middle and last={\draw[line width=0.5pt] (frame.south west) rectangle (frame.north east);
     \draw[line width=0.5pt] ([xshift=\ri@L]frame.north west) -- ([xshift=\ri@L]frame.south west);
     \draw[line width=0.5pt] ([xshift=-\ri@C]frame.north east) -- ([xshift=-\ri@C]frame.south east);},
  after upper={\ri@endrow\protected@write\@auxout{}{\string\ri@settotal{\ri@qid}{\the\ri@sumH}}},
  before upper={\global\ri@sumH=14mm\ri@soltrue\ri@firsttrue\ri@depth=0 \parindent=0pt \parskip=0.6ex
     \predisplaypenalty=100 
     \linespread{1.12}\selectfont\raggedright\@fleqntrue\@mathmargin=0pt
     \abovedisplayskip=5pt plus 2pt\belowdisplayskip=5pt plus 2pt
     \abovedisplayshortskip=3pt plus 1pt\belowdisplayshortskip=4pt plus 2pt
     \ifcsname ri@hp@\ri@qid\endcsname\else\ri@row{}\fi}}
\def\ri@qid{}
% one row: full-width rule (glue leaders, so it vanishes at a page break where the frame
% edge takes its place; not before the first row), then bold label + bold [marks]
% in the left column, on the baseline of the first line of the working
\ExplSyntaxOn
\tl_new:N \g__ri_tmp_tl
\cs_new_protected:Npn \ri_setkey:Nn #1#2
  { \group_begin: \cs_set_eq:NN \chg \use_ii:nn \cs_set_eq:NN \chgm \use_ii:nn
    \cs_set_eq:NN \textbf \use:n \cs_set_eq:NN \mathrm \use:n
    \tl_gset:Nx \g__ri_tmp_tl {#2} \group_end:
    \tl_set_eq:NN #1 \g__ri_tmp_tl \tl_remove_all:Nn #1 {~} }
\cs_set_eq:NN \ri@setkey \ri_setkey:Nn
\ExplSyntaxOff
% ---- page-break rules (positions from the previous run, zref labels in the .aux) ----
% * a row (one part) never breaks across pages: if it does not fit on the rest of the
%   page it starts the next page (forced break between rows); only a row taller than a
%   whole page may break inside;
% * a whole solution that fits on one page is kept together (needspace before the table);
%   a longer one starts where at least its first row fits;
% * a row is at least as tall as its label + [marks].
\newcount\ri@n \newbox\ri@labbox \def\ri@open{0}
\newdimen\ri@ybot \newdimen\ri@ytop \newdimen\ri@h \newdimen\ri@sumH
\AtBeginDocument{\ri@ybot=\dimexpr\paperheight-1in-\voffset-\topmargin-\headheight-\headsep-\textheight\relax
  \ri@ytop=\dimexpr\ri@ybot+\textheight\relax}
\def\ri@pg#1{\zref@extractdefault{#1}{abspage}{-1}}
% \ri@rowh{n}: height of row n in the previous run -> \ri@h (0pt if unknown)
\newcommand{\ri@rowh}[1]{\ri@h=0pt
  \ifnum\ri@pg{ri-s#1}>0 \ifnum\ri@pg{ri-e#1}>0
    \ifnum\ri@pg{ri-s#1}=\ri@pg{ri-e#1}\relax \ri@h=\dimexpr\zposy{ri-s#1}sp-\zposy{ri-e#1}sp\relax
    \else \ri@h=\dimexpr\zposy{ri-s#1}sp-\ri@ybot+\ri@ytop-\zposy{ri-e#1}sp-6mm\relax\fi\fi\fi}
\newcommand{\ri@setfirst}[2]{\expandafter\gdef\csname ri@first@#1\endcsname{#2}}
\newcommand{\ri@setkept}[1]{\expandafter\gdef\csname ri@kept@#1\endcsname{}}
\newcommand{\ri@settotal}[2]{\expandafter\gdef\csname ri@total@#1\endcsname{#2}}
\newcommand{\ri@endrow}{\ifnum\ri@open=1 \par\zsavepos{ri-e\the\ri@n}%
  \ifnum\ri@pg{ri-s\the\ri@n}=\ri@pg{ri-e\the\ri@n}\relax
    \dimen@=\dimexpr\ri@curH sp+1.6mm-\zposy{ri-s\the\ri@n}sp+\zposy{ri-e\the\ri@n}sp\relax
    \ifdim\dimen@>0pt \nobreak\kern\dimen@\penalty\z@\fi\fi
  \ri@rowh{\the\ri@n}\global\advance\ri@sumH\dimexpr\ri@h+3.5mm\relax
  \gdef\ri@open{0}\fi}
\newcommand{\ri@row}[1]{\par\ri@endrow\global\advance\ri@n\@ne\gdef\ri@open{1}\@beginparpenalty\@M
  \ifri@first\global\ri@firstfalse\ri@firstfalse
    \protected@write\@auxout{}{\string\ri@setfirst{\ri@qid}{\the\ri@n}}%
  \else
    \zsavepos{ri-p\the\ri@n}%
    % push the row to the next page if it would not fit on the rest of this one
    \ri@rowh{\the\ri@n}%
    \ifnum\ri@pg{ri-p\the\ri@n}>0 \ifdim\ri@h>0pt
      \ifdim\dimexpr\ri@h+4mm\relax<\dimexpr\textheight-10mm\relax
        \ifdim\dimexpr\ri@h+7mm\relax>\dimexpr\zposy{ri-p\the\ri@n}sp-\ri@ybot\relax
          \penalty-\@M\fi\fi\fi\fi
    \vskip0.6ex\nointerlineskip
    \leaders\hbox{\hskip-\dimexpr\ri@L+\ri@pad\relax\vrule width\dimexpr\linewidth+\ri@L+\ri@C+2\ri@pad\relax height0.5pt depth0pt}\vskip0.5pt
    \nobreak\vskip1.4mm\fi
  \ri@setkey\ri@key{\ri@qid#1}\def\ri@lab{#1}%
  \setbox\ri@labbox=\vtop{\hsize=\dimexpr\ri@L-1.6mm\relax\bfseries\raggedright\parskip=0pt
      \strut\ri@lab\ifcsname ri@mk@\ri@key\endcsname\ifx\ri@lab\@empty\else\\\fi
      [\csname ri@mk@\ri@key\endcsname]\fi\strut\par}%
  \xdef\ri@curH{\number\dimexpr\ht\ri@labbox+\dp\ri@labbox\relax}%
  \nointerlineskip\zsavepos{ri-s\the\ri@n}\nointerlineskip\nobreak
  \vbox to 0pt{\hbox{\hskip-\dimexpr\ri@L+\ri@pad-1.2mm\relax\box\ri@labbox}\vss}\nobreak\prevdepth=\dp\strutbox
  \vskip-\parskip\ignorespaces}
% ---- qparts / \pt: question mode records labels (for the marks); solution mode makes rows
\NewCommandCopy\ri@qparts@begin\qparts \let\ri@qparts@end\endqparts
\renewenvironment{qparts}[1][2.6em]{\advance\ri@depth\@ne
  \ifri@sol\else\ri@qparts@begin[#1]\fi}{\ifri@sol\par\else\ri@qparts@end\fi}
\def\ri@qpartsname{qparts}
\renewcommand{\pt}[1]{\ifnum\ri@depth<2 \gdef\ri@outer{#1}\gdef\ri@cur{#1}\else\expandafter\gdef\expandafter\ri@cur\expandafter{\ri@outer#1}\fi
  \ifri@sol\expandafter\ri@solpt\else\expandafter\ri@qpt\fi{#1}}
\def\ri@qpt#1{\expandafter\gdef\csname ri@hp@\ri@qid\endcsname{}\item[#1]}
\def\ri@solpt#1{\@ifnextchar\begin\ri@ptpeek{\ri@row{\ri@cur}}}
\def\ri@ptpeek\begin#1{\def\ri@tmp{#1}\ifx\ri@tmp\ri@qpartsname
   \expandafter\@firstoftwo\else\expandafter\@secondoftwo\fi
  {\begin{qparts}}{\ri@row{\ri@cur}\begin{#1}}}
% marks: record [n] against the current question/part
\let\ri@mk\mk
\renewcommand{\mk}[1]{\ri@setkey\ri@k{\ri@qid\ri@cur}\ri@setkey\ri@mv{#1}\expandafter\xdef\csname ri@mk@\ri@k\endcsname{\ri@mv}\ri@mk{#1}}
% \slabel (Ch10 Test answers): letter labels set the outer part, roman ones are appended;
% consecutive \slabel calls merge into one row label
\def\ri@romans{(i)(ii)(iii)(iv)(v)(vi)(vii)}
\def\ri@firstpart(#1)#2\@nil{(#1)}
\def\slabel#1{\ifri@sol\expandafter\ri@solslabel\else\expandafter\ri@qslabel\fi{#1}}
\def\ri@qslabel#1{\par\smallskip\noindent\textbf{#1}\par\nobreak}
\def\ri@solslabel#1{\@expandtwoargs\in@{#1}{\ri@romans}%
  \ifin@\expandafter\gdef\expandafter\ri@cur\expandafter{\ri@outer#1}\else\gdef\ri@cur{#1}\xdef\ri@outer{\ri@firstpart#1\@nil}\fi
  \@ifnextchar\slabel{}{\ri@row{\ri@cur}}}
% Ch11 spart{label}: a row; soln: the table
\def\spart#1{\gdef\ri@cur{#1}\ri@row{#1}}\def\endspart{\par}
\def\soln{\@ifnextchar[{\ri@solnopt}{\risol}}\def\ri@solnopt[#1]{\risol}\def\endsoln{\endrisol}
% ---- working that started on the label line is left-aligned in the working column
\AtBeginDocument{\@ifundefined{topalign}{}{%
  \renewenvironment{topalign}{\par\noindent$\displaystyle\begin{aligned}[t]}{\end{aligned}$\par\smallskip}%
  \renewcommand{\topdisplay}[1]{\par\noindent$\displaystyle #1$\par\smallskip}}}
% ---- Ch11 side-by-side (text | sketch): the sketch shrinks to fit its half of the working column
\AtBeginDocument{\@ifundefined{sbsplain}{}{%
  \renewenvironment{sbsplain}[1]{%
    \sbs@start\begin{minipage}[t]{#1\linewidth}\vspace{0pt}\sbs@shift\raggedright%
    \def\sbsnext{\end{minipage}\hfill\begin{minipage}[t]{\dimexpr0.97\linewidth-#1\linewidth\relax}\vspace{0pt}\sbs@shift\centering
      \begin{adjustbox}{max width=\linewidth}}%
  }{\end{adjustbox}\end{minipage}\par}}}
% ---- figures in the working column shrink to fit it (only if wider), colours kept
\newcommand{\rifitfig}[1]{\expandafter\let\csname ri@orig@\string#1\endcsname#1%
  \def#1{\@ifnextchar\bgroup{\ri@fitarg#1}{\ri@fitnoarg#1}}}
\def\ri@fitarg#1#2{\ifri@sol\adjustbox{max width=\linewidth}{\csname ri@orig@\string#1\endcsname{#2}}\else\csname ri@orig@\string#1\endcsname{#2}\fi}
\def\ri@fitnoarg#1{\ifri@sol\adjustbox{max width=\linewidth}{\csname ri@orig@\string#1\endcsname}\else\csname ri@orig@\string#1\endcsname\fi}
\@ifundefined{solfig}{}{\renewcommand{\solfig}[2][]{\par\begin{center}\adjustbox{max width=\linewidth}{\begin{tikzpicture}\begin{axis}[qfig, width=6.6cm, height=5cm, #1]#2\end{axis}\end{tikzpicture}}\end{center}}}
% \ricomment{text}: text in the Comments column of the current row (call right after \pt,
% \slabel or \spart). The row is made at least as tall as the comment.
\newbox\ri@cmtbox
\newcommand{\ricomment}[1]{\par\setbox\ri@cmtbox=\vtop{\hsize=\dimexpr\ri@C-2\ri@pad\relax
    \parskip=0pt\raggedright\small{\color{black!60}\strut#1\strut\par}}%
  \ifdim\dimexpr\ht\ri@cmtbox+\dp\ri@cmtbox\relax>\ri@curH sp\xdef\ri@curH{\number\dimexpr\ht\ri@cmtbox+\dp\ri@cmtbox\relax}\fi
  \nointerlineskip\vbox to 0pt{\hbox to\linewidth{\hskip\dimexpr\linewidth+2\ri@pad\relax\box\ri@cmtbox\hss}\vss}\nobreak\prevdepth=\dp\strutbox
  \vskip-\parskip\ignorespaces}
% \putTFlong{n}: the True/False question n with its answer table allowed to break across
% pages (the body of \defQ{n} must be: intro text, then \begin{tftable}...\end{tftable});
% the solutions copy prints it with tflongtable.
\renewcommand{\putTFlong}[1]{\ri@newq\begingroup\let\tftable\tflongtable\let\endtftable\endtflongtable
  \stepcounter{diq}\expandafter\xdef\csname diLab@#1\endcsname{\thediq}%
  \par\noindent\hbox to 2.6em{\textbf{\thediq}\hfil}\setbox0=\hbox{\di@tag{#1}}\ifdim\wd0=0pt\else{\small\textcolor{black!60}{[\di@tag{#1}]}}\par\noindent\hskip2.6em\fi\ignorespaces\csname diQbody@#1\endcsname\par
  \endgroup\vspace{3ex}}
% ---- question + solution; each question starts on a new page (as RI)
\def\ri@qid{}
\let\ri@putQ\putQ
\renewcommand{\putQ}[1]{\gdef\ri@qid{#1}\global\ri@depth=0 \gdef\ri@cur{}\ri@putQ{#1}}
% \risolQ{n}{body macro name}: question n + its solution body in the table;
% \risolQraw: body already contains its own soln environment (Ch11 Test)
\newcommand{\ri@newq}{\ifri@fresh\global\ri@freshfalse\else\clearpage\fi}
% keep the whole table together if it fits on a page; otherwise make sure at least its
% header and first row fit before starting it
\newcommand{\ri@keep}[1]{\par
  \ifcsname ri@total@#1\endcsname
    % hysteresis: a table kept whole in the previous run stays whole up to textheight-4.5mm,
    % so a total near the limit cannot flip between runs
    \ifdim\csname ri@total@#1\endcsname<\dimexpr\textheight-\ifcsname ri@kept@#1\endcsname4.5\else8\fi mm\relax
      \protected@write\@auxout{}{\string\ri@setkept{#1}}%
      \needspace{\dimexpr\csname ri@total@#1\endcsname+4mm\relax}%
    \else\ifcsname ri@first@#1\endcsname
      \ri@rowh{\csname ri@first@#1\endcsname}%
      \ifdim\ri@h<\dimexpr\textheight-20mm\relax\needspace{\dimexpr\ri@h+18mm\relax}\fi\fi\fi\fi}
\newcommand{\risolQ}[2]{\ri@newq\putQ{#1}\vspace{-3ex}\ri@keep{#1}\begin{risol}\csname #2@#1\endcsname\par\end{risol}}
\newcommand{\risolQraw}[2]{\ri@newq\putQ{#1}\vspace{-3ex}\ri@keep{#1}\csname #2@#1\endcsname\par}
\renewcommand{\putQS}[1]{\risolQ{#1}{diS}}
% end marker only if it fits on the last page (no page holding just the marker)
\renewcommand{\diend}[1][END OF SOLUTIONS]{\par\ifdim\dimexpr\pagegoal-\pagetotal\relax<4\baselineskip\else\vspace{1.5ex}\begin{center}\textbf{#1}\end{center}\fi}
\makeatother
\tfanswerstrue\fi
\fi
, set:
%   \def\housemode{notes}    lecture notes: book class, fncychap Lenny chapter page, boxes
%   \def\housemode{hw}       homework / additional / workshop QUESTIONS (+ Answers section)
%   \def\housemode{sol}      worked SOLUTIONS of a hw (RI "Solution | Comments" table)
%   \def\housemode{test}     test question paper (prelim cover, prelim layout)
%   \def\housemode{testsol}  test worked solutions (prelim layout + RI table)
%   \def\housechapterno{10}  (notes: chapter number; optional)
%   \def\housenotesLM{1}     (notes only, optional) keep Latin Modern as in the current
%                            DI/DE notes instead of Times New Roman + TeX Gyre Termes Math
% then \housechapter{Definite Integrals} and \ditle{Homework 1}.
% Document class (put BEFORE % =====================================================================
% house.tex: ONE preamble for every Maximus Pung Jun H2 Math document,
% modelled on Ch10 Definite Integrals + Ch11 Differential Equations (see SPEC.md).
%
% Before % =====================================================================
% house.tex: ONE preamble for every Maximus Pung Jun H2 Math document,
% modelled on Ch10 Definite Integrals + Ch11 Differential Equations (see SPEC.md).
%
% Before \input{house}, set:
%   \def\housemode{notes}    lecture notes: book class, fncychap Lenny chapter page, boxes
%   \def\housemode{hw}       homework / additional / workshop QUESTIONS (+ Answers section)
%   \def\housemode{sol}      worked SOLUTIONS of a hw (RI "Solution | Comments" table)
%   \def\housemode{test}     test question paper (prelim cover, prelim layout)
%   \def\housemode{testsol}  test worked solutions (prelim layout + RI table)
%   \def\housechapterno{10}  (notes: chapter number; optional)
%   \def\housenotesLM{1}     (notes only, optional) keep Latin Modern as in the current
%                            DI/DE notes instead of Times New Roman + TeX Gyre Termes Math
% then \housechapter{Definite Integrals} and \ditle{Homework 1}.
% Document class (put BEFORE \input{house}):
%   notes: \documentclass[a4paper]{book}
%   others: \documentclass[a4paper,11pt,oneside]{book}
% A HW file uses \providecommand\housemode{hw}; its solutions file has NO \documentclass:
%   \def\housemode{sol}\def\diquestionsonly{1}\input{HW1}  (as DI/DE HWn_solutions.tex)
% Compile: TEXINPUTS=.:/workspace/house_format: xelatex (2 runs; solutions until .aux stable)
% =====================================================================
\providecommand\housemode{hw}
\providecommand\housedir{/workspace/house_format}
\newcommand\housechap{}
\newcommand\housechapter[1]{\renewcommand\housechap{#1}}
\def\housexnotes{notes}\def\housexhw{hw}\def\housexsol{sol}\def\housextest{test}\def\housextestsol{testsol}
\edef\housexm{\housemode}
\input{house_base}
\input{house_local}
\ifx\housexm\housexnotes
  \ifdefined\housenotesLM\else\input{house_fonts}\fi
  \input{house_tf}
  \providecommand{\units}[1]{\ \text{units}^{#1}}
  \raggedbottom
\else
  \input{house_fonts}
  \input{house_hw}
  \input{house_solmacros}
  \input{house_tf}
  \ifx\housexm\housexsol \def\prelimkeepcolour{1}\input{house_prelim}\input{house_sol}\tfanswerstrue\fi
  \ifx\housexm\housextest \input{house_prelim}\fi
  \ifx\housexm\housextestsol \def\prelimkeepcolour{1}\input{house_prelim}\input{house_sol}\tfanswerstrue\fi
\fi
, set:
%   \def\housemode{notes}    lecture notes: book class, fncychap Lenny chapter page, boxes
%   \def\housemode{hw}       homework / additional / workshop QUESTIONS (+ Answers section)
%   \def\housemode{sol}      worked SOLUTIONS of a hw (RI "Solution | Comments" table)
%   \def\housemode{test}     test question paper (prelim cover, prelim layout)
%   \def\housemode{testsol}  test worked solutions (prelim layout + RI table)
%   \def\housechapterno{10}  (notes: chapter number; optional)
%   \def\housenotesLM{1}     (notes only, optional) keep Latin Modern as in the current
%                            DI/DE notes instead of Times New Roman + TeX Gyre Termes Math
% then \housechapter{Definite Integrals} and \ditle{Homework 1}.
% Document class (put BEFORE % =====================================================================
% house.tex: ONE preamble for every Maximus Pung Jun H2 Math document,
% modelled on Ch10 Definite Integrals + Ch11 Differential Equations (see SPEC.md).
%
% Before \input{house}, set:
%   \def\housemode{notes}    lecture notes: book class, fncychap Lenny chapter page, boxes
%   \def\housemode{hw}       homework / additional / workshop QUESTIONS (+ Answers section)
%   \def\housemode{sol}      worked SOLUTIONS of a hw (RI "Solution | Comments" table)
%   \def\housemode{test}     test question paper (prelim cover, prelim layout)
%   \def\housemode{testsol}  test worked solutions (prelim layout + RI table)
%   \def\housechapterno{10}  (notes: chapter number; optional)
%   \def\housenotesLM{1}     (notes only, optional) keep Latin Modern as in the current
%                            DI/DE notes instead of Times New Roman + TeX Gyre Termes Math
% then \housechapter{Definite Integrals} and \ditle{Homework 1}.
% Document class (put BEFORE \input{house}):
%   notes: \documentclass[a4paper]{book}
%   others: \documentclass[a4paper,11pt,oneside]{book}
% A HW file uses \providecommand\housemode{hw}; its solutions file has NO \documentclass:
%   \def\housemode{sol}\def\diquestionsonly{1}\input{HW1}  (as DI/DE HWn_solutions.tex)
% Compile: TEXINPUTS=.:/workspace/house_format: xelatex (2 runs; solutions until .aux stable)
% =====================================================================
\providecommand\housemode{hw}
\providecommand\housedir{/workspace/house_format}
\newcommand\housechap{}
\newcommand\housechapter[1]{\renewcommand\housechap{#1}}
\def\housexnotes{notes}\def\housexhw{hw}\def\housexsol{sol}\def\housextest{test}\def\housextestsol{testsol}
\edef\housexm{\housemode}
\input{house_base}
\input{house_local}
\ifx\housexm\housexnotes
  \ifdefined\housenotesLM\else\input{house_fonts}\fi
  \input{house_tf}
  \providecommand{\units}[1]{\ \text{units}^{#1}}
  \raggedbottom
\else
  \input{house_fonts}
  \input{house_hw}
  \input{house_solmacros}
  \input{house_tf}
  \ifx\housexm\housexsol \def\prelimkeepcolour{1}\input{house_prelim}\input{house_sol}\tfanswerstrue\fi
  \ifx\housexm\housextest \input{house_prelim}\fi
  \ifx\housexm\housextestsol \def\prelimkeepcolour{1}\input{house_prelim}\input{house_sol}\tfanswerstrue\fi
\fi
):
%   notes: \documentclass[a4paper]{book}
%   others: \documentclass[a4paper,11pt,oneside]{book}
% A HW file uses \providecommand\housemode{hw}; its solutions file has NO \documentclass:
%   \def\housemode{sol}\def\diquestionsonly{1}\input{HW1}  (as DI/DE HWn_solutions.tex)
% Compile: TEXINPUTS=.:/workspace/house_format: xelatex (2 runs; solutions until .aux stable)
% =====================================================================
\providecommand\housemode{hw}
\providecommand\housedir{/workspace/house_format}
\newcommand\housechap{}
\newcommand\housechapter[1]{\renewcommand\housechap{#1}}
\def\housexnotes{notes}\def\housexhw{hw}\def\housexsol{sol}\def\housextest{test}\def\housextestsol{testsol}
\edef\housexm{\housemode}
% house_base.tex: HOUSE COPY of definite_integrals/common/ws_preamble_build.tex
% (= Maximus's own preamble.tex + the DE v3 / AoI additions). Only change: GC font path. [HOUSE]
% WORKSHEET COPY of ../notes/preamble_build.tex (2026-10-09 version), unchanged.
% preamble_build.tex: copy of Maximus's preamble.tex for the DE v3 build.
% The blocks marked [BUILD COPY CHANGE] are the earlier build fixes.
% The blocks marked [V3 ADDITION] are new in this version (example box,
% spacing). Everything else is unchanged from preamble.tex.

% For margins
\usepackage[a4paper, margin = 20mm]{geometry}

\usepackage[Lenny]{fncychap}
% Rejne, Glenn, Lenny

\usepackage{etoc}

% [V3 ADDITION] cancel: used to show the dx's "cancelling" in Section 11.3
\usepackage{cancel}
\renewcommand{\CancelColor}{\color{red}}

% For enumeration & math symbols
\usepackage{enumerate}
\usepackage{amsmath,amsfonts,amssymb,amsthm,graphicx,lipsum}

% For diagrams, plots and boxes
\usepackage{pgfplots}
\pgfplotsset{compat=1.18}
\usepackage{tikz}
\usepackage[most]{tcolorbox}
\tcbuselibrary{skins}
\usepackage{xcolor}
\pgfplotsset{
    standard/.style={clip = false,
    axis line style = thick,
    trig format=rad,
    axis lines = middle,
    every axis x label/.style={at={(current axis.right of origin)},anchor= west},
    every axis y label/.style={at={(current axis.above origin)},anchor=south}
    }
}

% For GC font
\usepackage{fontspec}
% [BUILD COPY CHANGE 1] The TI-84 key font is loaded only if the file exists
% (checked in ./Fonts/... and ../Fonts/...). Otherwise \ti{} falls back to a
% small boxed sans-serif key. Known letters of the TI font are mapped to their
% key names in the fallback: o = [y=], z = [mode], f = [L3].
\newcommand{\tikey}[1]{\tcbox[on line, enhanced, boxsep=0pt, left=2pt, right=2pt,
  top=1pt, bottom=1pt, boxrule=0.5pt, arc=1.5pt, colback=white,
  colframe=black!70]{\sffamily\footnotesize\strut #1}}
\ExplSyntaxOn
\NewDocumentCommand{\tifallback}{m}{
  \str_case:nnF {#1}
    { {o}{\tikey{y=}} {z}{\tikey{mode}} {f}{\tikey{L3}} }
    { \tikey{#1} }
}
\ExplSyntaxOff
\def\tifontcase{0}
\IfFileExists{Fonts/TI-84 Plus CE Key Fonts/TI84PlusCEKeys.ttf}
  {\def\tifontcase{1}}{}
\IfFileExists{../Fonts/TI-84 Plus CE Key Fonts/TI84PlusCEKeys.ttf}
  {\def\tifontcase{2}}{}
\ifcase\tifontcase\relax
  % 0: font not found, use boxed fallback keys
  \newcommand{\ti}[1]{\tifallback{#1}}
\or
  % 1: original location (as in preamble.tex)
  \newfontfamily\tifont[
    Path = {Fonts/TI-84 Plus CE Key Fonts/},
    UprightFont = {TI84PlusCEKeys.ttf}
  ]{}
  \newcommand{\ti}[1]{{\tifont #1}}
\or
  % 2: one level up (when compiling inside draft-ch1/)
  \newfontfamily\tifont[
    Path = {../Fonts/TI-84 Plus CE Key Fonts/},
    UprightFont = {TI84PlusCEKeys.ttf}
  ]{}
  \newcommand{\ti}[1]{{\tifont #1}}
\fi

% [AoI items 17/18/37] GC DISPLAY FONT. Every GC key, command, menu item and
% screen line goes through \gcfont, so the font can be swapped HERE in one place.
%   \gc{fnInt(}      command / menu item / stored variable in the display font
%   \gckey{2nd}      a calculator key (display font in a small key outline)
%   \gcscreen{...}   a calculator screen (display font on a grey-green LCD panel)
% The font file lives in ../Fonts (shared by notes, slides and worksheet).
\def\gcfontfile{VT323-Regular.ttf}
% [HOUSE] GC font: always the absolute copy in \housedir/Fonts (relative Path= breaks under TEXINPUTS)
\providecommand\housedir{/workspace/house_format}
\IfFileExists{\housedir/Fonts/\gcfontfile}{\newfontfamily\gcfont[Path=\housedir/Fonts/, Scale=1.32]{\gcfontfile}}{%
  \newcommand\gcfont{\ttfamily}}
\newcommand{\gc}[1]{\mbox{\gcfont #1}}
\newcommand{\gckey}[1]{\tcbox[on line, enhanced, boxsep=0pt, left=2pt, right=2pt,
  top=0.6pt, bottom=0.2pt, boxrule=0.5pt, arc=1.5pt, colback=black!4,
  colframe=black!65]{\gcfont\strut #1}}
\definecolor{lcd}{RGB}{214,222,204}
\newcommand{\gcscreen}[1]{\tcbox[on line, enhanced, boxsep=1pt, left=4pt, right=4pt,
  top=2pt, bottom=2pt, boxrule=1.2pt, arc=1pt, colback=lcd, colframe=black!70]{%
  \gcfont\begin{tabular}{@{}l@{}}#1\end{tabular}}}
% integral template as shown on the TI-84 Plus CE (MathPrint): \gcint{0}{2}{e^(-X²)}{X}
\newcommand{\gcint}[4]{{\gcfont\raisebox{-0.35ex}{\Large ∫}\hspace{-0.15em}%
  \raisebox{-0.2ex}{\begin{tabular}{@{}l@{}}\footnotesize #2\\[-1.3ex]\footnotesize #1\end{tabular}}%
  \hspace{0.1em}(#3)\mathrm{d}#4}}
% old key macro kept as an alias so nothing breaks
\renewcommand{\ti}[1]{\gckey{#1}}

% SI units
\usepackage{siunitx}

% setspace
%\usepackage{setspace}
%\onehalfspacing

% For hyperlinks within and out of document
\usepackage{hyperref}
\hypersetup{
    colorlinks = true,
    linkcolor = black,
    urlcolor = blue,
}   

% For captions for pics
\usepackage{capt-of}


% tcolorbox custom environments
% Definition ========================================================
\newtcolorbox[auto counter,number within=section]{definition}[2][drop shadow]{%
enhanced, before skip=5mm, after skip=6mm, % [V3 ADDITION] spacing
colback=blue!5!white,colframe=blue!50!black,fonttitle=\bfseries,
title= Definition~\thetcbcounter. #2,#1}

% Theorem ==========================================================
\newtcolorbox[auto counter,number within=section]{theorem}[2][drop shadow]{%
enhanced, before skip=5mm, after skip=6mm, % [V3 ADDITION] spacing
colback=black!5!white,colframe=black,fonttitle=\bfseries,
title= Theorem~\thetcbcounter. #2,#1}

% Caution box ======================================================
\newtcolorbox{caution}[1][]{enhanced,
before skip=5mm,after skip=6mm, % [V3 ADDITION] was 2mm/3mm
boxrule=1pt,left=11mm,right=2mm,top=3mm,bottom=3mm,
colback=yellow!50,
colframe=yellow!20!black,
sharp corners,rounded corners=southeast,arc is angular,arc=3mm,
underlay={%
\path[fill=tcbcolback!80!black] ([yshift=3mm]interior.south east)--++(-0.4,-0.1)--++(0.1,-0.2);
\path[draw=tcbcolframe,shorten <=-0.05mm,shorten >=-0.05mm] ([yshift=3mm]interior.south east)--++(-0.4,-0.1)--++(0.1,-0.2);
\path[fill=yellow!50!black,draw=none] (interior.south west) rectangle node[white]{\Huge\bfseries !} ([xshift=9mm]interior.north west);
},
drop fuzzy shadow,#1}

% Question box ======================================================
\newtcolorbox{question}[1][]{enhanced,
before skip=5mm,after skip=6mm, % [V3 ADDITION] was 2mm/3mm
boxrule=1pt,left=11mm,right=2mm,top=3mm,bottom=3mm,
colback=blue!20!white,
colframe=blue!20!black,
sharp corners,rounded corners=southeast,arc is angular,arc=3mm,
underlay={%
\path[fill=tcbcolback!80!black] ([yshift=3mm]interior.south east)--++(-0.4,-0.1)--++(0.1,-0.2);
\path[draw=tcbcolframe,shorten <=-0.05mm,shorten >=-0.05mm] ([yshift=3mm]interior.south east)--++(-0.4,-0.1)--++(0.1,-0.2);
\path[fill=blue!50!black,draw=none] (interior.south west) rectangle node[white]{\Huge\bfseries ?} ([xshift=9mm]interior.north west);
},
drop fuzzy shadow,#1}

% Information box ======================================================
\newtcolorbox{info}[1][]{enhanced,
before skip=5mm,after skip=6mm, % [V3 ADDITION] was 2mm/3mm
boxrule=1pt,left=11mm,right=2mm,top=3mm,bottom=3mm,
colback=cyan!20!white,
colframe=cyan!20!black,
sharp corners,rounded corners=southeast,arc is angular,arc=3mm,
underlay={%
\path[fill=tcbcolback!80!black] ([yshift=3mm]interior.south east)--++(-0.4,-0.1)--++(0.1,-0.2);
\path[draw=tcbcolframe,shorten <=-0.05mm,shorten >=-0.05mm] ([yshift=3mm]interior.south east)--++(-0.4,-0.1)--++(0.1,-0.2);
\path[fill=cyan!50!black,draw=none] (interior.south west) rectangle node[white]{\Huge\bfseries i} ([xshift=9mm]interior.north west);
},
drop fuzzy shadow,#1}

% Mistake box ======================================================
\newtcolorbox{mistake}[1][]{enhanced,
before skip=5mm,after skip=6mm, % [V3 ADDITION] was 2mm/3mm
boxrule=1pt,left=11mm,right=2mm,top=3mm,bottom=3mm,
colback=red!20!white,
colframe=red!20!black,
sharp corners,rounded corners=southeast,arc is angular,arc=3mm,
underlay={%
\path[fill=tcbcolback!80!black] ([yshift=3mm]interior.south east)--++(-0.4,-0.1)--++(0.1,-0.2);
\path[draw=tcbcolframe,shorten <=-0.05mm,shorten >=-0.05mm] ([yshift=3mm]interior.south east)--++(-0.4,-0.1)--++(0.1,-0.2);
\path[fill=red!50!black,draw=none] (interior.south west) rectangle node[white]{\Huge\bfseries $\times$} ([xshift=9mm]interior.north west);
},
drop fuzzy shadow,#1}



% [V3 ADDITION] Worked example box ==================================
% Own counter, numbered within the section (Example 11.2.1, 11.2.2, ...).
% For per-subsection numbering (11.2.1.1) change `section' to `subsection'. Blue `question' boxes are kept for thought-provoking questions
% only; worked examples always use this box. Usage:
%   \begin{example} ...question text... \end{example}
%   \solution{...}
\newtcolorbox[auto counter,number within=section]{example}[1][]{%
enhanced, breakable=false, before skip=5mm, after skip=4mm,
colback=teal!5!white, colframe=teal!55!black, boxrule=0.8pt,
left=3mm, right=3mm, top=2mm, bottom=2mm,
sharp corners, rounded corners=northwest, arc=3mm,
fonttitle=\bfseries, coltitle=white,
title=Example~\thetcbcounter, attach boxed title to top left={yshift=-2mm, xshift=3mm},
boxed title style={colback=teal!55!black, sharp corners, rounded corners=northwest, arc=2mm},
drop fuzzy shadow, #1}

% [V3 ADDITION] Spacing between sections, subsections, parts, paragraphs =
% (fncychap's chapter heading is untouched; only \section and below.)
\makeatletter
\renewcommand\section{\@startsection{section}{1}{\z@}%
  {-2ex \@plus -.5ex \@minus -.2ex}{3.2ex \@plus .3ex}%
  {\normalfont\LARGE\bfseries\color{blue!50!black}}}
\renewcommand\subsection{\@startsection{subsection}{2}{\z@}%
  {-1.5ex \@plus -.5ex \@minus -.2ex}{2.6ex \@plus .3ex}%
  {\normalfont\Large\bfseries}}
\renewcommand\subsubsection{\@startsection{subsubsection}{3}{\z@}%
  {-4ex \@plus -1ex \@minus -.5ex}{1.6ex \@plus .2ex}%
  {\normalfont\large\bfseries\color{teal!50!black}}}
\makeatother
\setlength{\parindent}{0pt}
\setlength{\parskip}{0.9ex plus 0.3ex minus 0.2ex}

% Headers and Footers
\usepackage{fancyhdr,lastpage}
% [V3 R2] grey left-margin line on the out-of-syllabus pages (item 13)
\usepackage{eso-pic}
% [AoI item 63] thicker line (2.6pt, was 1.6pt) that starts just BELOW the header rule
% (top 14mm of the page left clear). Flag: copy this line into every chapter's preamble.
\newcommand{\oosmargin}{\AddToShipoutPictureBG{\put(12mm,0){\color{gray!75}\rule{2.6pt}{\dimexpr\paperheight-14mm\relax}}}}
\pagestyle{fancy}
\lhead{Maximus Pung Jun}
% [BUILD COPY CHANGE 2] Header bug fix: \chaptername is the word "Chapter",
% so the original header printed "Chapter 1: Chapter". \chaptermark now stores
% the bare chapter title in \leftmark, and the header uses \leftmark.
\renewcommand{\chaptermark}[1]{\markboth{#1}{}}
\rhead{Chapter \thechapter: \leftmark}
% \cfoot{Page \thepage\  of \pageref{LastPage}} % buggy, does not appear on chapter pages.

% [V3 R2] tick and cross markers for mistake boxes (G3)
\newcommand{\wrong}{\textcolor{red!70!black}{\ensuremath{\boldsymbol{\times}}}}
\newcommand{\cmark}{\textcolor{green!50!black}{\ensuremath{\checkmark}}}

\newcommand{\xmark}{\wrong}   % alias: red cross
% [V3 R2] out-of-syllabus helpers (DE notes item 13)
%   \oosmargin    : grey line in the left margin on every following page
%   \oosmarginoff : stop the margin line
%   \oosnote      : grey banner "Out of syllabus. For your interest only."
%   outofsyllabus : environment = new page + margin line + banner; ends with
%                   \clearpage and switches the margin line off again
\newcommand{\oosmarginoff}{\ClearShipoutPictureBG}
\newcommand{\oosnote}[1][Nothing in this section is examined in H2 Mathematics.]{\par\noindent\colorbox{gray!15}{\parbox{\dimexpr\textwidth-2\fboxsep}{\textbf{Out of syllabus. For your interest only.} #1}}\par}
\newenvironment{outofsyllabus}{\clearpage\oosmargin}{\clearpage\oosmarginoff}

% Notation as in Maximus's notes: italic d, display style everywhere (G1).
% Fractions inside exponents stay small, so never use these macros inside ^{}.
\newcommand{\dydx}{\dfrac{\mathrm{d}y}{\mathrm{d}x}}
\newcommand{\dvdx}{\dfrac{\mathrm{d}v}{\mathrm{d}x}}
\newcommand{\dudx}{\dfrac{\mathrm{d}u}{\mathrm{d}x}}
\newcommand{\ddydx}{\dfrac{\mathrm{d}^{2}y}{\mathrm{d}x^{2}}}
\newcommand{\ddt}[1]{\dfrac{\mathrm{d}#1}{\mathrm{d}t}}
% Small syllabus flag (no longer used in the DE notes, kept for other chapters)
\newcommand{\syllabusflag}[1]{\par\noindent{\footnotesize\color{red!70!black}\textbf{Syllabus note.} #1}\par}


% Teacher's copy toggle ================================================
% [BUILD COPY CHANGE 3] \ifx compares the whole word (the original
% \if\conform test was always false). DE_notes_student.tex predefines
% \conform as HideSolutions.
% A student copy defines \conform as HideSolutions BEFORE loading this preamble.
\providecommand\conform{ShowSolutions}
\newcommand\ShowSolutionsKey{ShowSolutions}
% [V3 ADDITION] In the student copy the hidden solution leaves stretchable
% blank space (\vfill), so the free space on each page is shared out as
% working space below the examples instead of collecting at the bottom.
% [AoI item 10] "Solution." on its own line, so that (a), (b), (c) line up below it.
\newcommand\solution[1]{\ifx\conform\ShowSolutionsKey
  \par\noindent\textbf{Solution.}\par\nobreak\vspace{0.5ex}\noindent #1\par\medskip\else\par\vspace{\workspace}\fi}
% [V3 R2] minimum working space left under each example in the student copy
\newlength\workspace\setlength\workspace{3cm plus 1fill}


\newcommand{\NN}{\mathbb{N}}
\newcommand{\ZZ}{\mathbb{Z}}
\newcommand{\QQ}{\mathbb{Q}}
\newcommand{\RR}{\mathbb{R}}
\newcommand{\CC}{\mathbb{C}}


% house_local.tex: HOUSE COPY of definite_integrals/common/ws_preamble_local.tex ([HOUSE] figstyles renamed)
% WORKSHEET COPY of ../notes/preamble_local.tex (version the worksheet was last built with,
% from ../_backup_pre_limitsum, 2026-10-09 17:31). Local so the worksheet no longer depends
% on ../notes; changes: \input{house_figstyles}; [2026-10-09 follow-up] \mk marks now bold, as in the notes.
% =====================================================================
% preamble_local.tex: extra packages, styles and macros for the
% Definite Integrals chapter ONLY. Loaded after preamble_build.tex.
% Plot styles and the mini 3D macros are in figstyles.tex (shared with the
% slides and worksheets).
% =====================================================================
\usepgfplotslibrary{fillbetween}
\usepackage{adjustbox} % mini 3D pictures (figstyles.tex)
\usetikzlibrary{arrows.meta, decorations.pathreplacing, calc}
\input{house_figstyles}

% Small fractions are allowed only in exponents and in the limits of an
% integral (which are superscripts/subscripts); everywhere else use \dfrac.
\newcommand{\dx}{\,\mathrm{d}x}
\newcommand{\dy}{\,\mathrm{d}y}
\newcommand{\evalat}[3]{\left[\,#1\,\right]_{#2}^{#3}}

% [item 78] right-aligned A-level style marks at the end of a question part
\newcommand{\mk}[1]{{\unskip\nobreak\hfil\penalty50\hskip1.5em\hbox{}\nobreak\hfill\textbf{[#1]}%
  \parfillskip=0pt\finalhyphendemerits=0\par}}
% [item 11] visually emphasised step labels in worked solutions
\newcommand{\step}[1]{\textcolor{teal!55!black}{\textbf{Step #1}}\quad}
% (a), (b), (c) part lists inside questions and solutions
\newenvironment{parts}{\begin{enumerate}[(a)]\setlength{\itemsep}{0.4ex}}{\end{enumerate}}
% SEAB-style solution layout (approved 10.5.8 template): working on the
% left, region sketch on the right (mini 3D pictures go under the sketch)
\newcommand{\solpair}[3][0.56]{\noindent\begin{minipage}[t]{#1\textwidth}\vspace{0pt}#2\end{minipage}\hfill
  \begin{minipage}[t]{\dimexpr0.98\textwidth-#1\textwidth\relax}\vspace{0pt}\centering #3\end{minipage}\par}
% minus-diagram style (drafts/minus_style_template.tex, items 29/31/62)
\pgfplotsset{minus/.style={axis lines=middle, width=5.0cm, height=4.2cm, ticks=none,
  xlabel={$x$}, ylabel={$y$}, every axis x label/.style={at={(ticklabel* cs:1)},anchor=west},
  every axis y label/.style={at={(ticklabel* cs:1)},anchor=south}, clip=false, axis on top,
  axis line style={gray!70}, every axis label/.append style={gray},
  every node/.append style={text=gray!80!black, font=\footnotesize}}}
\tikzset{hlreg/.style={fill=hl!35, draw=hl, thick}}

\ifx\housexm\housexnotes
  \ifdefined\housenotesLM\else% house_fonts.tex: HOUSE COPY of definite_integrals/common/prelim_fonts.tex (unchanged)
% =====================================================================
% prelim_fonts.tex: fonts of the real 9758 prelim papers (Word + MathType; pdffonts on
% ACJC 2026 P1/P2, ACJC 2025 P1, JPJC 2026, RVHS 2026, CJC/DHS/TJC/NYJC/HCI/EJC 2025):
%   text Times New Roman; cover/footer Arial; RI page footer Calibri (here Carlito);
%   maths = MathType: Times New Roman letters/digits (italic variables; upright digits,
%   function names, d, e, f), SymbolMT for the integral, sigma, pi and = + < >,
%   scripts 7 pt on 12 pt (about 58 %), tight operator spacing.
% Candidate C (teacher-approved 13 Oct 2026): the real Times New Roman and Arial
% (ttf-mscorefonts-installer) + URW Standard Symbols PS (fonts-urw-base35, a SymbolMT
% clone with the same Symbol encoding) for the integral, sigma, pi, = + < >;
% TeX Gyre Termes Math only for what MathType draws itself (radicals, stretchy brackets,
% fraction bars, arrows). House style: upright d (\mathrm{d}x), upright e and f.
% Loaded once by di_style / de_core (HW sheets) and prelim_style (Tests, solutions).
% =====================================================================
\ifdefined\prelimfontsloaded\endinput\fi
\def\prelimfontsloaded{1}
\edef\pfrestoreat{\catcode`\noexpand\@=\the\catcode`\@\relax}\makeatletter
% MathType script sizes: 7 pt superscripts on 12 pt text (58 %), 5 pt second level
\DeclareMathSizes{12}{12}{7}{5}
\DeclareMathSizes{11}{11}{6.5}{5}
\DeclareMathSizes{10.95}{10.95}{6.5}{5}
\DeclareMathSizes{10}{10}{6}{5}
\DeclareMathSizes{9}{9}{5.5}{5}
\DeclareMathSizes{14.4}{14.4}{8.5}{6}
\setmainfont{Times New Roman}
\setsansfont{Arial}
\newfontfamily\coverfont{Arial}
\newfontfamily\calibrifont{Carlito}
\usepackage{unicode-math}
\setmathfont{TeX Gyre Termes Math}
\setmathfont[range={up/{latin,Latin,num}}]{Times New Roman}
\setmathfont[range={it/{latin,Latin}}]{Times New Roman Italic}
\setmathfont[range={bfup/{latin,Latin,num}}]{Times New Roman Bold}
\setmathfont[range={"002B,"003D,"003C,"003E}]{Standard Symbols PS}
\newfontfamily\pf@sym{Standard Symbols PS}
\newfontfamily\pf@symi[FakeSlant=0.21]{Standard Symbols PS}
% MathType-like spacing (tighter than TeX around = + -)
\AtBeginDocument{\thickmuskip=2.5mu plus 1mu\medmuskip=1.5mu plus 1mu\thinmuskip=3mu
% Symbol integral and sigma (18 pt at 12 pt text, as MathType), slanted Symbol pi
 \renewcommand{\int}{\mathop{\vcenter{\hbox{\pf@sym\fontsize{\fpeval{1.5*\f@size}}{0pt}\selectfont\char"F2}}}\nolimits}%
 \renewcommand{\sum}{\mathop{\vcenter{\hbox{\pf@sym\fontsize{\fpeval{1.5*\f@size}}{0pt}\selectfont\char"E5}}}\limits}%
 \renewcommand{\pi}{\mathord{\text{\pf@symi\char"70}}}}
\pfrestoreat
\fi
  % =====================================================================
% house_tf.tex: True/False table (standing rules 2026-11-05/06/07):
%   * a T/F set is a NUMBERED QUESTION (no "Check Your Understanding" heading);
%   * table: row label (a),(b),... | Statement (left) | T/F column (right);
%   * row spacing slightly larger than default (extrarowheight 6pt, arraystretch 1.18,
%     as the Vectors Basics Worksheet after the 2026-11-07 change), font size unchanged;
%   * question doc: T/F column blank; solutions doc (\tfanswerstrue): bold T or F,
%     optional short reason under the statement in small grey.
% Usage:
%   \begin{tftable}  \tfrow{statement}{T}{optional reason}  ...  \end{tftable}
% Plain tabular: a question never splits across pages (house rule), so the table is kept whole.
% =====================================================================
\RequirePackage{array}
\makeatletter
\newif\iftfanswers \tfanswersfalse
\newcounter{tfrow}
\newlength{\tf@lab}\setlength{\tf@lab}{0.9cm}
\newlength{\tf@ans}\setlength{\tf@ans}{1.6cm}
\newenvironment{tftable}{\setcounter{tfrow}{0}\par\smallskip
  \setlength{\extrarowheight}{6pt}\renewcommand{\arraystretch}{1.18}%
  \noindent\begin{tabular}{|>{\centering\arraybackslash}p{\tf@lab}|p{\dimexpr\linewidth-\tf@lab-\tf@ans-6\tabcolsep-4\arrayrulewidth\relax}|>{\centering\arraybackslash}p{\tf@ans}|}
  \hline &\textbf{Statement}&\textbf{T/F}\\\hline}
 {\end{tabular}\par\smallskip}
\newcommand{\tfrow}[3]{\stepcounter{tfrow}\textbf{(\alph{tfrow})} & #1%
  \iftfanswers\if\relax\detokenize{#3}\relax\else\newline{\small\color{black!60}#3}\fi\fi
  & \iftfanswers\textbf{#2}\fi\\\hline}
% explicit row label (a set that keeps the Basics Worksheet letters, e.g. (e) omitted)
\newcommand{\tfrowl}[4]{\textbf{#1} & #2%
  \iftfanswers\if\relax\detokenize{#4}\relax\else\newline{\small\color{black!60}#4}\fi\fi
  & \iftfanswers\textbf{#3}\fi\\\hline}
% \tfcompact: allowed exception (Vectors I HW2/HW3) - compressed rows so the whole set
% stays on page 1 of the question sheet (extrarowheight 1.5pt, tabcolsep 4pt, 0.85/1.45 cm)
\newif\iftf@compact
\newcommand{\tfcompact}{\tf@compacttrue\setlength{\tf@lab}{0.85cm}\setlength{\tf@ans}{1.45cm}}
\renewenvironment{tftable}{\setcounter{tfrow}{0}\par\smallskip
  \iftf@compact\setlength{\extrarowheight}{1.5pt}\setlength{\tabcolsep}{4pt}\else\setlength{\extrarowheight}{6pt}\fi
  \renewcommand{\arraystretch}{1.18}%
  \noindent\begin{tabular}{|>{\centering\arraybackslash}p{\tf@lab}|p{\dimexpr\linewidth-\tf@lab-\tf@ans-6\tabcolsep-4\arrayrulewidth\relax}|>{\centering\arraybackslash}p{\tf@ans}|}
  \hline &\textbf{Statement}&\textbf{T/F}\\\hline}
 {\end{tabular}\par\smallskip}
% tflongtable: same table, may break across pages with the header repeated (solutions
% copies whose reasons make the set longer than a page). Needs longtable.
\RequirePackage{longtable}
\newenvironment{tflongtable}{\setcounter{tfrow}{0}\par\smallskip
  \setlength{\extrarowheight}{6pt}\renewcommand{\arraystretch}{1.18}%
  \setlength{\LTleft}{0pt}\setlength{\LTright}{0pt}%
  \begin{longtable}{|>{\centering\arraybackslash}p{\tf@lab}|p{\dimexpr\linewidth-\tf@lab-\tf@ans-6\tabcolsep-4\arrayrulewidth\relax}|>{\centering\arraybackslash}p{\tf@ans}|}
  \hline &\textbf{Statement}&\textbf{T/F}\\\hline\endhead}
 {\end{longtable}\par}
\makeatother

  \providecommand{\units}[1]{\ \text{units}^{#1}}
  \raggedbottom
\else
  % house_fonts.tex: HOUSE COPY of definite_integrals/common/prelim_fonts.tex (unchanged)
% =====================================================================
% prelim_fonts.tex: fonts of the real 9758 prelim papers (Word + MathType; pdffonts on
% ACJC 2026 P1/P2, ACJC 2025 P1, JPJC 2026, RVHS 2026, CJC/DHS/TJC/NYJC/HCI/EJC 2025):
%   text Times New Roman; cover/footer Arial; RI page footer Calibri (here Carlito);
%   maths = MathType: Times New Roman letters/digits (italic variables; upright digits,
%   function names, d, e, f), SymbolMT for the integral, sigma, pi and = + < >,
%   scripts 7 pt on 12 pt (about 58 %), tight operator spacing.
% Candidate C (teacher-approved 13 Oct 2026): the real Times New Roman and Arial
% (ttf-mscorefonts-installer) + URW Standard Symbols PS (fonts-urw-base35, a SymbolMT
% clone with the same Symbol encoding) for the integral, sigma, pi, = + < >;
% TeX Gyre Termes Math only for what MathType draws itself (radicals, stretchy brackets,
% fraction bars, arrows). House style: upright d (\mathrm{d}x), upright e and f.
% Loaded once by di_style / de_core (HW sheets) and prelim_style (Tests, solutions).
% =====================================================================
\ifdefined\prelimfontsloaded\endinput\fi
\def\prelimfontsloaded{1}
\edef\pfrestoreat{\catcode`\noexpand\@=\the\catcode`\@\relax}\makeatletter
% MathType script sizes: 7 pt superscripts on 12 pt text (58 %), 5 pt second level
\DeclareMathSizes{12}{12}{7}{5}
\DeclareMathSizes{11}{11}{6.5}{5}
\DeclareMathSizes{10.95}{10.95}{6.5}{5}
\DeclareMathSizes{10}{10}{6}{5}
\DeclareMathSizes{9}{9}{5.5}{5}
\DeclareMathSizes{14.4}{14.4}{8.5}{6}
\setmainfont{Times New Roman}
\setsansfont{Arial}
\newfontfamily\coverfont{Arial}
\newfontfamily\calibrifont{Carlito}
\usepackage{unicode-math}
\setmathfont{TeX Gyre Termes Math}
\setmathfont[range={up/{latin,Latin,num}}]{Times New Roman}
\setmathfont[range={it/{latin,Latin}}]{Times New Roman Italic}
\setmathfont[range={bfup/{latin,Latin,num}}]{Times New Roman Bold}
\setmathfont[range={"002B,"003D,"003C,"003E}]{Standard Symbols PS}
\newfontfamily\pf@sym{Standard Symbols PS}
\newfontfamily\pf@symi[FakeSlant=0.21]{Standard Symbols PS}
% MathType-like spacing (tighter than TeX around = + -)
\AtBeginDocument{\thickmuskip=2.5mu plus 1mu\medmuskip=1.5mu plus 1mu\thinmuskip=3mu
% Symbol integral and sigma (18 pt at 12 pt text, as MathType), slanted Symbol pi
 \renewcommand{\int}{\mathop{\vcenter{\hbox{\pf@sym\fontsize{\fpeval{1.5*\f@size}}{0pt}\selectfont\char"F2}}}\nolimits}%
 \renewcommand{\sum}{\mathop{\vcenter{\hbox{\pf@sym\fontsize{\fpeval{1.5*\f@size}}{0pt}\selectfont\char"E5}}}\limits}%
 \renewcommand{\pi}{\mathord{\text{\pf@symi\char"70}}}}
\pfrestoreat

  % house_hw.tex: HOUSE COPY of definite_integrals/common/di_style.tex; [HOUSE] chapter name from \housechapter{...}
% =====================================================================
% di_style.tex: Definite Integrals homework / test layout.
% Loads the house preamble (local copies of the AoI worksheet files:
% ws_preamble_build.tex, ws_preamble_local.tex -> ws_figstyles.tex),
% then adds the homework layout used by HW1-HW3 and the Test.
% Usage (before \begin{document}):  \ditle{Homework 1}
% =====================================================================
% [HOUSE] base, local and fonts are loaded by house.tex
\usepackage{needspace}
\usepackage{zref-savepos,zref-abspage}\makeatletter\zref@addprop{savepos}{abspage}\makeatother
\usepackage{array}

% ---- shortlist toggle ------------------------------------------------
% \ifshortlist true  -> only the questions in SHORTLIST.md (+ challenges)
% A wrapper file sets it with \def\setshortlist{1} before \input-ing the
% main file; the main file can also be edited to say \shortlisttrue.
\newif\ifshortlist
% Task 5: the FINAL homework (= former shortlist) is the default build. HWn.tex -> HWn.pdf.
% A wrapper that sets \def\setfull{1} builds the archive with all original questions.
\ifdefined\setfull\else\shortlisttrue\fi

\makeatletter
% ---- REVIEW build: highlight amended content ---------------------------
% A wrapper sets \def\setreview{1}. In review mode amended text is shaded
% yellow and carries a hover note (PDF tooltip / sticky note); otherwise the
% macros print the content plainly.
%   \chg{note}{inline content}        short inline change (no line breaks), text mode
%   \chgm{note}{math}                 the same inside math mode
%   \begin{chgblock}{note} ... \end{chgblock}   paragraph/display-level change
\newif\ifreview
\ifdefined\setreview \reviewtrue \fi
% hover notes are raw PDF annotations (xdvipdfmx \special; no extra packages).
% Keep notes plain ASCII; balanced parentheses are fine.
\newsavebox{\chg@box}
\newcommand{\chg@ann}[4]{% width, height, depth, note -> invisible markup annotation with a hover note
  \special{pdf:ann width #1 height #2 depth #3 << /Type /Annot /Subtype /Square /Border [0 0 0] /BS << /W 0 >> /C [1 0.85 0] /CA 0.01
    /T (Review) /Subj (Amendment) /Contents (#4) >>}}
\newcommand{\chg@sticky}[1]{\special{pdf:ann width 12pt height 12pt depth 0pt << /Type /Annot /Subtype /Text /Name /Comment
    /C [1 0.85 0] /T (Review) /Subj (Amendment) /Contents (#1) >>}}
\newcommand{\chg}[2]{\ifreview\sbox\chg@box{\colorbox{yellow!60}{#2}}%
  \leavevmode\rlap{\chg@ann{\the\wd\chg@box}{\the\ht\chg@box}{\the\dp\chg@box}{#1}}\usebox\chg@box\else#2\fi}
% \chgm{note}{math}: the same, for use INSIDE math mode (display style kept)
\newcommand{\chgm}[2]{\ifreview\mathchoice{\chg{#1}{$\displaystyle#2$}}{\chg{#1}{$\textstyle#2$}}{\chg{#1}{$\scriptstyle#2$}}{\chg{#1}{$\scriptscriptstyle#2$}}\else#2\fi}
\newtcolorbox{chgbox}[1]{enhanced, breakable, colback=yellow!35, colframe=yellow!35, boxrule=0pt, sharp corners,
  left=2pt, right=14pt, top=1pt, bottom=1pt, before skip=2pt, after skip=2pt,
  overlay first and middle={}, overlay last={}, overlay unbroken and first={\node[anchor=north east, inner sep=1pt] at (frame.north east) {\chg@sticky{#1}\hbox to 12pt{\vrule width 0pt height 12pt}};}}
% \chgnote{note}: just the sticky-note icon (e.g. inside a figure); nothing outside review mode
\newcommand{\chgnote}[1]{\ifreview\chg@sticky{#1}\hbox to 12pt{\vrule width 0pt height 12pt}\fi}
\newenvironment{chgblock}[1]{\ifreview\begin{chgbox}{#1}\fi}{\ifreview\end{chgbox}\fi}

% ---- header ---------------------------------------------------------
\newcommand\di@title{}
\makeatletter
\newcommand{\ditle}[1]{\gdef\di@title{#1}}
\fancypagestyle{di}{\fancyhf{}%
  \lhead{\small \housechap: \di@title}%
  \rhead{\small Compiled by Maximus Pung Jun}%
  \cfoot{\small\thepage}%
  \renewcommand{\headrulewidth}{0.4pt}}
\pagestyle{di}
% title block on page 1, as in his Typst sheets
\newcommand{\dititleblock}[1][]{\thispagestyle{di}%
  {\LARGE\scshape \housechap: \di@title}\par\smallskip
  {Compiled by Maximus Pung Jun}\par\ifx&#1&\else\smallskip{\small\textcolor{black!60}{#1}}\par\fi\medskip}
\makeatother
\newcommand{\disection}[1]{\par\bigskip\noindent{\Large\scshape #1}\par\medskip}

% ---- questions --------------------------------------------------------
% \Qhead{label}{source tag}: question heading; grey [tag] straight after the number, on the left (rule 3B, 9 Nov 2026)
\newcommand{\Qhead}[2]{\par\needspace{8\baselineskip}\bigskip\noindent
  {\large\bfseries #1}\setbox0=\hbox{#2}\ifdim\wd0=0pt\else\enspace{\small\textcolor{black!60}{[#2]}}\fi\par\smallskip}
% part lists with bold labels as he writes them: \pt{(a)}, \pt{(i)}, \pt{(b) (i)}
\newenvironment{qparts}[1][2.4em]{\begin{list}{}{\setlength\leftmargin{#1}%
  \setlength\labelwidth{\dimexpr#1-0.5em\relax}\setlength\labelsep{0.5em}%
  \setlength\itemsep{0.8ex}\setlength\topsep{0.6ex}\setlength\parsep{0.6ex}%
  \renewcommand\makelabel[1]{\textbf{##1}\hfil}}}{\end{list}}
\newcommand{\pt}[1]{\item[#1]}
% thin horizontal label for a challenge question
\newcommand{\challengebanner}{\par\bigskip\noindent\textcolor{black!50}{\rule[0.5ex]{\linewidth}{0.4pt}}\par
  \noindent{\large\scshape Challenge (optional)}\par}
% upright "units"
\newcommand{\units}[1]{\ \text{units}^{#1}}
% END line
\newcommand{\diend}[1][END OF WORKSHEET]{\par\vspace{2\baselineskip}\begin{center}\textbf{#1}\end{center}}

% ---- question registry (full / shortlist from one source) -------------
% \defQ{n}{tag}{body}   question n of the original sheet
% \defA{n}{answer body} its final answer, as printed in his answer section
% \putQ{n} / \putC{n}   print question n as the next (challenge) question
% \putA{n}              print its answer under the same new label
\newcounter{diq}
\makeatletter
\newcommand{\defQ}[3]{\expandafter\gdef\csname diQtag@#1\endcsname{#2}%
  \expandafter\long\expandafter\gdef\csname diQbody@#1\endcsname{#3}}
\newcommand{\defA}[2]{\expandafter\long\expandafter\gdef\csname diA@#1\endcsname{#2}}
\newcommand{\di@tag}[1]{\csname diQtag@#1\endcsname}% source tag only (no original Q number)
% A question never breaks across pages: it is measured and starts on a new page unless it
% fits whole on the rest of this one (only a question taller than a page may break, and then
% only between parts). Page positions of its start/end (zref labels q-s<n>/q-e<n> in the .aux)
% let the checker verify this.
\newsavebox{\di@qmeas}
\newcommand{\di@placeQ}[1]{\par\setbox\di@qmeas=\vbox{#1}%
  \ifdim\dimexpr\ht\di@qmeas+\dp\di@qmeas\relax<0.95\textheight
    \needspace{\dimexpr\ht\di@qmeas+\dp\di@qmeas+\bigskipamount\relax}\else\needspace{14\baselineskip}\fi
  \zsavepos{q-s\thediq}#1\par\zsavepos{q-e\thediq}}
\newcommand{\putQ}[1]{\stepcounter{diq}%
  \expandafter\xdef\csname diLab@#1\endcsname{\thediq}%
  \di@placeQ{\Qhead{\thediq}{\di@tag{#1}}\csname diQbody@#1\endcsname\par}}
\newcommand{\putC}[1]{\stepcounter{diq}%
  \expandafter\xdef\csname diLab@#1\endcsname{\thediq}%
  \expandafter\gdef\csname diChal@#1\endcsname{1}%
  \di@placeQ{\Qhead{\thediq\quad\normalsize\mdseries\scshape Challenge}{\di@tag{#1}}\csname diQbody@#1\endcsname\par}}
\newcommand{\putA}[1]{\par\needspace{4\baselineskip}%
  \begin{qparts}[2.4em]\item[{\csname diLab@#1\endcsname}]\csname diA@#1\endcsname\end{qparts}\medskip}
% ---- worked solutions: \defS{n}{body}, \putQS{n} prints question n + its solution ----
\newcommand{\defS}[2]{\expandafter\long\expandafter\gdef\csname diS@#1\endcsname{#2}}
\newcommand{\disolhead}{\par\medskip\noindent{\bfseries Solution}\par\nobreak\vspace{-0.6ex}\noindent\textcolor{black!35}{\rule{\linewidth}{0.4pt}}\par\nobreak}
\newcommand{\putQS}[1]{\putQ{#1}\disolhead\csname diS@#1\endcsname\par}
% source of a worked solution; visible only in REVIEW builds. #1: his / pub / pip
\newcommand{\solsrc}[2]{\ifreview\par\noindent\def\di@a{#1}\def\di@b{pip}%
  \ifx\di@a\di@b\colorbox{yellow!60}{\parbox{\dimexpr\linewidth-2\fboxsep-14pt}{\small\textbf{Pip-worked:} #2}}\chgnote{Pip-worked solution. No worked solution by the teacher and no published solution was found. Please check it.}%
  \else{\small\textcolor{black!55}{Source of this solution: #2}}\fi\par\smallskip\fi}
\makeatother
% compact sketch for worked solutions: \solfig[axis options]{plot commands}
\newcommand{\solfig}[2][]{\par\begin{center}\begin{tikzpicture}\begin{axis}[qfig, width=6.6cm, height=5cm, #1]#2\end{axis}\end{tikzpicture}\end{center}}
% a sketch added to an existing worked solution (teacher rule), shaded in REVIEW
\newenvironment{addedsketch}{\begin{chgblock}{Sketch added: teacher rule - every area or volume solution has a sketch.}}{\end{chgblock}}
% answer part list (compact)
\newenvironment{aparts}[1][2.4em]{\begin{qparts}[#1]\setlength\itemsep{0.3ex}}{\end{qparts}}

% ---- figure helpers ---------------------------------------------------------
\tikzset{rect/.style={fill=blue!10, draw=blue!55!black, thin}}
\pgfplotsset{qfig/.style={sketch, every node/.append style={font=\scriptsize}}}
% aligned working that starts on the same line as a part label: \pt{(iii)} \begin{topalign} ... \end{topalign}
\newenvironment{topalign}{\hfill$\displaystyle\begin{aligned}[t]}{\end{aligned}$\hfill\null\par\smallskip}
% one display that starts on the same line as a part label
\newcommand{\topdisplay}[1]{\hfill$\displaystyle #1$\hfill\null\par\smallskip}

  % house_solmacros.tex: HOUSE COPY of de-docs/common/de_style.tex minus \input{de_core} and the DE test titles.
% Adds the Ch11 solution helpers: soln/spart tables, sbs/sbsplain side-by-side, \note, \award, \tagm, \dstart.
% =====================================================================
% de_style.tex: Differential Equations (Ch11) homework / test layout.
% = Ch10 di_style (copied as de_core.tex; header text changed) + the DE
% notation and the framed solution boxes of the v1 DE documents
% (soln / spart, used by the Test worked solutions).
% =====================================================================
\providecommand{\e}{\mathrm{e}}
\providecommand{\cosec}{\operatorname{cosec}}
\providecommand{\dt}{\,\mathrm{d}t}
\providecommand{\abs}[1]{\left|#1\right|}
\providecommand{\dd}[2]{\dfrac{\mathrm{d}#1}{\mathrm{d}#2}}
\providecommand{\dvdx}{\dfrac{\mathrm{d}v}{\mathrm{d}x}}
\providecommand{\dudx}{\dfrac{\mathrm{d}u}{\mathrm{d}x}}
\colorlet{notegrey}{black!60}
\makeatletter
\protected\def\note#1{\ifmmode\text{\small\color{notegrey}\itshape #1}\else{\small\color{notegrey}\itshape #1}\fi}
\newcommand{\award}[1]{\unskip\nobreak\hfill\textbf{\small #1}}
\newcommand{\tagm}[1]{\tag*{\textbf{\small #1}}}
\newcommand{\mfAminus}{\[\begin{aligned}&\int\dfrac{1}{a^{2}-x^{2}}\dx\\&\quad=\dfrac{1}{2a}\ln\abs{\dfrac{a+x}{a-x}}+c\end{aligned}\]}
\newcommand{\mfXminus}{\[\begin{aligned}&\int\dfrac{1}{x^{2}-a^{2}}\dx\\&\quad=\dfrac{1}{2a}\ln\abs{\dfrac{x-a}{x+a}}+c\end{aligned}\]}
\newcommand{\sbs@start}{\if@inlabel\def\sbs@shift{\vspace{-\ht\strutbox}}\else
  \ifvmode\else\par\fi\addvspace{0.4ex}\def\sbs@shift{}\fi\noindent}
\newsavebox{\sbs@box}
\newenvironment{sbs}[1]{%
  \sbs@start\begin{minipage}[t]{#1\linewidth}\vspace{0pt}\sbs@shift\raggedright%
  \def\sbsnext{\end{minipage}\hfill\begin{minipage}[t]{\dimexpr0.97\linewidth-#1\linewidth\relax}\vspace{0pt}\sbs@shift\setlength{\fboxsep}{1.5mm}\begin{lrbox}{\sbs@box}\begin{minipage}[t]{\dimexpr\linewidth-2\fboxsep-1pt\relax}\small\raggedright}%
}{\end{minipage}\end{lrbox}\setlength{\fboxrule}{0.5pt}\fbox{\usebox{\sbs@box}}\end{minipage}\par}
\newenvironment{sbsplain}[1]{%
  \sbs@start\begin{minipage}[t]{#1\linewidth}\vspace{0pt}\sbs@shift\raggedright%
  \def\sbsnext{\end{minipage}\hfill\begin{minipage}[t]{\dimexpr0.97\linewidth-#1\linewidth\relax}\vspace{0pt}\sbs@shift\centering}%
}{\end{minipage}\par}
% ---- framed worked-solution box (v1 DE mark-scheme layout, kept for the Test) ----
\newlength{\solcol}\setlength{\solcol}{15mm}
\newcounter{spartno}
\newcommand{\solskips}{\abovedisplayskip=4pt plus 2pt\belowdisplayskip=4pt plus 2pt\abovedisplayshortskip=2pt\belowdisplayshortskip=3pt}
\newtcolorbox{soln}[1][]{enhanced, breakable, sharp corners,
  colback=white, colframe=black, boxrule=0.6pt,
  boxsep=0pt, left=0pt, right=3mm, top=1.5mm, bottom=2mm,
  before skip=1.2ex, after skip=1.5ex,
  before upper={\setcounter{spartno}{0}\parskip=0.5ex\solskips\@fleqntrue\@mathmargin=0pt\raggedright},
  overlay={\draw[line width=0.4pt] ([xshift=\solcol]frame.north west) -- ([xshift=\solcol]frame.south west);},
  #1}
\newenvironment{spart}[1]{%
  \stepcounter{spartno}%
  \ifnum\value{spartno}>1 \par\vspace{0.2ex}\noindent
    \makebox[\linewidth][l]{\rule{\dimexpr\linewidth+3mm\relax}{0.4pt}}\par\nobreak\vspace{0.3ex}\fi
  \begin{list}{}{\leftmargin=\dimexpr\solcol+2.5mm\relax
    \labelwidth=\solcol \labelsep=2.5mm \itemindent=0pt
    \topsep=0.3ex \parsep=0.5ex \rightmargin=0pt \@beginparpenalty=\@M}%
  \item[{\makebox[\solcol][l]{\strut\hspace{1.5mm}\textbf{#1}}}]%
  \ignorespaces}{\end{list}}
\makeatother
% ---- DE sketches: thin light axes (teacher rule), as the Ch10 'sketch' style ----
\pgfplotsset{dsketch/.style={sketch, axis lines=middle, axis line style={thin, black!45, -{Stealth[length=1.6mm]}},
    ticks=none, clip=false, width=6.5cm, height=4.6cm,
    every axis x label/.style={at={(current axis.right of origin)},anchor=west},
    every axis y label/.style={at={(current axis.above origin)},anchor=south},
    every node/.append style={font=\small}},
  dcurve/.style={black, semithick, smooth, samples=120}}
\tikzset{asym/.style={black, dashed, thin}}
\makeatletter
\newcommand{\dstart}{\leavevmode\vadjust{\nobreak\vskip\dimexpr-\baselineskip-\abovedisplayskip+1pt\relax}}
\newcommand{\sdstart}{\leavevmode\vadjust{\nobreak\vskip\dimexpr-\baselineskip-\abovedisplayskip+4pt\relax}}
\makeatother

  % =====================================================================
% house_tf.tex: True/False table (standing rules 2026-11-05/06/07):
%   * a T/F set is a NUMBERED QUESTION (no "Check Your Understanding" heading);
%   * table: row label (a),(b),... | Statement (left) | T/F column (right);
%   * row spacing slightly larger than default (extrarowheight 6pt, arraystretch 1.18,
%     as the Vectors Basics Worksheet after the 2026-11-07 change), font size unchanged;
%   * question doc: T/F column blank; solutions doc (\tfanswerstrue): bold T or F,
%     optional short reason under the statement in small grey.
% Usage:
%   \begin{tftable}  \tfrow{statement}{T}{optional reason}  ...  \end{tftable}
% Plain tabular: a question never splits across pages (house rule), so the table is kept whole.
% =====================================================================
\RequirePackage{array}
\makeatletter
\newif\iftfanswers \tfanswersfalse
\newcounter{tfrow}
\newlength{\tf@lab}\setlength{\tf@lab}{0.9cm}
\newlength{\tf@ans}\setlength{\tf@ans}{1.6cm}
\newenvironment{tftable}{\setcounter{tfrow}{0}\par\smallskip
  \setlength{\extrarowheight}{6pt}\renewcommand{\arraystretch}{1.18}%
  \noindent\begin{tabular}{|>{\centering\arraybackslash}p{\tf@lab}|p{\dimexpr\linewidth-\tf@lab-\tf@ans-6\tabcolsep-4\arrayrulewidth\relax}|>{\centering\arraybackslash}p{\tf@ans}|}
  \hline &\textbf{Statement}&\textbf{T/F}\\\hline}
 {\end{tabular}\par\smallskip}
\newcommand{\tfrow}[3]{\stepcounter{tfrow}\textbf{(\alph{tfrow})} & #1%
  \iftfanswers\if\relax\detokenize{#3}\relax\else\newline{\small\color{black!60}#3}\fi\fi
  & \iftfanswers\textbf{#2}\fi\\\hline}
% explicit row label (a set that keeps the Basics Worksheet letters, e.g. (e) omitted)
\newcommand{\tfrowl}[4]{\textbf{#1} & #2%
  \iftfanswers\if\relax\detokenize{#4}\relax\else\newline{\small\color{black!60}#4}\fi\fi
  & \iftfanswers\textbf{#3}\fi\\\hline}
% \tfcompact: allowed exception (Vectors I HW2/HW3) - compressed rows so the whole set
% stays on page 1 of the question sheet (extrarowheight 1.5pt, tabcolsep 4pt, 0.85/1.45 cm)
\newif\iftf@compact
\newcommand{\tfcompact}{\tf@compacttrue\setlength{\tf@lab}{0.85cm}\setlength{\tf@ans}{1.45cm}}
\renewenvironment{tftable}{\setcounter{tfrow}{0}\par\smallskip
  \iftf@compact\setlength{\extrarowheight}{1.5pt}\setlength{\tabcolsep}{4pt}\else\setlength{\extrarowheight}{6pt}\fi
  \renewcommand{\arraystretch}{1.18}%
  \noindent\begin{tabular}{|>{\centering\arraybackslash}p{\tf@lab}|p{\dimexpr\linewidth-\tf@lab-\tf@ans-6\tabcolsep-4\arrayrulewidth\relax}|>{\centering\arraybackslash}p{\tf@ans}|}
  \hline &\textbf{Statement}&\textbf{T/F}\\\hline}
 {\end{tabular}\par\smallskip}
% tflongtable: same table, may break across pages with the header repeated (solutions
% copies whose reasons make the set longer than a page). Needs longtable.
\RequirePackage{longtable}
\newenvironment{tflongtable}{\setcounter{tfrow}{0}\par\smallskip
  \setlength{\extrarowheight}{6pt}\renewcommand{\arraystretch}{1.18}%
  \setlength{\LTleft}{0pt}\setlength{\LTright}{0pt}%
  \begin{longtable}{|>{\centering\arraybackslash}p{\tf@lab}|p{\dimexpr\linewidth-\tf@lab-\tf@ans-6\tabcolsep-4\arrayrulewidth\relax}|>{\centering\arraybackslash}p{\tf@ans}|}
  \hline &\textbf{Statement}&\textbf{T/F}\\\hline\endhead}
 {\end{longtable}\par}
\makeatother

  \ifx\housexm\housexsol \def\prelimkeepcolour{1}% house_prelim.tex: HOUSE COPY of definite_integrals/common/prelim_style.tex; [HOUSE] titles from \housechapter
% =====================================================================
% prelim_style.tex: makes the Test look like a Singapore JC H2 Mathematics
% 9758 prelim paper (studied from ACJC/RI 2026 P1/P2 on the box:
% body Times New Roman + Times-style maths, cover in Arial).
% Fonts: prelim_fonts.tex (real Times New Roman/Arial + Times-style maths).
% Shared by Ch10 (di_style) and Ch11 (de_style; sets \prelimtitle/\prelimsubject). Loaded AFTER the chapter style (overrides its layout for the test only).
% =====================================================================
% [HOUSE] fonts loaded by house.tex
\providecommand{\prelimtitle}{\MakeUppercase{\housechap} TEST}
\providecommand{\prelimsubject}{\housechap}
\geometry{a4paper, top=22mm, bottom=24mm, left=25mm, right=25mm, headsep=6mm, footskip=12mm}
\setlength{\parindent}{0pt}
\setlength{\parskip}{0.9ex plus 0.3ex minus 0.2ex}
\linespread{1.3}
\setlength{\emergencystretch}{3em}
\clubpenalty=10000 \widowpenalty=10000 \interlinepenalty=10000
\AtBeginDocument{\fontsize{12}{14.4}\selectfont}
\renewcommand{\normalsize}{\fontsize{12}{14.4}\selectfont}

% ---- figures: plain black/grey like a printed paper ----------------------
\ifdefined\prelimkeepcolour\else % question paper: black/grey figures; solutions keep colour
\pgfplotsset{curve/.style={black, semithick, smooth, samples=150},
  line2/.style={black, semithick, samples=2},
  sketch/.append style={axis line style={thin, black}, every axis label/.append style={black}}}
\tikzset{region/.style={fill=black!15, draw=none}, regionB/.style={fill=black!30, draw=none},
  rect/.style={fill=black!8, draw=black, thin}, ax3/.style={-{Stealth[length=1.5mm]}, thin, black}}
\colorlet{blue}{black}\colorlet{red}{black}\colorlet{orange}{black}
\fi

% ---- header / footer: page number centred top (bold), [Turn over at bottom right
\makeatletter
\providecommand{\prelimfootlabel}{\prelimtitle\ 2026}
\newcommand{\pp@footer}{{\coverfont\scriptsize \prelimfootlabel\hspace{3em}H2 MATHEMATICS 9758}}
\fancypagestyle{prelim}{\fancyhf{}%
  \chead{\bfseries\thepage}%
  \lfoot{\pp@footer}%
  \rfoot{\ifnum\value{page}<\getpagerefnumber{LastPage}\relax\textbf{[Turn over}\fi}%
  \renewcommand{\headrulewidth}{0pt}\renewcommand{\footrulewidth}{0pt}}
\fancypagestyle{prelimcover}{\fancyhf{}%
  \rfoot{\textbf{[Turn over}}%
  \renewcommand{\headrulewidth}{0pt}\renewcommand{\footrulewidth}{0pt}}
\pagestyle{prelim}

% ---- marks: [n] regular weight, right-aligned at the end of the part ----------
\renewcommand{\mk}[1]{{\unskip\nobreak\hfil\penalty50\hskip1.5em\hbox{}\nobreak\hfill[#1]%
  \parfillskip=0pt\finalhyphendemerits=0\par}}

% ---- questions: bold number at the left margin, text indented; grey [tag] straight after the number (3B)
\newsavebox{\pp@qmeas}
\newcommand{\pp@qlist}[1]{%
  \begin{list}{}{\setlength\leftmargin{2.6em}\setlength\labelwidth{2.1em}\setlength\labelsep{0.5em}%
     \setlength\topsep{0pt}\setlength\itemsep{0pt}\setlength\parsep{0.9ex}%
     \renewcommand\makelabel[1]{\textbf{##1}\hfil}}%
  \item[\thediq]\@endparpenalty\@M\setbox0=\hbox{\di@tag{#1}}\ifdim\wd0=0pt\else{\small\textcolor{black!60}{[\di@tag{#1}]}}\par\nobreak\fi\csname diQbody@#1\endcsname\par\penalty\@M
  \end{list}}
% a question starts on a new page unless it fits whole on the current one
% (questions taller than a page may break, but only between parts)
\renewcommand{\putQ}[1]{\stepcounter{diq}%
  \expandafter\xdef\csname diLab@#1\endcsname{\thediq}%
  \par\setbox\pp@qmeas=\vbox{\pp@qlist{#1}}%
  \ifdim\dimexpr\ht\pp@qmeas+\dp\pp@qmeas\relax<0.95\textheight
    \needspace{\dimexpr\ht\pp@qmeas+\dp\pp@qmeas\relax}\else\needspace{14\baselineskip}\fi
  \zsavepos{q-s\thediq}\pp@qlist{#1}\par\zsavepos{q-e\thediq}\vspace{5ex plus 2ex minus 1ex}}
% part labels (a)/(i) bold, indented under the question text
\renewenvironment{qparts}[1][2.6em]{\begin{list}{}{\setlength\leftmargin{#1}%
  \setlength\labelwidth{\dimexpr#1-0.5em\relax}\setlength\labelsep{0.5em}%
  \setlength\itemsep{1.4ex}\setlength\topsep{1.2ex}\setlength\parsep{0.8ex}%
  \renewcommand\makelabel[1]{\textbf{##1}\hfil}}}{\end{list}}
\makeatother

% ---- cover page --------------------------------------------------------------
\newcommand{\prelimcover}[3]{% #1 total marks, #2 duration, #3 date line
\thispagestyle{prelimcover}%
{\coverfont\fontsize{10}{12}\selectfont\linespread{1}\selectfont
\noindent\begin{minipage}[t]{\linewidth}
  {\bfseries \prelimtitle}\\[0.2ex]
  {\bfseries Compiled by Maximus Pung Jun}\\[1.6ex]
  Higher 2
\end{minipage}\par\vspace{2ex}
\noindent\rule{\linewidth}{1.2pt}\par\vspace{0.4ex}
\noindent{\Large\bfseries MATHEMATICS}\hfill{\Large\bfseries 9758}\par\vspace{1.2ex}
\noindent \prelimsubject\hfill\textbf{#3}\\
QUESTION PAPER\hfill\textbf{#2}\par\vspace{2.5ex}
\noindent\begin{tabular}{@{}p{46mm}l@{}}Additional Materials: & Writing paper\\ & List of Formulae and Results (MF27)\end{tabular}\par\vspace{2ex}
\noindent\rule{\linewidth}{0.6pt}\par\vspace{0.6ex}
{\bfseries READ THESE INSTRUCTIONS FIRST}\par\vspace{2ex}
\setlength{\parskip}{0pt}
Write your name and class on all the work you hand in.\\
Write in dark blue or black pen.\\
You may use an HB pencil for any diagrams or graphs.\\
Do not use staples, paper clips, glue or correction fluid.\par\vspace{2ex}
Answer \textbf{all} the questions.\\
Give non-exact numerical answers correct to 3 significant figures, or 1 decimal place in the case of angles in degrees, unless a different level of accuracy is specified in the question.\\
You are expected to use an approved graphing calculator.\\
Unsupported answers from a graphing calculator are allowed unless a question specifically states otherwise.\\
Where unsupported answers from a graphing calculator are \textbf{not} allowed in a question, you are required to present the mathematical steps using mathematical notations and not calculator commands.\\
You are reminded of the need for clear presentation in your answers.\par\vspace{2ex}
The number of marks is given in brackets [\ ] at the end of each question or part question.\\
The total number of marks for this paper is #1.
\vfill
\noindent\rule{\linewidth}{0.6pt}\par\vspace{0.4ex}
\centerline{This document consists of \textbf{\pageref{LastPage}} printed pages.}\par\vspace{1ex}
}\clearpage}

% ---- solutions header block (page 1 of the worked solutions; same cover style) ----
\newcommand{\prelimsolhead}[1]{% #1 total marks
{\coverfont\fontsize{10}{12}\selectfont\linespread{1}\selectfont
\noindent{\bfseries \prelimtitle}\\[0.2ex]{\bfseries Compiled by Maximus Pung Jun}\\[1.6ex]Higher 2\par\vspace{2ex}
\noindent\rule{\linewidth}{1.2pt}\par\vspace{0.4ex}
\noindent{\Large\bfseries MATHEMATICS}\hfill{\Large\bfseries 9758}\par\vspace{1.2ex}
\noindent \prelimsubject\hfill\textbf{2026}\\
WORKED SOLUTIONS\hfill\textbf{Total marks: #1}\par\vspace{1ex}
\noindent\rule{\linewidth}{0.6pt}\par}\vspace{4ex}}
% house_sol.tex: HOUSE COPY of definite_integrals/common/ri_solstyle.tex (unchanged)
% =====================================================================
% ri_solstyle.tex: worked solutions in the layout of the RI 2026 Year 6 Prelim
% "Solutions with comments" (grail.moe 26612/26613), loaded AFTER prelim_style.
%   * page 1: bold school-style block + bold underlined title; later pages: running
%     title top-left with a rule under it; footer "Page N of M" bottom right (Calibri).
%   * each question: prelim-format question text (bold number, marks right, source tag
%     bottom right), then a bordered table:  Solution | Comments  (bold header row);
%     body rows  label + [marks] (bold, narrow left column) | working | comments,
%     a full-width rule between parts, working left-aligned (fleqn).
% Rows come from the existing sources: \pt in qparts (Ch10 HW), \slabel (Ch10 Test),
% spart (Ch11). Marks under each label are taken from the question's \mk.
% =====================================================================
\makeatletter
\usepackage{lastpage}
\usepackage{adjustbox}
\usepackage{zref-savepos,zref-abspage}
\zref@addprop{savepos}{abspage}
\geometry{a4paper, top=22mm, bottom=24mm, left=22mm, right=22mm, headsep=8mm, footskip=12mm, headheight=24pt}
\providecommand{\risoltitle}{\prelimtitle: Solutions with comments}
\providecommand{\risolschool}{Compiled by Maximus Pung Jun}
\fancypagestyle{risol}{\fancyhf{}%
  \lhead{\fontsize{11}{12}\selectfont 2026 \risoltitle\\[-1.2ex]\rule{\textwidth}{0.6pt}}%
  \rfoot{\calibrifont\fontsize{11}{13}\selectfont Page \textbf{\thepage} of \textbf{\pageref*{LastPage}}}%
  \renewcommand{\headrulewidth}{0pt}\renewcommand{\footrulewidth}{0pt}}
\fancypagestyle{risolfirst}{\fancyhf{}%
  \rfoot{\calibrifont\fontsize{11}{13}\selectfont Page \textbf{\thepage} of \textbf{\pageref*{LastPage}}}%
  \renewcommand{\headrulewidth}{0pt}\renewcommand{\footrulewidth}{0pt}}
\pagestyle{risol}
\newif\ifri@fresh
\newcommand{\risolfront}{\thispagestyle{risolfirst}\global\ri@freshtrue%
  {\noindent\bfseries\linespread{1}\selectfont \risolschool\\ H2 Mathematics (9758)\\ 2026\par}%
  \vspace{4ex}{\noindent\bfseries\underline{2026 \risoltitle}\par}\vspace{3ex}}
% ---- table geometry (RI: label 42pt | working 284pt | comments 142pt of 468pt) ----
\newlength{\ri@L}\setlength{\ri@L}{15mm}
\newlength{\ri@C}\setlength{\ri@C}{40mm}
\newlength{\ri@pad}\setlength{\ri@pad}{1.8mm}
\newif\ifri@sol \newif\ifri@first
\newcount\ri@depth
\def\ri@outer{}
\newtcolorbox{risol}{enhanced, breakable, sharp corners, frame hidden, colback=white,
  boxrule=0pt, boxsep=0pt, left=\dimexpr\ri@L+\ri@pad\relax, right=\dimexpr\ri@C+\ri@pad\relax,
  top=1.2mm, bottom=2.5mm, toptitle=0.6mm, bottomtitle=0.6mm, lefttitle=\ri@pad, righttitle=0pt,
  colbacktitle=white, coltitle=black, fonttitle=\bfseries,
  title={Solution\hfill\makebox[\ri@C][l]{\hspace{\ri@pad}Comments}},
  before skip=2ex, after skip=2ex,
  overlay unbroken and first={\draw[line width=0.5pt] (frame.south west) rectangle (frame.north east);
     \draw[line width=0.5pt] (title.south west) -- (title.south east);
     \draw[line width=0.5pt] ([xshift=\ri@L]interior.north west) -- ([xshift=\ri@L]frame.south west);
     \draw[line width=0.5pt] ([xshift=-\ri@C]frame.north east) -- ([xshift=-\ri@C]frame.south east);},
  overlay middle and last={\draw[line width=0.5pt] (frame.south west) rectangle (frame.north east);
     \draw[line width=0.5pt] ([xshift=\ri@L]frame.north west) -- ([xshift=\ri@L]frame.south west);
     \draw[line width=0.5pt] ([xshift=-\ri@C]frame.north east) -- ([xshift=-\ri@C]frame.south east);},
  after upper={\ri@endrow\protected@write\@auxout{}{\string\ri@settotal{\ri@qid}{\the\ri@sumH}}},
  before upper={\global\ri@sumH=14mm\ri@soltrue\ri@firsttrue\ri@depth=0 \parindent=0pt \parskip=0.6ex
     \predisplaypenalty=100 
     \linespread{1.12}\selectfont\raggedright\@fleqntrue\@mathmargin=0pt
     \abovedisplayskip=5pt plus 2pt\belowdisplayskip=5pt plus 2pt
     \abovedisplayshortskip=3pt plus 1pt\belowdisplayshortskip=4pt plus 2pt
     \ifcsname ri@hp@\ri@qid\endcsname\else\ri@row{}\fi}}
\def\ri@qid{}
% one row: full-width rule (glue leaders, so it vanishes at a page break where the frame
% edge takes its place; not before the first row), then bold label + bold [marks]
% in the left column, on the baseline of the first line of the working
\ExplSyntaxOn
\tl_new:N \g__ri_tmp_tl
\cs_new_protected:Npn \ri_setkey:Nn #1#2
  { \group_begin: \cs_set_eq:NN \chg \use_ii:nn \cs_set_eq:NN \chgm \use_ii:nn
    \cs_set_eq:NN \textbf \use:n \cs_set_eq:NN \mathrm \use:n
    \tl_gset:Nx \g__ri_tmp_tl {#2} \group_end:
    \tl_set_eq:NN #1 \g__ri_tmp_tl \tl_remove_all:Nn #1 {~} }
\cs_set_eq:NN \ri@setkey \ri_setkey:Nn
\ExplSyntaxOff
% ---- page-break rules (positions from the previous run, zref labels in the .aux) ----
% * a row (one part) never breaks across pages: if it does not fit on the rest of the
%   page it starts the next page (forced break between rows); only a row taller than a
%   whole page may break inside;
% * a whole solution that fits on one page is kept together (needspace before the table);
%   a longer one starts where at least its first row fits;
% * a row is at least as tall as its label + [marks].
\newcount\ri@n \newbox\ri@labbox \def\ri@open{0}
\newdimen\ri@ybot \newdimen\ri@ytop \newdimen\ri@h \newdimen\ri@sumH
\AtBeginDocument{\ri@ybot=\dimexpr\paperheight-1in-\voffset-\topmargin-\headheight-\headsep-\textheight\relax
  \ri@ytop=\dimexpr\ri@ybot+\textheight\relax}
\def\ri@pg#1{\zref@extractdefault{#1}{abspage}{-1}}
% \ri@rowh{n}: height of row n in the previous run -> \ri@h (0pt if unknown)
\newcommand{\ri@rowh}[1]{\ri@h=0pt
  \ifnum\ri@pg{ri-s#1}>0 \ifnum\ri@pg{ri-e#1}>0
    \ifnum\ri@pg{ri-s#1}=\ri@pg{ri-e#1}\relax \ri@h=\dimexpr\zposy{ri-s#1}sp-\zposy{ri-e#1}sp\relax
    \else \ri@h=\dimexpr\zposy{ri-s#1}sp-\ri@ybot+\ri@ytop-\zposy{ri-e#1}sp-6mm\relax\fi\fi\fi}
\newcommand{\ri@setfirst}[2]{\expandafter\gdef\csname ri@first@#1\endcsname{#2}}
\newcommand{\ri@setkept}[1]{\expandafter\gdef\csname ri@kept@#1\endcsname{}}
\newcommand{\ri@settotal}[2]{\expandafter\gdef\csname ri@total@#1\endcsname{#2}}
\newcommand{\ri@endrow}{\ifnum\ri@open=1 \par\zsavepos{ri-e\the\ri@n}%
  \ifnum\ri@pg{ri-s\the\ri@n}=\ri@pg{ri-e\the\ri@n}\relax
    \dimen@=\dimexpr\ri@curH sp+1.6mm-\zposy{ri-s\the\ri@n}sp+\zposy{ri-e\the\ri@n}sp\relax
    \ifdim\dimen@>0pt \nobreak\kern\dimen@\penalty\z@\fi\fi
  \ri@rowh{\the\ri@n}\global\advance\ri@sumH\dimexpr\ri@h+3.5mm\relax
  \gdef\ri@open{0}\fi}
\newcommand{\ri@row}[1]{\par\ri@endrow\global\advance\ri@n\@ne\gdef\ri@open{1}\@beginparpenalty\@M
  \ifri@first\global\ri@firstfalse\ri@firstfalse
    \protected@write\@auxout{}{\string\ri@setfirst{\ri@qid}{\the\ri@n}}%
  \else
    \zsavepos{ri-p\the\ri@n}%
    % push the row to the next page if it would not fit on the rest of this one
    \ri@rowh{\the\ri@n}%
    \ifnum\ri@pg{ri-p\the\ri@n}>0 \ifdim\ri@h>0pt
      \ifdim\dimexpr\ri@h+4mm\relax<\dimexpr\textheight-10mm\relax
        \ifdim\dimexpr\ri@h+7mm\relax>\dimexpr\zposy{ri-p\the\ri@n}sp-\ri@ybot\relax
          \penalty-\@M\fi\fi\fi\fi
    \vskip0.6ex\nointerlineskip
    \leaders\hbox{\hskip-\dimexpr\ri@L+\ri@pad\relax\vrule width\dimexpr\linewidth+\ri@L+\ri@C+2\ri@pad\relax height0.5pt depth0pt}\vskip0.5pt
    \nobreak\vskip1.4mm\fi
  \ri@setkey\ri@key{\ri@qid#1}\def\ri@lab{#1}%
  \setbox\ri@labbox=\vtop{\hsize=\dimexpr\ri@L-1.6mm\relax\bfseries\raggedright\parskip=0pt
      \strut\ri@lab\ifcsname ri@mk@\ri@key\endcsname\ifx\ri@lab\@empty\else\\\fi
      [\csname ri@mk@\ri@key\endcsname]\fi\strut\par}%
  \xdef\ri@curH{\number\dimexpr\ht\ri@labbox+\dp\ri@labbox\relax}%
  \nointerlineskip\zsavepos{ri-s\the\ri@n}\nointerlineskip\nobreak
  \vbox to 0pt{\hbox{\hskip-\dimexpr\ri@L+\ri@pad-1.2mm\relax\box\ri@labbox}\vss}\nobreak\prevdepth=\dp\strutbox
  \vskip-\parskip\ignorespaces}
% ---- qparts / \pt: question mode records labels (for the marks); solution mode makes rows
\NewCommandCopy\ri@qparts@begin\qparts \let\ri@qparts@end\endqparts
\renewenvironment{qparts}[1][2.6em]{\advance\ri@depth\@ne
  \ifri@sol\else\ri@qparts@begin[#1]\fi}{\ifri@sol\par\else\ri@qparts@end\fi}
\def\ri@qpartsname{qparts}
\renewcommand{\pt}[1]{\ifnum\ri@depth<2 \gdef\ri@outer{#1}\gdef\ri@cur{#1}\else\expandafter\gdef\expandafter\ri@cur\expandafter{\ri@outer#1}\fi
  \ifri@sol\expandafter\ri@solpt\else\expandafter\ri@qpt\fi{#1}}
\def\ri@qpt#1{\expandafter\gdef\csname ri@hp@\ri@qid\endcsname{}\item[#1]}
\def\ri@solpt#1{\@ifnextchar\begin\ri@ptpeek{\ri@row{\ri@cur}}}
\def\ri@ptpeek\begin#1{\def\ri@tmp{#1}\ifx\ri@tmp\ri@qpartsname
   \expandafter\@firstoftwo\else\expandafter\@secondoftwo\fi
  {\begin{qparts}}{\ri@row{\ri@cur}\begin{#1}}}
% marks: record [n] against the current question/part
\let\ri@mk\mk
\renewcommand{\mk}[1]{\ri@setkey\ri@k{\ri@qid\ri@cur}\ri@setkey\ri@mv{#1}\expandafter\xdef\csname ri@mk@\ri@k\endcsname{\ri@mv}\ri@mk{#1}}
% \slabel (Ch10 Test answers): letter labels set the outer part, roman ones are appended;
% consecutive \slabel calls merge into one row label
\def\ri@romans{(i)(ii)(iii)(iv)(v)(vi)(vii)}
\def\ri@firstpart(#1)#2\@nil{(#1)}
\def\slabel#1{\ifri@sol\expandafter\ri@solslabel\else\expandafter\ri@qslabel\fi{#1}}
\def\ri@qslabel#1{\par\smallskip\noindent\textbf{#1}\par\nobreak}
\def\ri@solslabel#1{\@expandtwoargs\in@{#1}{\ri@romans}%
  \ifin@\expandafter\gdef\expandafter\ri@cur\expandafter{\ri@outer#1}\else\gdef\ri@cur{#1}\xdef\ri@outer{\ri@firstpart#1\@nil}\fi
  \@ifnextchar\slabel{}{\ri@row{\ri@cur}}}
% Ch11 spart{label}: a row; soln: the table
\def\spart#1{\gdef\ri@cur{#1}\ri@row{#1}}\def\endspart{\par}
\def\soln{\@ifnextchar[{\ri@solnopt}{\risol}}\def\ri@solnopt[#1]{\risol}\def\endsoln{\endrisol}
% ---- working that started on the label line is left-aligned in the working column
\AtBeginDocument{\@ifundefined{topalign}{}{%
  \renewenvironment{topalign}{\par\noindent$\displaystyle\begin{aligned}[t]}{\end{aligned}$\par\smallskip}%
  \renewcommand{\topdisplay}[1]{\par\noindent$\displaystyle #1$\par\smallskip}}}
% ---- Ch11 side-by-side (text | sketch): the sketch shrinks to fit its half of the working column
\AtBeginDocument{\@ifundefined{sbsplain}{}{%
  \renewenvironment{sbsplain}[1]{%
    \sbs@start\begin{minipage}[t]{#1\linewidth}\vspace{0pt}\sbs@shift\raggedright%
    \def\sbsnext{\end{minipage}\hfill\begin{minipage}[t]{\dimexpr0.97\linewidth-#1\linewidth\relax}\vspace{0pt}\sbs@shift\centering
      \begin{adjustbox}{max width=\linewidth}}%
  }{\end{adjustbox}\end{minipage}\par}}}
% ---- figures in the working column shrink to fit it (only if wider), colours kept
\newcommand{\rifitfig}[1]{\expandafter\let\csname ri@orig@\string#1\endcsname#1%
  \def#1{\@ifnextchar\bgroup{\ri@fitarg#1}{\ri@fitnoarg#1}}}
\def\ri@fitarg#1#2{\ifri@sol\adjustbox{max width=\linewidth}{\csname ri@orig@\string#1\endcsname{#2}}\else\csname ri@orig@\string#1\endcsname{#2}\fi}
\def\ri@fitnoarg#1{\ifri@sol\adjustbox{max width=\linewidth}{\csname ri@orig@\string#1\endcsname}\else\csname ri@orig@\string#1\endcsname\fi}
\@ifundefined{solfig}{}{\renewcommand{\solfig}[2][]{\par\begin{center}\adjustbox{max width=\linewidth}{\begin{tikzpicture}\begin{axis}[qfig, width=6.6cm, height=5cm, #1]#2\end{axis}\end{tikzpicture}}\end{center}}}
% \ricomment{text}: text in the Comments column of the current row (call right after \pt,
% \slabel or \spart). The row is made at least as tall as the comment.
\newbox\ri@cmtbox
\newcommand{\ricomment}[1]{\par\setbox\ri@cmtbox=\vtop{\hsize=\dimexpr\ri@C-2\ri@pad\relax
    \parskip=0pt\raggedright\small{\color{black!60}\strut#1\strut\par}}%
  \ifdim\dimexpr\ht\ri@cmtbox+\dp\ri@cmtbox\relax>\ri@curH sp\xdef\ri@curH{\number\dimexpr\ht\ri@cmtbox+\dp\ri@cmtbox\relax}\fi
  \nointerlineskip\vbox to 0pt{\hbox to\linewidth{\hskip\dimexpr\linewidth+2\ri@pad\relax\box\ri@cmtbox\hss}\vss}\nobreak\prevdepth=\dp\strutbox
  \vskip-\parskip\ignorespaces}
% \putTFlong{n}: the True/False question n with its answer table allowed to break across
% pages (the body of \defQ{n} must be: intro text, then \begin{tftable}...\end{tftable});
% the solutions copy prints it with tflongtable.
\renewcommand{\putTFlong}[1]{\ri@newq\begingroup\let\tftable\tflongtable\let\endtftable\endtflongtable
  \stepcounter{diq}\expandafter\xdef\csname diLab@#1\endcsname{\thediq}%
  \par\noindent\hbox to 2.6em{\textbf{\thediq}\hfil}\setbox0=\hbox{\di@tag{#1}}\ifdim\wd0=0pt\else{\small\textcolor{black!60}{[\di@tag{#1}]}}\par\noindent\hskip2.6em\fi\ignorespaces\csname diQbody@#1\endcsname\par
  \endgroup\vspace{3ex}}
% ---- question + solution; each question starts on a new page (as RI)
\def\ri@qid{}
\let\ri@putQ\putQ
\renewcommand{\putQ}[1]{\gdef\ri@qid{#1}\global\ri@depth=0 \gdef\ri@cur{}\ri@putQ{#1}}
% \risolQ{n}{body macro name}: question n + its solution body in the table;
% \risolQraw: body already contains its own soln environment (Ch11 Test)
\newcommand{\ri@newq}{\ifri@fresh\global\ri@freshfalse\else\clearpage\fi}
% keep the whole table together if it fits on a page; otherwise make sure at least its
% header and first row fit before starting it
\newcommand{\ri@keep}[1]{\par
  \ifcsname ri@total@#1\endcsname
    % hysteresis: a table kept whole in the previous run stays whole up to textheight-4.5mm,
    % so a total near the limit cannot flip between runs
    \ifdim\csname ri@total@#1\endcsname<\dimexpr\textheight-\ifcsname ri@kept@#1\endcsname4.5\else8\fi mm\relax
      \protected@write\@auxout{}{\string\ri@setkept{#1}}%
      \needspace{\dimexpr\csname ri@total@#1\endcsname+4mm\relax}%
    \else\ifcsname ri@first@#1\endcsname
      \ri@rowh{\csname ri@first@#1\endcsname}%
      \ifdim\ri@h<\dimexpr\textheight-20mm\relax\needspace{\dimexpr\ri@h+18mm\relax}\fi\fi\fi\fi}
\newcommand{\risolQ}[2]{\ri@newq\putQ{#1}\vspace{-3ex}\ri@keep{#1}\begin{risol}\csname #2@#1\endcsname\par\end{risol}}
\newcommand{\risolQraw}[2]{\ri@newq\putQ{#1}\vspace{-3ex}\ri@keep{#1}\csname #2@#1\endcsname\par}
\renewcommand{\putQS}[1]{\risolQ{#1}{diS}}
% end marker only if it fits on the last page (no page holding just the marker)
\renewcommand{\diend}[1][END OF SOLUTIONS]{\par\ifdim\dimexpr\pagegoal-\pagetotal\relax<4\baselineskip\else\vspace{1.5ex}\begin{center}\textbf{#1}\end{center}\fi}
\makeatother
\tfanswerstrue\fi
  \ifx\housexm\housextest % house_prelim.tex: HOUSE COPY of definite_integrals/common/prelim_style.tex; [HOUSE] titles from \housechapter
% =====================================================================
% prelim_style.tex: makes the Test look like a Singapore JC H2 Mathematics
% 9758 prelim paper (studied from ACJC/RI 2026 P1/P2 on the box:
% body Times New Roman + Times-style maths, cover in Arial).
% Fonts: prelim_fonts.tex (real Times New Roman/Arial + Times-style maths).
% Shared by Ch10 (di_style) and Ch11 (de_style; sets \prelimtitle/\prelimsubject). Loaded AFTER the chapter style (overrides its layout for the test only).
% =====================================================================
% [HOUSE] fonts loaded by house.tex
\providecommand{\prelimtitle}{\MakeUppercase{\housechap} TEST}
\providecommand{\prelimsubject}{\housechap}
\geometry{a4paper, top=22mm, bottom=24mm, left=25mm, right=25mm, headsep=6mm, footskip=12mm}
\setlength{\parindent}{0pt}
\setlength{\parskip}{0.9ex plus 0.3ex minus 0.2ex}
\linespread{1.3}
\setlength{\emergencystretch}{3em}
\clubpenalty=10000 \widowpenalty=10000 \interlinepenalty=10000
\AtBeginDocument{\fontsize{12}{14.4}\selectfont}
\renewcommand{\normalsize}{\fontsize{12}{14.4}\selectfont}

% ---- figures: plain black/grey like a printed paper ----------------------
\ifdefined\prelimkeepcolour\else % question paper: black/grey figures; solutions keep colour
\pgfplotsset{curve/.style={black, semithick, smooth, samples=150},
  line2/.style={black, semithick, samples=2},
  sketch/.append style={axis line style={thin, black}, every axis label/.append style={black}}}
\tikzset{region/.style={fill=black!15, draw=none}, regionB/.style={fill=black!30, draw=none},
  rect/.style={fill=black!8, draw=black, thin}, ax3/.style={-{Stealth[length=1.5mm]}, thin, black}}
\colorlet{blue}{black}\colorlet{red}{black}\colorlet{orange}{black}
\fi

% ---- header / footer: page number centred top (bold), [Turn over at bottom right
\makeatletter
\providecommand{\prelimfootlabel}{\prelimtitle\ 2026}
\newcommand{\pp@footer}{{\coverfont\scriptsize \prelimfootlabel\hspace{3em}H2 MATHEMATICS 9758}}
\fancypagestyle{prelim}{\fancyhf{}%
  \chead{\bfseries\thepage}%
  \lfoot{\pp@footer}%
  \rfoot{\ifnum\value{page}<\getpagerefnumber{LastPage}\relax\textbf{[Turn over}\fi}%
  \renewcommand{\headrulewidth}{0pt}\renewcommand{\footrulewidth}{0pt}}
\fancypagestyle{prelimcover}{\fancyhf{}%
  \rfoot{\textbf{[Turn over}}%
  \renewcommand{\headrulewidth}{0pt}\renewcommand{\footrulewidth}{0pt}}
\pagestyle{prelim}

% ---- marks: [n] regular weight, right-aligned at the end of the part ----------
\renewcommand{\mk}[1]{{\unskip\nobreak\hfil\penalty50\hskip1.5em\hbox{}\nobreak\hfill[#1]%
  \parfillskip=0pt\finalhyphendemerits=0\par}}

% ---- questions: bold number at the left margin, text indented; grey [tag] straight after the number (3B)
\newsavebox{\pp@qmeas}
\newcommand{\pp@qlist}[1]{%
  \begin{list}{}{\setlength\leftmargin{2.6em}\setlength\labelwidth{2.1em}\setlength\labelsep{0.5em}%
     \setlength\topsep{0pt}\setlength\itemsep{0pt}\setlength\parsep{0.9ex}%
     \renewcommand\makelabel[1]{\textbf{##1}\hfil}}%
  \item[\thediq]\@endparpenalty\@M\setbox0=\hbox{\di@tag{#1}}\ifdim\wd0=0pt\else{\small\textcolor{black!60}{[\di@tag{#1}]}}\par\nobreak\fi\csname diQbody@#1\endcsname\par\penalty\@M
  \end{list}}
% a question starts on a new page unless it fits whole on the current one
% (questions taller than a page may break, but only between parts)
\renewcommand{\putQ}[1]{\stepcounter{diq}%
  \expandafter\xdef\csname diLab@#1\endcsname{\thediq}%
  \par\setbox\pp@qmeas=\vbox{\pp@qlist{#1}}%
  \ifdim\dimexpr\ht\pp@qmeas+\dp\pp@qmeas\relax<0.95\textheight
    \needspace{\dimexpr\ht\pp@qmeas+\dp\pp@qmeas\relax}\else\needspace{14\baselineskip}\fi
  \zsavepos{q-s\thediq}\pp@qlist{#1}\par\zsavepos{q-e\thediq}\vspace{5ex plus 2ex minus 1ex}}
% part labels (a)/(i) bold, indented under the question text
\renewenvironment{qparts}[1][2.6em]{\begin{list}{}{\setlength\leftmargin{#1}%
  \setlength\labelwidth{\dimexpr#1-0.5em\relax}\setlength\labelsep{0.5em}%
  \setlength\itemsep{1.4ex}\setlength\topsep{1.2ex}\setlength\parsep{0.8ex}%
  \renewcommand\makelabel[1]{\textbf{##1}\hfil}}}{\end{list}}
\makeatother

% ---- cover page --------------------------------------------------------------
\newcommand{\prelimcover}[3]{% #1 total marks, #2 duration, #3 date line
\thispagestyle{prelimcover}%
{\coverfont\fontsize{10}{12}\selectfont\linespread{1}\selectfont
\noindent\begin{minipage}[t]{\linewidth}
  {\bfseries \prelimtitle}\\[0.2ex]
  {\bfseries Compiled by Maximus Pung Jun}\\[1.6ex]
  Higher 2
\end{minipage}\par\vspace{2ex}
\noindent\rule{\linewidth}{1.2pt}\par\vspace{0.4ex}
\noindent{\Large\bfseries MATHEMATICS}\hfill{\Large\bfseries 9758}\par\vspace{1.2ex}
\noindent \prelimsubject\hfill\textbf{#3}\\
QUESTION PAPER\hfill\textbf{#2}\par\vspace{2.5ex}
\noindent\begin{tabular}{@{}p{46mm}l@{}}Additional Materials: & Writing paper\\ & List of Formulae and Results (MF27)\end{tabular}\par\vspace{2ex}
\noindent\rule{\linewidth}{0.6pt}\par\vspace{0.6ex}
{\bfseries READ THESE INSTRUCTIONS FIRST}\par\vspace{2ex}
\setlength{\parskip}{0pt}
Write your name and class on all the work you hand in.\\
Write in dark blue or black pen.\\
You may use an HB pencil for any diagrams or graphs.\\
Do not use staples, paper clips, glue or correction fluid.\par\vspace{2ex}
Answer \textbf{all} the questions.\\
Give non-exact numerical answers correct to 3 significant figures, or 1 decimal place in the case of angles in degrees, unless a different level of accuracy is specified in the question.\\
You are expected to use an approved graphing calculator.\\
Unsupported answers from a graphing calculator are allowed unless a question specifically states otherwise.\\
Where unsupported answers from a graphing calculator are \textbf{not} allowed in a question, you are required to present the mathematical steps using mathematical notations and not calculator commands.\\
You are reminded of the need for clear presentation in your answers.\par\vspace{2ex}
The number of marks is given in brackets [\ ] at the end of each question or part question.\\
The total number of marks for this paper is #1.
\vfill
\noindent\rule{\linewidth}{0.6pt}\par\vspace{0.4ex}
\centerline{This document consists of \textbf{\pageref{LastPage}} printed pages.}\par\vspace{1ex}
}\clearpage}

% ---- solutions header block (page 1 of the worked solutions; same cover style) ----
\newcommand{\prelimsolhead}[1]{% #1 total marks
{\coverfont\fontsize{10}{12}\selectfont\linespread{1}\selectfont
\noindent{\bfseries \prelimtitle}\\[0.2ex]{\bfseries Compiled by Maximus Pung Jun}\\[1.6ex]Higher 2\par\vspace{2ex}
\noindent\rule{\linewidth}{1.2pt}\par\vspace{0.4ex}
\noindent{\Large\bfseries MATHEMATICS}\hfill{\Large\bfseries 9758}\par\vspace{1.2ex}
\noindent \prelimsubject\hfill\textbf{2026}\\
WORKED SOLUTIONS\hfill\textbf{Total marks: #1}\par\vspace{1ex}
\noindent\rule{\linewidth}{0.6pt}\par}\vspace{4ex}}
\fi
  \ifx\housexm\housextestsol \def\prelimkeepcolour{1}% house_prelim.tex: HOUSE COPY of definite_integrals/common/prelim_style.tex; [HOUSE] titles from \housechapter
% =====================================================================
% prelim_style.tex: makes the Test look like a Singapore JC H2 Mathematics
% 9758 prelim paper (studied from ACJC/RI 2026 P1/P2 on the box:
% body Times New Roman + Times-style maths, cover in Arial).
% Fonts: prelim_fonts.tex (real Times New Roman/Arial + Times-style maths).
% Shared by Ch10 (di_style) and Ch11 (de_style; sets \prelimtitle/\prelimsubject). Loaded AFTER the chapter style (overrides its layout for the test only).
% =====================================================================
% [HOUSE] fonts loaded by house.tex
\providecommand{\prelimtitle}{\MakeUppercase{\housechap} TEST}
\providecommand{\prelimsubject}{\housechap}
\geometry{a4paper, top=22mm, bottom=24mm, left=25mm, right=25mm, headsep=6mm, footskip=12mm}
\setlength{\parindent}{0pt}
\setlength{\parskip}{0.9ex plus 0.3ex minus 0.2ex}
\linespread{1.3}
\setlength{\emergencystretch}{3em}
\clubpenalty=10000 \widowpenalty=10000 \interlinepenalty=10000
\AtBeginDocument{\fontsize{12}{14.4}\selectfont}
\renewcommand{\normalsize}{\fontsize{12}{14.4}\selectfont}

% ---- figures: plain black/grey like a printed paper ----------------------
\ifdefined\prelimkeepcolour\else % question paper: black/grey figures; solutions keep colour
\pgfplotsset{curve/.style={black, semithick, smooth, samples=150},
  line2/.style={black, semithick, samples=2},
  sketch/.append style={axis line style={thin, black}, every axis label/.append style={black}}}
\tikzset{region/.style={fill=black!15, draw=none}, regionB/.style={fill=black!30, draw=none},
  rect/.style={fill=black!8, draw=black, thin}, ax3/.style={-{Stealth[length=1.5mm]}, thin, black}}
\colorlet{blue}{black}\colorlet{red}{black}\colorlet{orange}{black}
\fi

% ---- header / footer: page number centred top (bold), [Turn over at bottom right
\makeatletter
\providecommand{\prelimfootlabel}{\prelimtitle\ 2026}
\newcommand{\pp@footer}{{\coverfont\scriptsize \prelimfootlabel\hspace{3em}H2 MATHEMATICS 9758}}
\fancypagestyle{prelim}{\fancyhf{}%
  \chead{\bfseries\thepage}%
  \lfoot{\pp@footer}%
  \rfoot{\ifnum\value{page}<\getpagerefnumber{LastPage}\relax\textbf{[Turn over}\fi}%
  \renewcommand{\headrulewidth}{0pt}\renewcommand{\footrulewidth}{0pt}}
\fancypagestyle{prelimcover}{\fancyhf{}%
  \rfoot{\textbf{[Turn over}}%
  \renewcommand{\headrulewidth}{0pt}\renewcommand{\footrulewidth}{0pt}}
\pagestyle{prelim}

% ---- marks: [n] regular weight, right-aligned at the end of the part ----------
\renewcommand{\mk}[1]{{\unskip\nobreak\hfil\penalty50\hskip1.5em\hbox{}\nobreak\hfill[#1]%
  \parfillskip=0pt\finalhyphendemerits=0\par}}

% ---- questions: bold number at the left margin, text indented; grey [tag] straight after the number (3B)
\newsavebox{\pp@qmeas}
\newcommand{\pp@qlist}[1]{%
  \begin{list}{}{\setlength\leftmargin{2.6em}\setlength\labelwidth{2.1em}\setlength\labelsep{0.5em}%
     \setlength\topsep{0pt}\setlength\itemsep{0pt}\setlength\parsep{0.9ex}%
     \renewcommand\makelabel[1]{\textbf{##1}\hfil}}%
  \item[\thediq]\@endparpenalty\@M\setbox0=\hbox{\di@tag{#1}}\ifdim\wd0=0pt\else{\small\textcolor{black!60}{[\di@tag{#1}]}}\par\nobreak\fi\csname diQbody@#1\endcsname\par\penalty\@M
  \end{list}}
% a question starts on a new page unless it fits whole on the current one
% (questions taller than a page may break, but only between parts)
\renewcommand{\putQ}[1]{\stepcounter{diq}%
  \expandafter\xdef\csname diLab@#1\endcsname{\thediq}%
  \par\setbox\pp@qmeas=\vbox{\pp@qlist{#1}}%
  \ifdim\dimexpr\ht\pp@qmeas+\dp\pp@qmeas\relax<0.95\textheight
    \needspace{\dimexpr\ht\pp@qmeas+\dp\pp@qmeas\relax}\else\needspace{14\baselineskip}\fi
  \zsavepos{q-s\thediq}\pp@qlist{#1}\par\zsavepos{q-e\thediq}\vspace{5ex plus 2ex minus 1ex}}
% part labels (a)/(i) bold, indented under the question text
\renewenvironment{qparts}[1][2.6em]{\begin{list}{}{\setlength\leftmargin{#1}%
  \setlength\labelwidth{\dimexpr#1-0.5em\relax}\setlength\labelsep{0.5em}%
  \setlength\itemsep{1.4ex}\setlength\topsep{1.2ex}\setlength\parsep{0.8ex}%
  \renewcommand\makelabel[1]{\textbf{##1}\hfil}}}{\end{list}}
\makeatother

% ---- cover page --------------------------------------------------------------
\newcommand{\prelimcover}[3]{% #1 total marks, #2 duration, #3 date line
\thispagestyle{prelimcover}%
{\coverfont\fontsize{10}{12}\selectfont\linespread{1}\selectfont
\noindent\begin{minipage}[t]{\linewidth}
  {\bfseries \prelimtitle}\\[0.2ex]
  {\bfseries Compiled by Maximus Pung Jun}\\[1.6ex]
  Higher 2
\end{minipage}\par\vspace{2ex}
\noindent\rule{\linewidth}{1.2pt}\par\vspace{0.4ex}
\noindent{\Large\bfseries MATHEMATICS}\hfill{\Large\bfseries 9758}\par\vspace{1.2ex}
\noindent \prelimsubject\hfill\textbf{#3}\\
QUESTION PAPER\hfill\textbf{#2}\par\vspace{2.5ex}
\noindent\begin{tabular}{@{}p{46mm}l@{}}Additional Materials: & Writing paper\\ & List of Formulae and Results (MF27)\end{tabular}\par\vspace{2ex}
\noindent\rule{\linewidth}{0.6pt}\par\vspace{0.6ex}
{\bfseries READ THESE INSTRUCTIONS FIRST}\par\vspace{2ex}
\setlength{\parskip}{0pt}
Write your name and class on all the work you hand in.\\
Write in dark blue or black pen.\\
You may use an HB pencil for any diagrams or graphs.\\
Do not use staples, paper clips, glue or correction fluid.\par\vspace{2ex}
Answer \textbf{all} the questions.\\
Give non-exact numerical answers correct to 3 significant figures, or 1 decimal place in the case of angles in degrees, unless a different level of accuracy is specified in the question.\\
You are expected to use an approved graphing calculator.\\
Unsupported answers from a graphing calculator are allowed unless a question specifically states otherwise.\\
Where unsupported answers from a graphing calculator are \textbf{not} allowed in a question, you are required to present the mathematical steps using mathematical notations and not calculator commands.\\
You are reminded of the need for clear presentation in your answers.\par\vspace{2ex}
The number of marks is given in brackets [\ ] at the end of each question or part question.\\
The total number of marks for this paper is #1.
\vfill
\noindent\rule{\linewidth}{0.6pt}\par\vspace{0.4ex}
\centerline{This document consists of \textbf{\pageref{LastPage}} printed pages.}\par\vspace{1ex}
}\clearpage}

% ---- solutions header block (page 1 of the worked solutions; same cover style) ----
\newcommand{\prelimsolhead}[1]{% #1 total marks
{\coverfont\fontsize{10}{12}\selectfont\linespread{1}\selectfont
\noindent{\bfseries \prelimtitle}\\[0.2ex]{\bfseries Compiled by Maximus Pung Jun}\\[1.6ex]Higher 2\par\vspace{2ex}
\noindent\rule{\linewidth}{1.2pt}\par\vspace{0.4ex}
\noindent{\Large\bfseries MATHEMATICS}\hfill{\Large\bfseries 9758}\par\vspace{1.2ex}
\noindent \prelimsubject\hfill\textbf{2026}\\
WORKED SOLUTIONS\hfill\textbf{Total marks: #1}\par\vspace{1ex}
\noindent\rule{\linewidth}{0.6pt}\par}\vspace{4ex}}
% house_sol.tex: HOUSE COPY of definite_integrals/common/ri_solstyle.tex (unchanged)
% =====================================================================
% ri_solstyle.tex: worked solutions in the layout of the RI 2026 Year 6 Prelim
% "Solutions with comments" (grail.moe 26612/26613), loaded AFTER prelim_style.
%   * page 1: bold school-style block + bold underlined title; later pages: running
%     title top-left with a rule under it; footer "Page N of M" bottom right (Calibri).
%   * each question: prelim-format question text (bold number, marks right, source tag
%     bottom right), then a bordered table:  Solution | Comments  (bold header row);
%     body rows  label + [marks] (bold, narrow left column) | working | comments,
%     a full-width rule between parts, working left-aligned (fleqn).
% Rows come from the existing sources: \pt in qparts (Ch10 HW), \slabel (Ch10 Test),
% spart (Ch11). Marks under each label are taken from the question's \mk.
% =====================================================================
\makeatletter
\usepackage{lastpage}
\usepackage{adjustbox}
\usepackage{zref-savepos,zref-abspage}
\zref@addprop{savepos}{abspage}
\geometry{a4paper, top=22mm, bottom=24mm, left=22mm, right=22mm, headsep=8mm, footskip=12mm, headheight=24pt}
\providecommand{\risoltitle}{\prelimtitle: Solutions with comments}
\providecommand{\risolschool}{Compiled by Maximus Pung Jun}
\fancypagestyle{risol}{\fancyhf{}%
  \lhead{\fontsize{11}{12}\selectfont 2026 \risoltitle\\[-1.2ex]\rule{\textwidth}{0.6pt}}%
  \rfoot{\calibrifont\fontsize{11}{13}\selectfont Page \textbf{\thepage} of \textbf{\pageref*{LastPage}}}%
  \renewcommand{\headrulewidth}{0pt}\renewcommand{\footrulewidth}{0pt}}
\fancypagestyle{risolfirst}{\fancyhf{}%
  \rfoot{\calibrifont\fontsize{11}{13}\selectfont Page \textbf{\thepage} of \textbf{\pageref*{LastPage}}}%
  \renewcommand{\headrulewidth}{0pt}\renewcommand{\footrulewidth}{0pt}}
\pagestyle{risol}
\newif\ifri@fresh
\newcommand{\risolfront}{\thispagestyle{risolfirst}\global\ri@freshtrue%
  {\noindent\bfseries\linespread{1}\selectfont \risolschool\\ H2 Mathematics (9758)\\ 2026\par}%
  \vspace{4ex}{\noindent\bfseries\underline{2026 \risoltitle}\par}\vspace{3ex}}
% ---- table geometry (RI: label 42pt | working 284pt | comments 142pt of 468pt) ----
\newlength{\ri@L}\setlength{\ri@L}{15mm}
\newlength{\ri@C}\setlength{\ri@C}{40mm}
\newlength{\ri@pad}\setlength{\ri@pad}{1.8mm}
\newif\ifri@sol \newif\ifri@first
\newcount\ri@depth
\def\ri@outer{}
\newtcolorbox{risol}{enhanced, breakable, sharp corners, frame hidden, colback=white,
  boxrule=0pt, boxsep=0pt, left=\dimexpr\ri@L+\ri@pad\relax, right=\dimexpr\ri@C+\ri@pad\relax,
  top=1.2mm, bottom=2.5mm, toptitle=0.6mm, bottomtitle=0.6mm, lefttitle=\ri@pad, righttitle=0pt,
  colbacktitle=white, coltitle=black, fonttitle=\bfseries,
  title={Solution\hfill\makebox[\ri@C][l]{\hspace{\ri@pad}Comments}},
  before skip=2ex, after skip=2ex,
  overlay unbroken and first={\draw[line width=0.5pt] (frame.south west) rectangle (frame.north east);
     \draw[line width=0.5pt] (title.south west) -- (title.south east);
     \draw[line width=0.5pt] ([xshift=\ri@L]interior.north west) -- ([xshift=\ri@L]frame.south west);
     \draw[line width=0.5pt] ([xshift=-\ri@C]frame.north east) -- ([xshift=-\ri@C]frame.south east);},
  overlay middle and last={\draw[line width=0.5pt] (frame.south west) rectangle (frame.north east);
     \draw[line width=0.5pt] ([xshift=\ri@L]frame.north west) -- ([xshift=\ri@L]frame.south west);
     \draw[line width=0.5pt] ([xshift=-\ri@C]frame.north east) -- ([xshift=-\ri@C]frame.south east);},
  after upper={\ri@endrow\protected@write\@auxout{}{\string\ri@settotal{\ri@qid}{\the\ri@sumH}}},
  before upper={\global\ri@sumH=14mm\ri@soltrue\ri@firsttrue\ri@depth=0 \parindent=0pt \parskip=0.6ex
     \predisplaypenalty=100 
     \linespread{1.12}\selectfont\raggedright\@fleqntrue\@mathmargin=0pt
     \abovedisplayskip=5pt plus 2pt\belowdisplayskip=5pt plus 2pt
     \abovedisplayshortskip=3pt plus 1pt\belowdisplayshortskip=4pt plus 2pt
     \ifcsname ri@hp@\ri@qid\endcsname\else\ri@row{}\fi}}
\def\ri@qid{}
% one row: full-width rule (glue leaders, so it vanishes at a page break where the frame
% edge takes its place; not before the first row), then bold label + bold [marks]
% in the left column, on the baseline of the first line of the working
\ExplSyntaxOn
\tl_new:N \g__ri_tmp_tl
\cs_new_protected:Npn \ri_setkey:Nn #1#2
  { \group_begin: \cs_set_eq:NN \chg \use_ii:nn \cs_set_eq:NN \chgm \use_ii:nn
    \cs_set_eq:NN \textbf \use:n \cs_set_eq:NN \mathrm \use:n
    \tl_gset:Nx \g__ri_tmp_tl {#2} \group_end:
    \tl_set_eq:NN #1 \g__ri_tmp_tl \tl_remove_all:Nn #1 {~} }
\cs_set_eq:NN \ri@setkey \ri_setkey:Nn
\ExplSyntaxOff
% ---- page-break rules (positions from the previous run, zref labels in the .aux) ----
% * a row (one part) never breaks across pages: if it does not fit on the rest of the
%   page it starts the next page (forced break between rows); only a row taller than a
%   whole page may break inside;
% * a whole solution that fits on one page is kept together (needspace before the table);
%   a longer one starts where at least its first row fits;
% * a row is at least as tall as its label + [marks].
\newcount\ri@n \newbox\ri@labbox \def\ri@open{0}
\newdimen\ri@ybot \newdimen\ri@ytop \newdimen\ri@h \newdimen\ri@sumH
\AtBeginDocument{\ri@ybot=\dimexpr\paperheight-1in-\voffset-\topmargin-\headheight-\headsep-\textheight\relax
  \ri@ytop=\dimexpr\ri@ybot+\textheight\relax}
\def\ri@pg#1{\zref@extractdefault{#1}{abspage}{-1}}
% \ri@rowh{n}: height of row n in the previous run -> \ri@h (0pt if unknown)
\newcommand{\ri@rowh}[1]{\ri@h=0pt
  \ifnum\ri@pg{ri-s#1}>0 \ifnum\ri@pg{ri-e#1}>0
    \ifnum\ri@pg{ri-s#1}=\ri@pg{ri-e#1}\relax \ri@h=\dimexpr\zposy{ri-s#1}sp-\zposy{ri-e#1}sp\relax
    \else \ri@h=\dimexpr\zposy{ri-s#1}sp-\ri@ybot+\ri@ytop-\zposy{ri-e#1}sp-6mm\relax\fi\fi\fi}
\newcommand{\ri@setfirst}[2]{\expandafter\gdef\csname ri@first@#1\endcsname{#2}}
\newcommand{\ri@setkept}[1]{\expandafter\gdef\csname ri@kept@#1\endcsname{}}
\newcommand{\ri@settotal}[2]{\expandafter\gdef\csname ri@total@#1\endcsname{#2}}
\newcommand{\ri@endrow}{\ifnum\ri@open=1 \par\zsavepos{ri-e\the\ri@n}%
  \ifnum\ri@pg{ri-s\the\ri@n}=\ri@pg{ri-e\the\ri@n}\relax
    \dimen@=\dimexpr\ri@curH sp+1.6mm-\zposy{ri-s\the\ri@n}sp+\zposy{ri-e\the\ri@n}sp\relax
    \ifdim\dimen@>0pt \nobreak\kern\dimen@\penalty\z@\fi\fi
  \ri@rowh{\the\ri@n}\global\advance\ri@sumH\dimexpr\ri@h+3.5mm\relax
  \gdef\ri@open{0}\fi}
\newcommand{\ri@row}[1]{\par\ri@endrow\global\advance\ri@n\@ne\gdef\ri@open{1}\@beginparpenalty\@M
  \ifri@first\global\ri@firstfalse\ri@firstfalse
    \protected@write\@auxout{}{\string\ri@setfirst{\ri@qid}{\the\ri@n}}%
  \else
    \zsavepos{ri-p\the\ri@n}%
    % push the row to the next page if it would not fit on the rest of this one
    \ri@rowh{\the\ri@n}%
    \ifnum\ri@pg{ri-p\the\ri@n}>0 \ifdim\ri@h>0pt
      \ifdim\dimexpr\ri@h+4mm\relax<\dimexpr\textheight-10mm\relax
        \ifdim\dimexpr\ri@h+7mm\relax>\dimexpr\zposy{ri-p\the\ri@n}sp-\ri@ybot\relax
          \penalty-\@M\fi\fi\fi\fi
    \vskip0.6ex\nointerlineskip
    \leaders\hbox{\hskip-\dimexpr\ri@L+\ri@pad\relax\vrule width\dimexpr\linewidth+\ri@L+\ri@C+2\ri@pad\relax height0.5pt depth0pt}\vskip0.5pt
    \nobreak\vskip1.4mm\fi
  \ri@setkey\ri@key{\ri@qid#1}\def\ri@lab{#1}%
  \setbox\ri@labbox=\vtop{\hsize=\dimexpr\ri@L-1.6mm\relax\bfseries\raggedright\parskip=0pt
      \strut\ri@lab\ifcsname ri@mk@\ri@key\endcsname\ifx\ri@lab\@empty\else\\\fi
      [\csname ri@mk@\ri@key\endcsname]\fi\strut\par}%
  \xdef\ri@curH{\number\dimexpr\ht\ri@labbox+\dp\ri@labbox\relax}%
  \nointerlineskip\zsavepos{ri-s\the\ri@n}\nointerlineskip\nobreak
  \vbox to 0pt{\hbox{\hskip-\dimexpr\ri@L+\ri@pad-1.2mm\relax\box\ri@labbox}\vss}\nobreak\prevdepth=\dp\strutbox
  \vskip-\parskip\ignorespaces}
% ---- qparts / \pt: question mode records labels (for the marks); solution mode makes rows
\NewCommandCopy\ri@qparts@begin\qparts \let\ri@qparts@end\endqparts
\renewenvironment{qparts}[1][2.6em]{\advance\ri@depth\@ne
  \ifri@sol\else\ri@qparts@begin[#1]\fi}{\ifri@sol\par\else\ri@qparts@end\fi}
\def\ri@qpartsname{qparts}
\renewcommand{\pt}[1]{\ifnum\ri@depth<2 \gdef\ri@outer{#1}\gdef\ri@cur{#1}\else\expandafter\gdef\expandafter\ri@cur\expandafter{\ri@outer#1}\fi
  \ifri@sol\expandafter\ri@solpt\else\expandafter\ri@qpt\fi{#1}}
\def\ri@qpt#1{\expandafter\gdef\csname ri@hp@\ri@qid\endcsname{}\item[#1]}
\def\ri@solpt#1{\@ifnextchar\begin\ri@ptpeek{\ri@row{\ri@cur}}}
\def\ri@ptpeek\begin#1{\def\ri@tmp{#1}\ifx\ri@tmp\ri@qpartsname
   \expandafter\@firstoftwo\else\expandafter\@secondoftwo\fi
  {\begin{qparts}}{\ri@row{\ri@cur}\begin{#1}}}
% marks: record [n] against the current question/part
\let\ri@mk\mk
\renewcommand{\mk}[1]{\ri@setkey\ri@k{\ri@qid\ri@cur}\ri@setkey\ri@mv{#1}\expandafter\xdef\csname ri@mk@\ri@k\endcsname{\ri@mv}\ri@mk{#1}}
% \slabel (Ch10 Test answers): letter labels set the outer part, roman ones are appended;
% consecutive \slabel calls merge into one row label
\def\ri@romans{(i)(ii)(iii)(iv)(v)(vi)(vii)}
\def\ri@firstpart(#1)#2\@nil{(#1)}
\def\slabel#1{\ifri@sol\expandafter\ri@solslabel\else\expandafter\ri@qslabel\fi{#1}}
\def\ri@qslabel#1{\par\smallskip\noindent\textbf{#1}\par\nobreak}
\def\ri@solslabel#1{\@expandtwoargs\in@{#1}{\ri@romans}%
  \ifin@\expandafter\gdef\expandafter\ri@cur\expandafter{\ri@outer#1}\else\gdef\ri@cur{#1}\xdef\ri@outer{\ri@firstpart#1\@nil}\fi
  \@ifnextchar\slabel{}{\ri@row{\ri@cur}}}
% Ch11 spart{label}: a row; soln: the table
\def\spart#1{\gdef\ri@cur{#1}\ri@row{#1}}\def\endspart{\par}
\def\soln{\@ifnextchar[{\ri@solnopt}{\risol}}\def\ri@solnopt[#1]{\risol}\def\endsoln{\endrisol}
% ---- working that started on the label line is left-aligned in the working column
\AtBeginDocument{\@ifundefined{topalign}{}{%
  \renewenvironment{topalign}{\par\noindent$\displaystyle\begin{aligned}[t]}{\end{aligned}$\par\smallskip}%
  \renewcommand{\topdisplay}[1]{\par\noindent$\displaystyle #1$\par\smallskip}}}
% ---- Ch11 side-by-side (text | sketch): the sketch shrinks to fit its half of the working column
\AtBeginDocument{\@ifundefined{sbsplain}{}{%
  \renewenvironment{sbsplain}[1]{%
    \sbs@start\begin{minipage}[t]{#1\linewidth}\vspace{0pt}\sbs@shift\raggedright%
    \def\sbsnext{\end{minipage}\hfill\begin{minipage}[t]{\dimexpr0.97\linewidth-#1\linewidth\relax}\vspace{0pt}\sbs@shift\centering
      \begin{adjustbox}{max width=\linewidth}}%
  }{\end{adjustbox}\end{minipage}\par}}}
% ---- figures in the working column shrink to fit it (only if wider), colours kept
\newcommand{\rifitfig}[1]{\expandafter\let\csname ri@orig@\string#1\endcsname#1%
  \def#1{\@ifnextchar\bgroup{\ri@fitarg#1}{\ri@fitnoarg#1}}}
\def\ri@fitarg#1#2{\ifri@sol\adjustbox{max width=\linewidth}{\csname ri@orig@\string#1\endcsname{#2}}\else\csname ri@orig@\string#1\endcsname{#2}\fi}
\def\ri@fitnoarg#1{\ifri@sol\adjustbox{max width=\linewidth}{\csname ri@orig@\string#1\endcsname}\else\csname ri@orig@\string#1\endcsname\fi}
\@ifundefined{solfig}{}{\renewcommand{\solfig}[2][]{\par\begin{center}\adjustbox{max width=\linewidth}{\begin{tikzpicture}\begin{axis}[qfig, width=6.6cm, height=5cm, #1]#2\end{axis}\end{tikzpicture}}\end{center}}}
% \ricomment{text}: text in the Comments column of the current row (call right after \pt,
% \slabel or \spart). The row is made at least as tall as the comment.
\newbox\ri@cmtbox
\newcommand{\ricomment}[1]{\par\setbox\ri@cmtbox=\vtop{\hsize=\dimexpr\ri@C-2\ri@pad\relax
    \parskip=0pt\raggedright\small{\color{black!60}\strut#1\strut\par}}%
  \ifdim\dimexpr\ht\ri@cmtbox+\dp\ri@cmtbox\relax>\ri@curH sp\xdef\ri@curH{\number\dimexpr\ht\ri@cmtbox+\dp\ri@cmtbox\relax}\fi
  \nointerlineskip\vbox to 0pt{\hbox to\linewidth{\hskip\dimexpr\linewidth+2\ri@pad\relax\box\ri@cmtbox\hss}\vss}\nobreak\prevdepth=\dp\strutbox
  \vskip-\parskip\ignorespaces}
% \putTFlong{n}: the True/False question n with its answer table allowed to break across
% pages (the body of \defQ{n} must be: intro text, then \begin{tftable}...\end{tftable});
% the solutions copy prints it with tflongtable.
\renewcommand{\putTFlong}[1]{\ri@newq\begingroup\let\tftable\tflongtable\let\endtftable\endtflongtable
  \stepcounter{diq}\expandafter\xdef\csname diLab@#1\endcsname{\thediq}%
  \par\noindent\hbox to 2.6em{\textbf{\thediq}\hfil}\setbox0=\hbox{\di@tag{#1}}\ifdim\wd0=0pt\else{\small\textcolor{black!60}{[\di@tag{#1}]}}\par\noindent\hskip2.6em\fi\ignorespaces\csname diQbody@#1\endcsname\par
  \endgroup\vspace{3ex}}
% ---- question + solution; each question starts on a new page (as RI)
\def\ri@qid{}
\let\ri@putQ\putQ
\renewcommand{\putQ}[1]{\gdef\ri@qid{#1}\global\ri@depth=0 \gdef\ri@cur{}\ri@putQ{#1}}
% \risolQ{n}{body macro name}: question n + its solution body in the table;
% \risolQraw: body already contains its own soln environment (Ch11 Test)
\newcommand{\ri@newq}{\ifri@fresh\global\ri@freshfalse\else\clearpage\fi}
% keep the whole table together if it fits on a page; otherwise make sure at least its
% header and first row fit before starting it
\newcommand{\ri@keep}[1]{\par
  \ifcsname ri@total@#1\endcsname
    % hysteresis: a table kept whole in the previous run stays whole up to textheight-4.5mm,
    % so a total near the limit cannot flip between runs
    \ifdim\csname ri@total@#1\endcsname<\dimexpr\textheight-\ifcsname ri@kept@#1\endcsname4.5\else8\fi mm\relax
      \protected@write\@auxout{}{\string\ri@setkept{#1}}%
      \needspace{\dimexpr\csname ri@total@#1\endcsname+4mm\relax}%
    \else\ifcsname ri@first@#1\endcsname
      \ri@rowh{\csname ri@first@#1\endcsname}%
      \ifdim\ri@h<\dimexpr\textheight-20mm\relax\needspace{\dimexpr\ri@h+18mm\relax}\fi\fi\fi\fi}
\newcommand{\risolQ}[2]{\ri@newq\putQ{#1}\vspace{-3ex}\ri@keep{#1}\begin{risol}\csname #2@#1\endcsname\par\end{risol}}
\newcommand{\risolQraw}[2]{\ri@newq\putQ{#1}\vspace{-3ex}\ri@keep{#1}\csname #2@#1\endcsname\par}
\renewcommand{\putQS}[1]{\risolQ{#1}{diS}}
% end marker only if it fits on the last page (no page holding just the marker)
\renewcommand{\diend}[1][END OF SOLUTIONS]{\par\ifdim\dimexpr\pagegoal-\pagetotal\relax<4\baselineskip\else\vspace{1.5ex}\begin{center}\textbf{#1}\end{center}\fi}
\makeatother
\tfanswerstrue\fi
\fi
):
%   notes: \documentclass[a4paper]{book}
%   others: \documentclass[a4paper,11pt,oneside]{book}
% A HW file uses \providecommand\housemode{hw}; its solutions file has NO \documentclass:
%   \def\housemode{sol}\def\diquestionsonly{1}\input{HW1}  (as DI/DE HWn_solutions.tex)
% Compile: TEXINPUTS=.:/workspace/house_format: xelatex (2 runs; solutions until .aux stable)
% =====================================================================
\providecommand\housemode{hw}
\providecommand\housedir{/workspace/house_format}
\newcommand\housechap{}
\newcommand\housechapter[1]{\renewcommand\housechap{#1}}
\def\housexnotes{notes}\def\housexhw{hw}\def\housexsol{sol}\def\housextest{test}\def\housextestsol{testsol}
\edef\housexm{\housemode}
% house_base.tex: HOUSE COPY of definite_integrals/common/ws_preamble_build.tex
% (= Maximus's own preamble.tex + the DE v3 / AoI additions). Only change: GC font path. [HOUSE]
% WORKSHEET COPY of ../notes/preamble_build.tex (2026-10-09 version), unchanged.
% preamble_build.tex: copy of Maximus's preamble.tex for the DE v3 build.
% The blocks marked [BUILD COPY CHANGE] are the earlier build fixes.
% The blocks marked [V3 ADDITION] are new in this version (example box,
% spacing). Everything else is unchanged from preamble.tex.

% For margins
\usepackage[a4paper, margin = 20mm]{geometry}

\usepackage[Lenny]{fncychap}
% Rejne, Glenn, Lenny

\usepackage{etoc}

% [V3 ADDITION] cancel: used to show the dx's "cancelling" in Section 11.3
\usepackage{cancel}
\renewcommand{\CancelColor}{\color{red}}

% For enumeration & math symbols
\usepackage{enumerate}
\usepackage{amsmath,amsfonts,amssymb,amsthm,graphicx,lipsum}

% For diagrams, plots and boxes
\usepackage{pgfplots}
\pgfplotsset{compat=1.18}
\usepackage{tikz}
\usepackage[most]{tcolorbox}
\tcbuselibrary{skins}
\usepackage{xcolor}
\pgfplotsset{
    standard/.style={clip = false,
    axis line style = thick,
    trig format=rad,
    axis lines = middle,
    every axis x label/.style={at={(current axis.right of origin)},anchor= west},
    every axis y label/.style={at={(current axis.above origin)},anchor=south}
    }
}

% For GC font
\usepackage{fontspec}
% [BUILD COPY CHANGE 1] The TI-84 key font is loaded only if the file exists
% (checked in ./Fonts/... and ../Fonts/...). Otherwise \ti{} falls back to a
% small boxed sans-serif key. Known letters of the TI font are mapped to their
% key names in the fallback: o = [y=], z = [mode], f = [L3].
\newcommand{\tikey}[1]{\tcbox[on line, enhanced, boxsep=0pt, left=2pt, right=2pt,
  top=1pt, bottom=1pt, boxrule=0.5pt, arc=1.5pt, colback=white,
  colframe=black!70]{\sffamily\footnotesize\strut #1}}
\ExplSyntaxOn
\NewDocumentCommand{\tifallback}{m}{
  \str_case:nnF {#1}
    { {o}{\tikey{y=}} {z}{\tikey{mode}} {f}{\tikey{L3}} }
    { \tikey{#1} }
}
\ExplSyntaxOff
\def\tifontcase{0}
\IfFileExists{Fonts/TI-84 Plus CE Key Fonts/TI84PlusCEKeys.ttf}
  {\def\tifontcase{1}}{}
\IfFileExists{../Fonts/TI-84 Plus CE Key Fonts/TI84PlusCEKeys.ttf}
  {\def\tifontcase{2}}{}
\ifcase\tifontcase\relax
  % 0: font not found, use boxed fallback keys
  \newcommand{\ti}[1]{\tifallback{#1}}
\or
  % 1: original location (as in preamble.tex)
  \newfontfamily\tifont[
    Path = {Fonts/TI-84 Plus CE Key Fonts/},
    UprightFont = {TI84PlusCEKeys.ttf}
  ]{}
  \newcommand{\ti}[1]{{\tifont #1}}
\or
  % 2: one level up (when compiling inside draft-ch1/)
  \newfontfamily\tifont[
    Path = {../Fonts/TI-84 Plus CE Key Fonts/},
    UprightFont = {TI84PlusCEKeys.ttf}
  ]{}
  \newcommand{\ti}[1]{{\tifont #1}}
\fi

% [AoI items 17/18/37] GC DISPLAY FONT. Every GC key, command, menu item and
% screen line goes through \gcfont, so the font can be swapped HERE in one place.
%   \gc{fnInt(}      command / menu item / stored variable in the display font
%   \gckey{2nd}      a calculator key (display font in a small key outline)
%   \gcscreen{...}   a calculator screen (display font on a grey-green LCD panel)
% The font file lives in ../Fonts (shared by notes, slides and worksheet).
\def\gcfontfile{VT323-Regular.ttf}
% [HOUSE] GC font: always the absolute copy in \housedir/Fonts (relative Path= breaks under TEXINPUTS)
\providecommand\housedir{/workspace/house_format}
\IfFileExists{\housedir/Fonts/\gcfontfile}{\newfontfamily\gcfont[Path=\housedir/Fonts/, Scale=1.32]{\gcfontfile}}{%
  \newcommand\gcfont{\ttfamily}}
\newcommand{\gc}[1]{\mbox{\gcfont #1}}
\newcommand{\gckey}[1]{\tcbox[on line, enhanced, boxsep=0pt, left=2pt, right=2pt,
  top=0.6pt, bottom=0.2pt, boxrule=0.5pt, arc=1.5pt, colback=black!4,
  colframe=black!65]{\gcfont\strut #1}}
\definecolor{lcd}{RGB}{214,222,204}
\newcommand{\gcscreen}[1]{\tcbox[on line, enhanced, boxsep=1pt, left=4pt, right=4pt,
  top=2pt, bottom=2pt, boxrule=1.2pt, arc=1pt, colback=lcd, colframe=black!70]{%
  \gcfont\begin{tabular}{@{}l@{}}#1\end{tabular}}}
% integral template as shown on the TI-84 Plus CE (MathPrint): \gcint{0}{2}{e^(-X²)}{X}
\newcommand{\gcint}[4]{{\gcfont\raisebox{-0.35ex}{\Large ∫}\hspace{-0.15em}%
  \raisebox{-0.2ex}{\begin{tabular}{@{}l@{}}\footnotesize #2\\[-1.3ex]\footnotesize #1\end{tabular}}%
  \hspace{0.1em}(#3)\mathrm{d}#4}}
% old key macro kept as an alias so nothing breaks
\renewcommand{\ti}[1]{\gckey{#1}}

% SI units
\usepackage{siunitx}

% setspace
%\usepackage{setspace}
%\onehalfspacing

% For hyperlinks within and out of document
\usepackage{hyperref}
\hypersetup{
    colorlinks = true,
    linkcolor = black,
    urlcolor = blue,
}   

% For captions for pics
\usepackage{capt-of}


% tcolorbox custom environments
% Definition ========================================================
\newtcolorbox[auto counter,number within=section]{definition}[2][drop shadow]{%
enhanced, before skip=5mm, after skip=6mm, % [V3 ADDITION] spacing
colback=blue!5!white,colframe=blue!50!black,fonttitle=\bfseries,
title= Definition~\thetcbcounter. #2,#1}

% Theorem ==========================================================
\newtcolorbox[auto counter,number within=section]{theorem}[2][drop shadow]{%
enhanced, before skip=5mm, after skip=6mm, % [V3 ADDITION] spacing
colback=black!5!white,colframe=black,fonttitle=\bfseries,
title= Theorem~\thetcbcounter. #2,#1}

% Caution box ======================================================
\newtcolorbox{caution}[1][]{enhanced,
before skip=5mm,after skip=6mm, % [V3 ADDITION] was 2mm/3mm
boxrule=1pt,left=11mm,right=2mm,top=3mm,bottom=3mm,
colback=yellow!50,
colframe=yellow!20!black,
sharp corners,rounded corners=southeast,arc is angular,arc=3mm,
underlay={%
\path[fill=tcbcolback!80!black] ([yshift=3mm]interior.south east)--++(-0.4,-0.1)--++(0.1,-0.2);
\path[draw=tcbcolframe,shorten <=-0.05mm,shorten >=-0.05mm] ([yshift=3mm]interior.south east)--++(-0.4,-0.1)--++(0.1,-0.2);
\path[fill=yellow!50!black,draw=none] (interior.south west) rectangle node[white]{\Huge\bfseries !} ([xshift=9mm]interior.north west);
},
drop fuzzy shadow,#1}

% Question box ======================================================
\newtcolorbox{question}[1][]{enhanced,
before skip=5mm,after skip=6mm, % [V3 ADDITION] was 2mm/3mm
boxrule=1pt,left=11mm,right=2mm,top=3mm,bottom=3mm,
colback=blue!20!white,
colframe=blue!20!black,
sharp corners,rounded corners=southeast,arc is angular,arc=3mm,
underlay={%
\path[fill=tcbcolback!80!black] ([yshift=3mm]interior.south east)--++(-0.4,-0.1)--++(0.1,-0.2);
\path[draw=tcbcolframe,shorten <=-0.05mm,shorten >=-0.05mm] ([yshift=3mm]interior.south east)--++(-0.4,-0.1)--++(0.1,-0.2);
\path[fill=blue!50!black,draw=none] (interior.south west) rectangle node[white]{\Huge\bfseries ?} ([xshift=9mm]interior.north west);
},
drop fuzzy shadow,#1}

% Information box ======================================================
\newtcolorbox{info}[1][]{enhanced,
before skip=5mm,after skip=6mm, % [V3 ADDITION] was 2mm/3mm
boxrule=1pt,left=11mm,right=2mm,top=3mm,bottom=3mm,
colback=cyan!20!white,
colframe=cyan!20!black,
sharp corners,rounded corners=southeast,arc is angular,arc=3mm,
underlay={%
\path[fill=tcbcolback!80!black] ([yshift=3mm]interior.south east)--++(-0.4,-0.1)--++(0.1,-0.2);
\path[draw=tcbcolframe,shorten <=-0.05mm,shorten >=-0.05mm] ([yshift=3mm]interior.south east)--++(-0.4,-0.1)--++(0.1,-0.2);
\path[fill=cyan!50!black,draw=none] (interior.south west) rectangle node[white]{\Huge\bfseries i} ([xshift=9mm]interior.north west);
},
drop fuzzy shadow,#1}

% Mistake box ======================================================
\newtcolorbox{mistake}[1][]{enhanced,
before skip=5mm,after skip=6mm, % [V3 ADDITION] was 2mm/3mm
boxrule=1pt,left=11mm,right=2mm,top=3mm,bottom=3mm,
colback=red!20!white,
colframe=red!20!black,
sharp corners,rounded corners=southeast,arc is angular,arc=3mm,
underlay={%
\path[fill=tcbcolback!80!black] ([yshift=3mm]interior.south east)--++(-0.4,-0.1)--++(0.1,-0.2);
\path[draw=tcbcolframe,shorten <=-0.05mm,shorten >=-0.05mm] ([yshift=3mm]interior.south east)--++(-0.4,-0.1)--++(0.1,-0.2);
\path[fill=red!50!black,draw=none] (interior.south west) rectangle node[white]{\Huge\bfseries $\times$} ([xshift=9mm]interior.north west);
},
drop fuzzy shadow,#1}



% [V3 ADDITION] Worked example box ==================================
% Own counter, numbered within the section (Example 11.2.1, 11.2.2, ...).
% For per-subsection numbering (11.2.1.1) change `section' to `subsection'. Blue `question' boxes are kept for thought-provoking questions
% only; worked examples always use this box. Usage:
%   \begin{example} ...question text... \end{example}
%   \solution{...}
\newtcolorbox[auto counter,number within=section]{example}[1][]{%
enhanced, breakable=false, before skip=5mm, after skip=4mm,
colback=teal!5!white, colframe=teal!55!black, boxrule=0.8pt,
left=3mm, right=3mm, top=2mm, bottom=2mm,
sharp corners, rounded corners=northwest, arc=3mm,
fonttitle=\bfseries, coltitle=white,
title=Example~\thetcbcounter, attach boxed title to top left={yshift=-2mm, xshift=3mm},
boxed title style={colback=teal!55!black, sharp corners, rounded corners=northwest, arc=2mm},
drop fuzzy shadow, #1}

% [V3 ADDITION] Spacing between sections, subsections, parts, paragraphs =
% (fncychap's chapter heading is untouched; only \section and below.)
\makeatletter
\renewcommand\section{\@startsection{section}{1}{\z@}%
  {-2ex \@plus -.5ex \@minus -.2ex}{3.2ex \@plus .3ex}%
  {\normalfont\LARGE\bfseries\color{blue!50!black}}}
\renewcommand\subsection{\@startsection{subsection}{2}{\z@}%
  {-1.5ex \@plus -.5ex \@minus -.2ex}{2.6ex \@plus .3ex}%
  {\normalfont\Large\bfseries}}
\renewcommand\subsubsection{\@startsection{subsubsection}{3}{\z@}%
  {-4ex \@plus -1ex \@minus -.5ex}{1.6ex \@plus .2ex}%
  {\normalfont\large\bfseries\color{teal!50!black}}}
\makeatother
\setlength{\parindent}{0pt}
\setlength{\parskip}{0.9ex plus 0.3ex minus 0.2ex}

% Headers and Footers
\usepackage{fancyhdr,lastpage}
% [V3 R2] grey left-margin line on the out-of-syllabus pages (item 13)
\usepackage{eso-pic}
% [AoI item 63] thicker line (2.6pt, was 1.6pt) that starts just BELOW the header rule
% (top 14mm of the page left clear). Flag: copy this line into every chapter's preamble.
\newcommand{\oosmargin}{\AddToShipoutPictureBG{\put(12mm,0){\color{gray!75}\rule{2.6pt}{\dimexpr\paperheight-14mm\relax}}}}
\pagestyle{fancy}
\lhead{Maximus Pung Jun}
% [BUILD COPY CHANGE 2] Header bug fix: \chaptername is the word "Chapter",
% so the original header printed "Chapter 1: Chapter". \chaptermark now stores
% the bare chapter title in \leftmark, and the header uses \leftmark.
\renewcommand{\chaptermark}[1]{\markboth{#1}{}}
\rhead{Chapter \thechapter: \leftmark}
% \cfoot{Page \thepage\  of \pageref{LastPage}} % buggy, does not appear on chapter pages.

% [V3 R2] tick and cross markers for mistake boxes (G3)
\newcommand{\wrong}{\textcolor{red!70!black}{\ensuremath{\boldsymbol{\times}}}}
\newcommand{\cmark}{\textcolor{green!50!black}{\ensuremath{\checkmark}}}

\newcommand{\xmark}{\wrong}   % alias: red cross
% [V3 R2] out-of-syllabus helpers (DE notes item 13)
%   \oosmargin    : grey line in the left margin on every following page
%   \oosmarginoff : stop the margin line
%   \oosnote      : grey banner "Out of syllabus. For your interest only."
%   outofsyllabus : environment = new page + margin line + banner; ends with
%                   \clearpage and switches the margin line off again
\newcommand{\oosmarginoff}{\ClearShipoutPictureBG}
\newcommand{\oosnote}[1][Nothing in this section is examined in H2 Mathematics.]{\par\noindent\colorbox{gray!15}{\parbox{\dimexpr\textwidth-2\fboxsep}{\textbf{Out of syllabus. For your interest only.} #1}}\par}
\newenvironment{outofsyllabus}{\clearpage\oosmargin}{\clearpage\oosmarginoff}

% Notation as in Maximus's notes: italic d, display style everywhere (G1).
% Fractions inside exponents stay small, so never use these macros inside ^{}.
\newcommand{\dydx}{\dfrac{\mathrm{d}y}{\mathrm{d}x}}
\newcommand{\dvdx}{\dfrac{\mathrm{d}v}{\mathrm{d}x}}
\newcommand{\dudx}{\dfrac{\mathrm{d}u}{\mathrm{d}x}}
\newcommand{\ddydx}{\dfrac{\mathrm{d}^{2}y}{\mathrm{d}x^{2}}}
\newcommand{\ddt}[1]{\dfrac{\mathrm{d}#1}{\mathrm{d}t}}
% Small syllabus flag (no longer used in the DE notes, kept for other chapters)
\newcommand{\syllabusflag}[1]{\par\noindent{\footnotesize\color{red!70!black}\textbf{Syllabus note.} #1}\par}


% Teacher's copy toggle ================================================
% [BUILD COPY CHANGE 3] \ifx compares the whole word (the original
% \if\conform test was always false). DE_notes_student.tex predefines
% \conform as HideSolutions.
% A student copy defines \conform as HideSolutions BEFORE loading this preamble.
\providecommand\conform{ShowSolutions}
\newcommand\ShowSolutionsKey{ShowSolutions}
% [V3 ADDITION] In the student copy the hidden solution leaves stretchable
% blank space (\vfill), so the free space on each page is shared out as
% working space below the examples instead of collecting at the bottom.
% [AoI item 10] "Solution." on its own line, so that (a), (b), (c) line up below it.
\newcommand\solution[1]{\ifx\conform\ShowSolutionsKey
  \par\noindent\textbf{Solution.}\par\nobreak\vspace{0.5ex}\noindent #1\par\medskip\else\par\vspace{\workspace}\fi}
% [V3 R2] minimum working space left under each example in the student copy
\newlength\workspace\setlength\workspace{3cm plus 1fill}


\newcommand{\NN}{\mathbb{N}}
\newcommand{\ZZ}{\mathbb{Z}}
\newcommand{\QQ}{\mathbb{Q}}
\newcommand{\RR}{\mathbb{R}}
\newcommand{\CC}{\mathbb{C}}


% house_local.tex: HOUSE COPY of definite_integrals/common/ws_preamble_local.tex ([HOUSE] figstyles renamed)
% WORKSHEET COPY of ../notes/preamble_local.tex (version the worksheet was last built with,
% from ../_backup_pre_limitsum, 2026-10-09 17:31). Local so the worksheet no longer depends
% on ../notes; changes: % house_figstyles.tex: HOUSE COPY of definite_integrals/common/ws_figstyles.tex (unchanged)
% WORKSHEET COPY of ../notes/figstyles.tex (2026-10-09 17:31 version). Only change (follow-up):
% thin light axes in `sketch' and `ax3' (marked [2026-10-09 follow-up]).
% =====================================================================
% figstyles.tex: plot styles and figure macros shared by the notes, the
% slides (../slides) and the worksheets (../worksheet). Needs pgfplots,
% the fillbetween library and tikz libraries arrows.meta and calc.
% House rules built in:
%   * `axis on top' in every sketch: shading never covers an axis, and the
%     fill starts at the axis line (items 31, 32, 34);
%   * regions are filled from the curve itself (fill between / closed plot);
%   * 3D solids in the page-13 style (translucent blue, light end caps).
% =====================================================================
\pgfplotsset{
  sketch/.style={clip=false, axis line style={thin, black!45}, every axis label/.append style={black!60}, trig format=rad, % [2026-10-09 follow-up] was axis line style=thick
    axis lines=middle, axis on top,
    every axis x label/.style={at={(current axis.right of origin)},anchor=west},
    every axis y label/.style={at={(current axis.above origin)},anchor=south},
    ticks=none, xlabel=$x$, ylabel=$y$, width=7.6cm, height=6cm,
    every node/.append style={font=\footnotesize}},
  curve/.style={blue, semithick, smooth, samples=150},
  line2/.style={red!75!black, semithick, samples=2},
  % 3D solids of revolution (angles in DEGREES: no `trig format=rad' here)
  solid3d/.style={hide axis, trig format=deg, width=8.4cm, height=6.4cm, clip=false,
    z buffer=sort, colormap={bl}{rgb255=(150,185,235) rgb255=(40,90,180)},
    every node/.append style={font=\footnotesize}},
  solidsurf/.style={surf, shader=interp, opacity=0.75, point meta=y},
  red3d/.style={surf, shader=flat, draw=none, point meta=1,
    colormap={r}{rgb255=(215,50,50) rgb255=(215,50,50)}},
  orange3d/.style={surf, shader=interp, point meta=z,
    colormap={o}{rgb255=(250,210,160) rgb255=(220,120,30)}},
  cap3d/.style={surf, shader=flat, draw=none, point meta=1,
    colormap={c}{rgb255=(185,205,238) rgb255=(185,205,238)}},
  mini3d/.style={solid3d, width=2.6cm, height=2.4cm},
}
\tikzset{
  region/.style={fill=blue!22, draw=none},
  regionB/.style={fill=orange!35, draw=none},
  strip/.style={fill=red!55, draw=red!70!black, thin},
  asym/.style={black, dashed, thin},
  pt/.style={circle, fill=black, inner sep=1.1pt},
  dimarrow/.style={<->, >=Stealth, thin},
  ax3/.style={-{Stealth[length=1.5mm]}, thin, black!50}, % [2026-10-09 follow-up] was -{Stealth[length=2mm]}, thick
}
\definecolor{hl}{RGB}{30,110,200}

% ---------------------------------------------------------------------
% Mini 3D solids for SOLUTIONS (items 54, 62, 66, 74): solid only, no
% washer, every picture 1.5 cm tall; a "big minus small" pair is 4.4 cm wide.
%  \solX[x unit cm]{R(x)}{a}{b}{R(a)}{R(b)}{x-axis end}{y-axis top}{curve colour}
%     solid swept out by y = R(x), a <= x <= b, about the x-axis
%  \solY[y unit cm]{R(y)}{c}{d}{R(c)}{R(d)}{x-axis end}{y-axis top}{curve colour}
%     solid swept out by x = R(y), c <= y <= d, about the y-axis
%     (R is written in terms of the letter x, which here stands for y)
% ---------------------------------------------------------------------
\newcommand{\solX}[9][0.5]{%
\begin{tikzpicture}[baseline=(current bounding box.center)]
\begin{axis}[mini3d, x={(#1cm,-0.06cm)}, y={(0.12cm,0.09cm)}, z={(0cm,0.4cm)}]
  \addplot3[cap3d, domain=0:1, y domain=0:360, samples=2, samples y=40] ({#3}, {x*(#5)*cos(y)}, {x*(#5)*sin(y)});
  \addplot3[solidsurf, opacity=0.5, point meta=z, domain=#3:#4, y domain=0:360, samples=24, samples y=40] ({x}, {(#2)*cos(y)}, {(#2)*sin(y)});
  \addplot3[cap3d, domain=0:1, y domain=0:360, samples=2, samples y=40] ({#4}, {x*(#6)*cos(y)}, {x*(#6)*sin(y)});
  \addplot3[#9, thick, domain=#3:#4, samples=40, samples y=1] ({x}, {0}, {#2});
  \draw[ax3] (axis cs:{min(0,#3)-0.3},0,0) -- (axis cs:#7,0,0);
  \draw[ax3] (axis cs:0,0,{-0.15*(#8)}) -- (axis cs:0,0,#8);
\end{axis}
\end{tikzpicture}}
\newcommand{\solY}[9][0.6]{%
\begin{tikzpicture}[baseline=(current bounding box.center)]
\begin{axis}[mini3d, x={(0.3cm,-0.03cm)}, y={(0.13cm,0.11cm)}, z={(0cm,#1cm)}]
  \addplot3[cap3d, domain=0:1, y domain=0:360, samples=2, samples y=40] ({x*(#5)*cos(y)}, {x*(#5)*sin(y)}, {#3});
  \addplot3[solidsurf, opacity=0.5, point meta=z, domain=#3:#4, y domain=0:360, samples=24, samples y=40] ({(#2)*cos(y)}, {(#2)*sin(y)}, {x});
  \addplot3[cap3d, domain=0:1, y domain=0:360, samples=2, samples y=40] ({x*(#6)*cos(y)}, {x*(#6)*sin(y)}, {#4});
  \addplot3[#9, thick, domain=#3:#4, samples=40, samples y=1] ({#2}, {0}, {x});
  \draw[ax3] (axis cs:0,0,{min(0,#3)-0.15*(#8)}) -- (axis cs:0,0,#8);
  \draw[ax3] (axis cs:-0.3,0,0) -- (axis cs:#7,0,0);
\end{axis}
\end{tikzpicture}}
\newcommand{\minisolo}[1]{\adjustbox{totalheight=1.5cm, max width=4.6cm}{#1}}
\newcommand{\minipair}[2]{\adjustbox{totalheight=1.5cm, max width=4.6cm}{#1\;\raisebox{0.15ex}{\Large$-$}\;#2}}
\newcommand{\minilabel}[1]{\raisebox{0.65cm}{\small #1}\;}
; [2026-10-09 follow-up] \mk marks now bold, as in the notes.
% =====================================================================
% preamble_local.tex: extra packages, styles and macros for the
% Definite Integrals chapter ONLY. Loaded after preamble_build.tex.
% Plot styles and the mini 3D macros are in figstyles.tex (shared with the
% slides and worksheets).
% =====================================================================
\usepgfplotslibrary{fillbetween}
\usepackage{adjustbox} % mini 3D pictures (figstyles.tex)
\usetikzlibrary{arrows.meta, decorations.pathreplacing, calc}
% house_figstyles.tex: HOUSE COPY of definite_integrals/common/ws_figstyles.tex (unchanged)
% WORKSHEET COPY of ../notes/figstyles.tex (2026-10-09 17:31 version). Only change (follow-up):
% thin light axes in `sketch' and `ax3' (marked [2026-10-09 follow-up]).
% =====================================================================
% figstyles.tex: plot styles and figure macros shared by the notes, the
% slides (../slides) and the worksheets (../worksheet). Needs pgfplots,
% the fillbetween library and tikz libraries arrows.meta and calc.
% House rules built in:
%   * `axis on top' in every sketch: shading never covers an axis, and the
%     fill starts at the axis line (items 31, 32, 34);
%   * regions are filled from the curve itself (fill between / closed plot);
%   * 3D solids in the page-13 style (translucent blue, light end caps).
% =====================================================================
\pgfplotsset{
  sketch/.style={clip=false, axis line style={thin, black!45}, every axis label/.append style={black!60}, trig format=rad, % [2026-10-09 follow-up] was axis line style=thick
    axis lines=middle, axis on top,
    every axis x label/.style={at={(current axis.right of origin)},anchor=west},
    every axis y label/.style={at={(current axis.above origin)},anchor=south},
    ticks=none, xlabel=$x$, ylabel=$y$, width=7.6cm, height=6cm,
    every node/.append style={font=\footnotesize}},
  curve/.style={blue, semithick, smooth, samples=150},
  line2/.style={red!75!black, semithick, samples=2},
  % 3D solids of revolution (angles in DEGREES: no `trig format=rad' here)
  solid3d/.style={hide axis, trig format=deg, width=8.4cm, height=6.4cm, clip=false,
    z buffer=sort, colormap={bl}{rgb255=(150,185,235) rgb255=(40,90,180)},
    every node/.append style={font=\footnotesize}},
  solidsurf/.style={surf, shader=interp, opacity=0.75, point meta=y},
  red3d/.style={surf, shader=flat, draw=none, point meta=1,
    colormap={r}{rgb255=(215,50,50) rgb255=(215,50,50)}},
  orange3d/.style={surf, shader=interp, point meta=z,
    colormap={o}{rgb255=(250,210,160) rgb255=(220,120,30)}},
  cap3d/.style={surf, shader=flat, draw=none, point meta=1,
    colormap={c}{rgb255=(185,205,238) rgb255=(185,205,238)}},
  mini3d/.style={solid3d, width=2.6cm, height=2.4cm},
}
\tikzset{
  region/.style={fill=blue!22, draw=none},
  regionB/.style={fill=orange!35, draw=none},
  strip/.style={fill=red!55, draw=red!70!black, thin},
  asym/.style={black, dashed, thin},
  pt/.style={circle, fill=black, inner sep=1.1pt},
  dimarrow/.style={<->, >=Stealth, thin},
  ax3/.style={-{Stealth[length=1.5mm]}, thin, black!50}, % [2026-10-09 follow-up] was -{Stealth[length=2mm]}, thick
}
\definecolor{hl}{RGB}{30,110,200}

% ---------------------------------------------------------------------
% Mini 3D solids for SOLUTIONS (items 54, 62, 66, 74): solid only, no
% washer, every picture 1.5 cm tall; a "big minus small" pair is 4.4 cm wide.
%  \solX[x unit cm]{R(x)}{a}{b}{R(a)}{R(b)}{x-axis end}{y-axis top}{curve colour}
%     solid swept out by y = R(x), a <= x <= b, about the x-axis
%  \solY[y unit cm]{R(y)}{c}{d}{R(c)}{R(d)}{x-axis end}{y-axis top}{curve colour}
%     solid swept out by x = R(y), c <= y <= d, about the y-axis
%     (R is written in terms of the letter x, which here stands for y)
% ---------------------------------------------------------------------
\newcommand{\solX}[9][0.5]{%
\begin{tikzpicture}[baseline=(current bounding box.center)]
\begin{axis}[mini3d, x={(#1cm,-0.06cm)}, y={(0.12cm,0.09cm)}, z={(0cm,0.4cm)}]
  \addplot3[cap3d, domain=0:1, y domain=0:360, samples=2, samples y=40] ({#3}, {x*(#5)*cos(y)}, {x*(#5)*sin(y)});
  \addplot3[solidsurf, opacity=0.5, point meta=z, domain=#3:#4, y domain=0:360, samples=24, samples y=40] ({x}, {(#2)*cos(y)}, {(#2)*sin(y)});
  \addplot3[cap3d, domain=0:1, y domain=0:360, samples=2, samples y=40] ({#4}, {x*(#6)*cos(y)}, {x*(#6)*sin(y)});
  \addplot3[#9, thick, domain=#3:#4, samples=40, samples y=1] ({x}, {0}, {#2});
  \draw[ax3] (axis cs:{min(0,#3)-0.3},0,0) -- (axis cs:#7,0,0);
  \draw[ax3] (axis cs:0,0,{-0.15*(#8)}) -- (axis cs:0,0,#8);
\end{axis}
\end{tikzpicture}}
\newcommand{\solY}[9][0.6]{%
\begin{tikzpicture}[baseline=(current bounding box.center)]
\begin{axis}[mini3d, x={(0.3cm,-0.03cm)}, y={(0.13cm,0.11cm)}, z={(0cm,#1cm)}]
  \addplot3[cap3d, domain=0:1, y domain=0:360, samples=2, samples y=40] ({x*(#5)*cos(y)}, {x*(#5)*sin(y)}, {#3});
  \addplot3[solidsurf, opacity=0.5, point meta=z, domain=#3:#4, y domain=0:360, samples=24, samples y=40] ({(#2)*cos(y)}, {(#2)*sin(y)}, {x});
  \addplot3[cap3d, domain=0:1, y domain=0:360, samples=2, samples y=40] ({x*(#6)*cos(y)}, {x*(#6)*sin(y)}, {#4});
  \addplot3[#9, thick, domain=#3:#4, samples=40, samples y=1] ({#2}, {0}, {x});
  \draw[ax3] (axis cs:0,0,{min(0,#3)-0.15*(#8)}) -- (axis cs:0,0,#8);
  \draw[ax3] (axis cs:-0.3,0,0) -- (axis cs:#7,0,0);
\end{axis}
\end{tikzpicture}}
\newcommand{\minisolo}[1]{\adjustbox{totalheight=1.5cm, max width=4.6cm}{#1}}
\newcommand{\minipair}[2]{\adjustbox{totalheight=1.5cm, max width=4.6cm}{#1\;\raisebox{0.15ex}{\Large$-$}\;#2}}
\newcommand{\minilabel}[1]{\raisebox{0.65cm}{\small #1}\;}


% Small fractions are allowed only in exponents and in the limits of an
% integral (which are superscripts/subscripts); everywhere else use \dfrac.
\newcommand{\dx}{\,\mathrm{d}x}
\newcommand{\dy}{\,\mathrm{d}y}
\newcommand{\evalat}[3]{\left[\,#1\,\right]_{#2}^{#3}}

% [item 78] right-aligned A-level style marks at the end of a question part
\newcommand{\mk}[1]{{\unskip\nobreak\hfil\penalty50\hskip1.5em\hbox{}\nobreak\hfill\textbf{[#1]}%
  \parfillskip=0pt\finalhyphendemerits=0\par}}
% [item 11] visually emphasised step labels in worked solutions
\newcommand{\step}[1]{\textcolor{teal!55!black}{\textbf{Step #1}}\quad}
% (a), (b), (c) part lists inside questions and solutions
\newenvironment{parts}{\begin{enumerate}[(a)]\setlength{\itemsep}{0.4ex}}{\end{enumerate}}
% SEAB-style solution layout (approved 10.5.8 template): working on the
% left, region sketch on the right (mini 3D pictures go under the sketch)
\newcommand{\solpair}[3][0.56]{\noindent\begin{minipage}[t]{#1\textwidth}\vspace{0pt}#2\end{minipage}\hfill
  \begin{minipage}[t]{\dimexpr0.98\textwidth-#1\textwidth\relax}\vspace{0pt}\centering #3\end{minipage}\par}
% minus-diagram style (drafts/minus_style_template.tex, items 29/31/62)
\pgfplotsset{minus/.style={axis lines=middle, width=5.0cm, height=4.2cm, ticks=none,
  xlabel={$x$}, ylabel={$y$}, every axis x label/.style={at={(ticklabel* cs:1)},anchor=west},
  every axis y label/.style={at={(ticklabel* cs:1)},anchor=south}, clip=false, axis on top,
  axis line style={gray!70}, every axis label/.append style={gray},
  every node/.append style={text=gray!80!black, font=\footnotesize}}}
\tikzset{hlreg/.style={fill=hl!35, draw=hl, thick}}

\ifx\housexm\housexnotes
  \ifdefined\housenotesLM\else% house_fonts.tex: HOUSE COPY of definite_integrals/common/prelim_fonts.tex (unchanged)
% =====================================================================
% prelim_fonts.tex: fonts of the real 9758 prelim papers (Word + MathType; pdffonts on
% ACJC 2026 P1/P2, ACJC 2025 P1, JPJC 2026, RVHS 2026, CJC/DHS/TJC/NYJC/HCI/EJC 2025):
%   text Times New Roman; cover/footer Arial; RI page footer Calibri (here Carlito);
%   maths = MathType: Times New Roman letters/digits (italic variables; upright digits,
%   function names, d, e, f), SymbolMT for the integral, sigma, pi and = + < >,
%   scripts 7 pt on 12 pt (about 58 %), tight operator spacing.
% Candidate C (teacher-approved 13 Oct 2026): the real Times New Roman and Arial
% (ttf-mscorefonts-installer) + URW Standard Symbols PS (fonts-urw-base35, a SymbolMT
% clone with the same Symbol encoding) for the integral, sigma, pi, = + < >;
% TeX Gyre Termes Math only for what MathType draws itself (radicals, stretchy brackets,
% fraction bars, arrows). House style: upright d (\mathrm{d}x), upright e and f.
% Loaded once by di_style / de_core (HW sheets) and prelim_style (Tests, solutions).
% =====================================================================
\ifdefined\prelimfontsloaded\endinput\fi
\def\prelimfontsloaded{1}
\edef\pfrestoreat{\catcode`\noexpand\@=\the\catcode`\@\relax}\makeatletter
% MathType script sizes: 7 pt superscripts on 12 pt text (58 %), 5 pt second level
\DeclareMathSizes{12}{12}{7}{5}
\DeclareMathSizes{11}{11}{6.5}{5}
\DeclareMathSizes{10.95}{10.95}{6.5}{5}
\DeclareMathSizes{10}{10}{6}{5}
\DeclareMathSizes{9}{9}{5.5}{5}
\DeclareMathSizes{14.4}{14.4}{8.5}{6}
\setmainfont{Times New Roman}
\setsansfont{Arial}
\newfontfamily\coverfont{Arial}
\newfontfamily\calibrifont{Carlito}
\usepackage{unicode-math}
\setmathfont{TeX Gyre Termes Math}
\setmathfont[range={up/{latin,Latin,num}}]{Times New Roman}
\setmathfont[range={it/{latin,Latin}}]{Times New Roman Italic}
\setmathfont[range={bfup/{latin,Latin,num}}]{Times New Roman Bold}
\setmathfont[range={"002B,"003D,"003C,"003E}]{Standard Symbols PS}
\newfontfamily\pf@sym{Standard Symbols PS}
\newfontfamily\pf@symi[FakeSlant=0.21]{Standard Symbols PS}
% MathType-like spacing (tighter than TeX around = + -)
\AtBeginDocument{\thickmuskip=2.5mu plus 1mu\medmuskip=1.5mu plus 1mu\thinmuskip=3mu
% Symbol integral and sigma (18 pt at 12 pt text, as MathType), slanted Symbol pi
 \renewcommand{\int}{\mathop{\vcenter{\hbox{\pf@sym\fontsize{\fpeval{1.5*\f@size}}{0pt}\selectfont\char"F2}}}\nolimits}%
 \renewcommand{\sum}{\mathop{\vcenter{\hbox{\pf@sym\fontsize{\fpeval{1.5*\f@size}}{0pt}\selectfont\char"E5}}}\limits}%
 \renewcommand{\pi}{\mathord{\text{\pf@symi\char"70}}}}
\pfrestoreat
\fi
  % =====================================================================
% house_tf.tex: True/False table (standing rules 2026-11-05/06/07):
%   * a T/F set is a NUMBERED QUESTION (no "Check Your Understanding" heading);
%   * table: row label (a),(b),... | Statement (left) | T/F column (right);
%   * row spacing slightly larger than default (extrarowheight 6pt, arraystretch 1.18,
%     as the Vectors Basics Worksheet after the 2026-11-07 change), font size unchanged;
%   * question doc: T/F column blank; solutions doc (\tfanswerstrue): bold T or F,
%     optional short reason under the statement in small grey.
% Usage:
%   \begin{tftable}  \tfrow{statement}{T}{optional reason}  ...  \end{tftable}
% Plain tabular: a question never splits across pages (house rule), so the table is kept whole.
% =====================================================================
\RequirePackage{array}
\makeatletter
\newif\iftfanswers \tfanswersfalse
\newcounter{tfrow}
\newlength{\tf@lab}\setlength{\tf@lab}{0.9cm}
\newlength{\tf@ans}\setlength{\tf@ans}{1.6cm}
\newenvironment{tftable}{\setcounter{tfrow}{0}\par\smallskip
  \setlength{\extrarowheight}{6pt}\renewcommand{\arraystretch}{1.18}%
  \noindent\begin{tabular}{|>{\centering\arraybackslash}p{\tf@lab}|p{\dimexpr\linewidth-\tf@lab-\tf@ans-6\tabcolsep-4\arrayrulewidth\relax}|>{\centering\arraybackslash}p{\tf@ans}|}
  \hline &\textbf{Statement}&\textbf{T/F}\\\hline}
 {\end{tabular}\par\smallskip}
\newcommand{\tfrow}[3]{\stepcounter{tfrow}\textbf{(\alph{tfrow})} & #1%
  \iftfanswers\if\relax\detokenize{#3}\relax\else\newline{\small\color{black!60}#3}\fi\fi
  & \iftfanswers\textbf{#2}\fi\\\hline}
% explicit row label (a set that keeps the Basics Worksheet letters, e.g. (e) omitted)
\newcommand{\tfrowl}[4]{\textbf{#1} & #2%
  \iftfanswers\if\relax\detokenize{#4}\relax\else\newline{\small\color{black!60}#4}\fi\fi
  & \iftfanswers\textbf{#3}\fi\\\hline}
% \tfcompact: allowed exception (Vectors I HW2/HW3) - compressed rows so the whole set
% stays on page 1 of the question sheet (extrarowheight 1.5pt, tabcolsep 4pt, 0.85/1.45 cm)
\newif\iftf@compact
\newcommand{\tfcompact}{\tf@compacttrue\setlength{\tf@lab}{0.85cm}\setlength{\tf@ans}{1.45cm}}
\renewenvironment{tftable}{\setcounter{tfrow}{0}\par\smallskip
  \iftf@compact\setlength{\extrarowheight}{1.5pt}\setlength{\tabcolsep}{4pt}\else\setlength{\extrarowheight}{6pt}\fi
  \renewcommand{\arraystretch}{1.18}%
  \noindent\begin{tabular}{|>{\centering\arraybackslash}p{\tf@lab}|p{\dimexpr\linewidth-\tf@lab-\tf@ans-6\tabcolsep-4\arrayrulewidth\relax}|>{\centering\arraybackslash}p{\tf@ans}|}
  \hline &\textbf{Statement}&\textbf{T/F}\\\hline}
 {\end{tabular}\par\smallskip}
% tflongtable: same table, may break across pages with the header repeated (solutions
% copies whose reasons make the set longer than a page). Needs longtable.
\RequirePackage{longtable}
\newenvironment{tflongtable}{\setcounter{tfrow}{0}\par\smallskip
  \setlength{\extrarowheight}{6pt}\renewcommand{\arraystretch}{1.18}%
  \setlength{\LTleft}{0pt}\setlength{\LTright}{0pt}%
  \begin{longtable}{|>{\centering\arraybackslash}p{\tf@lab}|p{\dimexpr\linewidth-\tf@lab-\tf@ans-6\tabcolsep-4\arrayrulewidth\relax}|>{\centering\arraybackslash}p{\tf@ans}|}
  \hline &\textbf{Statement}&\textbf{T/F}\\\hline\endhead}
 {\end{longtable}\par}
\makeatother

  \providecommand{\units}[1]{\ \text{units}^{#1}}
  \raggedbottom
\else
  % house_fonts.tex: HOUSE COPY of definite_integrals/common/prelim_fonts.tex (unchanged)
% =====================================================================
% prelim_fonts.tex: fonts of the real 9758 prelim papers (Word + MathType; pdffonts on
% ACJC 2026 P1/P2, ACJC 2025 P1, JPJC 2026, RVHS 2026, CJC/DHS/TJC/NYJC/HCI/EJC 2025):
%   text Times New Roman; cover/footer Arial; RI page footer Calibri (here Carlito);
%   maths = MathType: Times New Roman letters/digits (italic variables; upright digits,
%   function names, d, e, f), SymbolMT for the integral, sigma, pi and = + < >,
%   scripts 7 pt on 12 pt (about 58 %), tight operator spacing.
% Candidate C (teacher-approved 13 Oct 2026): the real Times New Roman and Arial
% (ttf-mscorefonts-installer) + URW Standard Symbols PS (fonts-urw-base35, a SymbolMT
% clone with the same Symbol encoding) for the integral, sigma, pi, = + < >;
% TeX Gyre Termes Math only for what MathType draws itself (radicals, stretchy brackets,
% fraction bars, arrows). House style: upright d (\mathrm{d}x), upright e and f.
% Loaded once by di_style / de_core (HW sheets) and prelim_style (Tests, solutions).
% =====================================================================
\ifdefined\prelimfontsloaded\endinput\fi
\def\prelimfontsloaded{1}
\edef\pfrestoreat{\catcode`\noexpand\@=\the\catcode`\@\relax}\makeatletter
% MathType script sizes: 7 pt superscripts on 12 pt text (58 %), 5 pt second level
\DeclareMathSizes{12}{12}{7}{5}
\DeclareMathSizes{11}{11}{6.5}{5}
\DeclareMathSizes{10.95}{10.95}{6.5}{5}
\DeclareMathSizes{10}{10}{6}{5}
\DeclareMathSizes{9}{9}{5.5}{5}
\DeclareMathSizes{14.4}{14.4}{8.5}{6}
\setmainfont{Times New Roman}
\setsansfont{Arial}
\newfontfamily\coverfont{Arial}
\newfontfamily\calibrifont{Carlito}
\usepackage{unicode-math}
\setmathfont{TeX Gyre Termes Math}
\setmathfont[range={up/{latin,Latin,num}}]{Times New Roman}
\setmathfont[range={it/{latin,Latin}}]{Times New Roman Italic}
\setmathfont[range={bfup/{latin,Latin,num}}]{Times New Roman Bold}
\setmathfont[range={"002B,"003D,"003C,"003E}]{Standard Symbols PS}
\newfontfamily\pf@sym{Standard Symbols PS}
\newfontfamily\pf@symi[FakeSlant=0.21]{Standard Symbols PS}
% MathType-like spacing (tighter than TeX around = + -)
\AtBeginDocument{\thickmuskip=2.5mu plus 1mu\medmuskip=1.5mu plus 1mu\thinmuskip=3mu
% Symbol integral and sigma (18 pt at 12 pt text, as MathType), slanted Symbol pi
 \renewcommand{\int}{\mathop{\vcenter{\hbox{\pf@sym\fontsize{\fpeval{1.5*\f@size}}{0pt}\selectfont\char"F2}}}\nolimits}%
 \renewcommand{\sum}{\mathop{\vcenter{\hbox{\pf@sym\fontsize{\fpeval{1.5*\f@size}}{0pt}\selectfont\char"E5}}}\limits}%
 \renewcommand{\pi}{\mathord{\text{\pf@symi\char"70}}}}
\pfrestoreat

  % house_hw.tex: HOUSE COPY of definite_integrals/common/di_style.tex; [HOUSE] chapter name from \housechapter{...}
% =====================================================================
% di_style.tex: Definite Integrals homework / test layout.
% Loads the house preamble (local copies of the AoI worksheet files:
% ws_preamble_build.tex, ws_preamble_local.tex -> ws_figstyles.tex),
% then adds the homework layout used by HW1-HW3 and the Test.
% Usage (before \begin{document}):  \ditle{Homework 1}
% =====================================================================
% [HOUSE] base, local and fonts are loaded by house.tex
\usepackage{needspace}
\usepackage{zref-savepos,zref-abspage}\makeatletter\zref@addprop{savepos}{abspage}\makeatother
\usepackage{array}

% ---- shortlist toggle ------------------------------------------------
% \ifshortlist true  -> only the questions in SHORTLIST.md (+ challenges)
% A wrapper file sets it with \def\setshortlist{1} before \input-ing the
% main file; the main file can also be edited to say \shortlisttrue.
\newif\ifshortlist
% Task 5: the FINAL homework (= former shortlist) is the default build. HWn.tex -> HWn.pdf.
% A wrapper that sets \def\setfull{1} builds the archive with all original questions.
\ifdefined\setfull\else\shortlisttrue\fi

\makeatletter
% ---- REVIEW build: highlight amended content ---------------------------
% A wrapper sets \def\setreview{1}. In review mode amended text is shaded
% yellow and carries a hover note (PDF tooltip / sticky note); otherwise the
% macros print the content plainly.
%   \chg{note}{inline content}        short inline change (no line breaks), text mode
%   \chgm{note}{math}                 the same inside math mode
%   \begin{chgblock}{note} ... \end{chgblock}   paragraph/display-level change
\newif\ifreview
\ifdefined\setreview \reviewtrue \fi
% hover notes are raw PDF annotations (xdvipdfmx \special; no extra packages).
% Keep notes plain ASCII; balanced parentheses are fine.
\newsavebox{\chg@box}
\newcommand{\chg@ann}[4]{% width, height, depth, note -> invisible markup annotation with a hover note
  \special{pdf:ann width #1 height #2 depth #3 << /Type /Annot /Subtype /Square /Border [0 0 0] /BS << /W 0 >> /C [1 0.85 0] /CA 0.01
    /T (Review) /Subj (Amendment) /Contents (#4) >>}}
\newcommand{\chg@sticky}[1]{\special{pdf:ann width 12pt height 12pt depth 0pt << /Type /Annot /Subtype /Text /Name /Comment
    /C [1 0.85 0] /T (Review) /Subj (Amendment) /Contents (#1) >>}}
\newcommand{\chg}[2]{\ifreview\sbox\chg@box{\colorbox{yellow!60}{#2}}%
  \leavevmode\rlap{\chg@ann{\the\wd\chg@box}{\the\ht\chg@box}{\the\dp\chg@box}{#1}}\usebox\chg@box\else#2\fi}
% \chgm{note}{math}: the same, for use INSIDE math mode (display style kept)
\newcommand{\chgm}[2]{\ifreview\mathchoice{\chg{#1}{$\displaystyle#2$}}{\chg{#1}{$\textstyle#2$}}{\chg{#1}{$\scriptstyle#2$}}{\chg{#1}{$\scriptscriptstyle#2$}}\else#2\fi}
\newtcolorbox{chgbox}[1]{enhanced, breakable, colback=yellow!35, colframe=yellow!35, boxrule=0pt, sharp corners,
  left=2pt, right=14pt, top=1pt, bottom=1pt, before skip=2pt, after skip=2pt,
  overlay first and middle={}, overlay last={}, overlay unbroken and first={\node[anchor=north east, inner sep=1pt] at (frame.north east) {\chg@sticky{#1}\hbox to 12pt{\vrule width 0pt height 12pt}};}}
% \chgnote{note}: just the sticky-note icon (e.g. inside a figure); nothing outside review mode
\newcommand{\chgnote}[1]{\ifreview\chg@sticky{#1}\hbox to 12pt{\vrule width 0pt height 12pt}\fi}
\newenvironment{chgblock}[1]{\ifreview\begin{chgbox}{#1}\fi}{\ifreview\end{chgbox}\fi}

% ---- header ---------------------------------------------------------
\newcommand\di@title{}
\makeatletter
\newcommand{\ditle}[1]{\gdef\di@title{#1}}
\fancypagestyle{di}{\fancyhf{}%
  \lhead{\small \housechap: \di@title}%
  \rhead{\small Compiled by Maximus Pung Jun}%
  \cfoot{\small\thepage}%
  \renewcommand{\headrulewidth}{0.4pt}}
\pagestyle{di}
% title block on page 1, as in his Typst sheets
\newcommand{\dititleblock}[1][]{\thispagestyle{di}%
  {\LARGE\scshape \housechap: \di@title}\par\smallskip
  {Compiled by Maximus Pung Jun}\par\ifx&#1&\else\smallskip{\small\textcolor{black!60}{#1}}\par\fi\medskip}
\makeatother
\newcommand{\disection}[1]{\par\bigskip\noindent{\Large\scshape #1}\par\medskip}

% ---- questions --------------------------------------------------------
% \Qhead{label}{source tag}: question heading; grey [tag] straight after the number, on the left (rule 3B, 9 Nov 2026)
\newcommand{\Qhead}[2]{\par\needspace{8\baselineskip}\bigskip\noindent
  {\large\bfseries #1}\setbox0=\hbox{#2}\ifdim\wd0=0pt\else\enspace{\small\textcolor{black!60}{[#2]}}\fi\par\smallskip}
% part lists with bold labels as he writes them: \pt{(a)}, \pt{(i)}, \pt{(b) (i)}
\newenvironment{qparts}[1][2.4em]{\begin{list}{}{\setlength\leftmargin{#1}%
  \setlength\labelwidth{\dimexpr#1-0.5em\relax}\setlength\labelsep{0.5em}%
  \setlength\itemsep{0.8ex}\setlength\topsep{0.6ex}\setlength\parsep{0.6ex}%
  \renewcommand\makelabel[1]{\textbf{##1}\hfil}}}{\end{list}}
\newcommand{\pt}[1]{\item[#1]}
% thin horizontal label for a challenge question
\newcommand{\challengebanner}{\par\bigskip\noindent\textcolor{black!50}{\rule[0.5ex]{\linewidth}{0.4pt}}\par
  \noindent{\large\scshape Challenge (optional)}\par}
% upright "units"
\newcommand{\units}[1]{\ \text{units}^{#1}}
% END line
\newcommand{\diend}[1][END OF WORKSHEET]{\par\vspace{2\baselineskip}\begin{center}\textbf{#1}\end{center}}

% ---- question registry (full / shortlist from one source) -------------
% \defQ{n}{tag}{body}   question n of the original sheet
% \defA{n}{answer body} its final answer, as printed in his answer section
% \putQ{n} / \putC{n}   print question n as the next (challenge) question
% \putA{n}              print its answer under the same new label
\newcounter{diq}
\makeatletter
\newcommand{\defQ}[3]{\expandafter\gdef\csname diQtag@#1\endcsname{#2}%
  \expandafter\long\expandafter\gdef\csname diQbody@#1\endcsname{#3}}
\newcommand{\defA}[2]{\expandafter\long\expandafter\gdef\csname diA@#1\endcsname{#2}}
\newcommand{\di@tag}[1]{\csname diQtag@#1\endcsname}% source tag only (no original Q number)
% A question never breaks across pages: it is measured and starts on a new page unless it
% fits whole on the rest of this one (only a question taller than a page may break, and then
% only between parts). Page positions of its start/end (zref labels q-s<n>/q-e<n> in the .aux)
% let the checker verify this.
\newsavebox{\di@qmeas}
\newcommand{\di@placeQ}[1]{\par\setbox\di@qmeas=\vbox{#1}%
  \ifdim\dimexpr\ht\di@qmeas+\dp\di@qmeas\relax<0.95\textheight
    \needspace{\dimexpr\ht\di@qmeas+\dp\di@qmeas+\bigskipamount\relax}\else\needspace{14\baselineskip}\fi
  \zsavepos{q-s\thediq}#1\par\zsavepos{q-e\thediq}}
\newcommand{\putQ}[1]{\stepcounter{diq}%
  \expandafter\xdef\csname diLab@#1\endcsname{\thediq}%
  \di@placeQ{\Qhead{\thediq}{\di@tag{#1}}\csname diQbody@#1\endcsname\par}}
\newcommand{\putC}[1]{\stepcounter{diq}%
  \expandafter\xdef\csname diLab@#1\endcsname{\thediq}%
  \expandafter\gdef\csname diChal@#1\endcsname{1}%
  \di@placeQ{\Qhead{\thediq\quad\normalsize\mdseries\scshape Challenge}{\di@tag{#1}}\csname diQbody@#1\endcsname\par}}
\newcommand{\putA}[1]{\par\needspace{4\baselineskip}%
  \begin{qparts}[2.4em]\item[{\csname diLab@#1\endcsname}]\csname diA@#1\endcsname\end{qparts}\medskip}
% ---- worked solutions: \defS{n}{body}, \putQS{n} prints question n + its solution ----
\newcommand{\defS}[2]{\expandafter\long\expandafter\gdef\csname diS@#1\endcsname{#2}}
\newcommand{\disolhead}{\par\medskip\noindent{\bfseries Solution}\par\nobreak\vspace{-0.6ex}\noindent\textcolor{black!35}{\rule{\linewidth}{0.4pt}}\par\nobreak}
\newcommand{\putQS}[1]{\putQ{#1}\disolhead\csname diS@#1\endcsname\par}
% source of a worked solution; visible only in REVIEW builds. #1: his / pub / pip
\newcommand{\solsrc}[2]{\ifreview\par\noindent\def\di@a{#1}\def\di@b{pip}%
  \ifx\di@a\di@b\colorbox{yellow!60}{\parbox{\dimexpr\linewidth-2\fboxsep-14pt}{\small\textbf{Pip-worked:} #2}}\chgnote{Pip-worked solution. No worked solution by the teacher and no published solution was found. Please check it.}%
  \else{\small\textcolor{black!55}{Source of this solution: #2}}\fi\par\smallskip\fi}
\makeatother
% compact sketch for worked solutions: \solfig[axis options]{plot commands}
\newcommand{\solfig}[2][]{\par\begin{center}\begin{tikzpicture}\begin{axis}[qfig, width=6.6cm, height=5cm, #1]#2\end{axis}\end{tikzpicture}\end{center}}
% a sketch added to an existing worked solution (teacher rule), shaded in REVIEW
\newenvironment{addedsketch}{\begin{chgblock}{Sketch added: teacher rule - every area or volume solution has a sketch.}}{\end{chgblock}}
% answer part list (compact)
\newenvironment{aparts}[1][2.4em]{\begin{qparts}[#1]\setlength\itemsep{0.3ex}}{\end{qparts}}

% ---- figure helpers ---------------------------------------------------------
\tikzset{rect/.style={fill=blue!10, draw=blue!55!black, thin}}
\pgfplotsset{qfig/.style={sketch, every node/.append style={font=\scriptsize}}}
% aligned working that starts on the same line as a part label: \pt{(iii)} \begin{topalign} ... \end{topalign}
\newenvironment{topalign}{\hfill$\displaystyle\begin{aligned}[t]}{\end{aligned}$\hfill\null\par\smallskip}
% one display that starts on the same line as a part label
\newcommand{\topdisplay}[1]{\hfill$\displaystyle #1$\hfill\null\par\smallskip}

  % house_solmacros.tex: HOUSE COPY of de-docs/common/de_style.tex minus \input{de_core} and the DE test titles.
% Adds the Ch11 solution helpers: soln/spart tables, sbs/sbsplain side-by-side, \note, \award, \tagm, \dstart.
% =====================================================================
% de_style.tex: Differential Equations (Ch11) homework / test layout.
% = Ch10 di_style (copied as de_core.tex; header text changed) + the DE
% notation and the framed solution boxes of the v1 DE documents
% (soln / spart, used by the Test worked solutions).
% =====================================================================
\providecommand{\e}{\mathrm{e}}
\providecommand{\cosec}{\operatorname{cosec}}
\providecommand{\dt}{\,\mathrm{d}t}
\providecommand{\abs}[1]{\left|#1\right|}
\providecommand{\dd}[2]{\dfrac{\mathrm{d}#1}{\mathrm{d}#2}}
\providecommand{\dvdx}{\dfrac{\mathrm{d}v}{\mathrm{d}x}}
\providecommand{\dudx}{\dfrac{\mathrm{d}u}{\mathrm{d}x}}
\colorlet{notegrey}{black!60}
\makeatletter
\protected\def\note#1{\ifmmode\text{\small\color{notegrey}\itshape #1}\else{\small\color{notegrey}\itshape #1}\fi}
\newcommand{\award}[1]{\unskip\nobreak\hfill\textbf{\small #1}}
\newcommand{\tagm}[1]{\tag*{\textbf{\small #1}}}
\newcommand{\mfAminus}{\[\begin{aligned}&\int\dfrac{1}{a^{2}-x^{2}}\dx\\&\quad=\dfrac{1}{2a}\ln\abs{\dfrac{a+x}{a-x}}+c\end{aligned}\]}
\newcommand{\mfXminus}{\[\begin{aligned}&\int\dfrac{1}{x^{2}-a^{2}}\dx\\&\quad=\dfrac{1}{2a}\ln\abs{\dfrac{x-a}{x+a}}+c\end{aligned}\]}
\newcommand{\sbs@start}{\if@inlabel\def\sbs@shift{\vspace{-\ht\strutbox}}\else
  \ifvmode\else\par\fi\addvspace{0.4ex}\def\sbs@shift{}\fi\noindent}
\newsavebox{\sbs@box}
\newenvironment{sbs}[1]{%
  \sbs@start\begin{minipage}[t]{#1\linewidth}\vspace{0pt}\sbs@shift\raggedright%
  \def\sbsnext{\end{minipage}\hfill\begin{minipage}[t]{\dimexpr0.97\linewidth-#1\linewidth\relax}\vspace{0pt}\sbs@shift\setlength{\fboxsep}{1.5mm}\begin{lrbox}{\sbs@box}\begin{minipage}[t]{\dimexpr\linewidth-2\fboxsep-1pt\relax}\small\raggedright}%
}{\end{minipage}\end{lrbox}\setlength{\fboxrule}{0.5pt}\fbox{\usebox{\sbs@box}}\end{minipage}\par}
\newenvironment{sbsplain}[1]{%
  \sbs@start\begin{minipage}[t]{#1\linewidth}\vspace{0pt}\sbs@shift\raggedright%
  \def\sbsnext{\end{minipage}\hfill\begin{minipage}[t]{\dimexpr0.97\linewidth-#1\linewidth\relax}\vspace{0pt}\sbs@shift\centering}%
}{\end{minipage}\par}
% ---- framed worked-solution box (v1 DE mark-scheme layout, kept for the Test) ----
\newlength{\solcol}\setlength{\solcol}{15mm}
\newcounter{spartno}
\newcommand{\solskips}{\abovedisplayskip=4pt plus 2pt\belowdisplayskip=4pt plus 2pt\abovedisplayshortskip=2pt\belowdisplayshortskip=3pt}
\newtcolorbox{soln}[1][]{enhanced, breakable, sharp corners,
  colback=white, colframe=black, boxrule=0.6pt,
  boxsep=0pt, left=0pt, right=3mm, top=1.5mm, bottom=2mm,
  before skip=1.2ex, after skip=1.5ex,
  before upper={\setcounter{spartno}{0}\parskip=0.5ex\solskips\@fleqntrue\@mathmargin=0pt\raggedright},
  overlay={\draw[line width=0.4pt] ([xshift=\solcol]frame.north west) -- ([xshift=\solcol]frame.south west);},
  #1}
\newenvironment{spart}[1]{%
  \stepcounter{spartno}%
  \ifnum\value{spartno}>1 \par\vspace{0.2ex}\noindent
    \makebox[\linewidth][l]{\rule{\dimexpr\linewidth+3mm\relax}{0.4pt}}\par\nobreak\vspace{0.3ex}\fi
  \begin{list}{}{\leftmargin=\dimexpr\solcol+2.5mm\relax
    \labelwidth=\solcol \labelsep=2.5mm \itemindent=0pt
    \topsep=0.3ex \parsep=0.5ex \rightmargin=0pt \@beginparpenalty=\@M}%
  \item[{\makebox[\solcol][l]{\strut\hspace{1.5mm}\textbf{#1}}}]%
  \ignorespaces}{\end{list}}
\makeatother
% ---- DE sketches: thin light axes (teacher rule), as the Ch10 'sketch' style ----
\pgfplotsset{dsketch/.style={sketch, axis lines=middle, axis line style={thin, black!45, -{Stealth[length=1.6mm]}},
    ticks=none, clip=false, width=6.5cm, height=4.6cm,
    every axis x label/.style={at={(current axis.right of origin)},anchor=west},
    every axis y label/.style={at={(current axis.above origin)},anchor=south},
    every node/.append style={font=\small}},
  dcurve/.style={black, semithick, smooth, samples=120}}
\tikzset{asym/.style={black, dashed, thin}}
\makeatletter
\newcommand{\dstart}{\leavevmode\vadjust{\nobreak\vskip\dimexpr-\baselineskip-\abovedisplayskip+1pt\relax}}
\newcommand{\sdstart}{\leavevmode\vadjust{\nobreak\vskip\dimexpr-\baselineskip-\abovedisplayskip+4pt\relax}}
\makeatother

  % =====================================================================
% house_tf.tex: True/False table (standing rules 2026-11-05/06/07):
%   * a T/F set is a NUMBERED QUESTION (no "Check Your Understanding" heading);
%   * table: row label (a),(b),... | Statement (left) | T/F column (right);
%   * row spacing slightly larger than default (extrarowheight 6pt, arraystretch 1.18,
%     as the Vectors Basics Worksheet after the 2026-11-07 change), font size unchanged;
%   * question doc: T/F column blank; solutions doc (\tfanswerstrue): bold T or F,
%     optional short reason under the statement in small grey.
% Usage:
%   \begin{tftable}  \tfrow{statement}{T}{optional reason}  ...  \end{tftable}
% Plain tabular: a question never splits across pages (house rule), so the table is kept whole.
% =====================================================================
\RequirePackage{array}
\makeatletter
\newif\iftfanswers \tfanswersfalse
\newcounter{tfrow}
\newlength{\tf@lab}\setlength{\tf@lab}{0.9cm}
\newlength{\tf@ans}\setlength{\tf@ans}{1.6cm}
\newenvironment{tftable}{\setcounter{tfrow}{0}\par\smallskip
  \setlength{\extrarowheight}{6pt}\renewcommand{\arraystretch}{1.18}%
  \noindent\begin{tabular}{|>{\centering\arraybackslash}p{\tf@lab}|p{\dimexpr\linewidth-\tf@lab-\tf@ans-6\tabcolsep-4\arrayrulewidth\relax}|>{\centering\arraybackslash}p{\tf@ans}|}
  \hline &\textbf{Statement}&\textbf{T/F}\\\hline}
 {\end{tabular}\par\smallskip}
\newcommand{\tfrow}[3]{\stepcounter{tfrow}\textbf{(\alph{tfrow})} & #1%
  \iftfanswers\if\relax\detokenize{#3}\relax\else\newline{\small\color{black!60}#3}\fi\fi
  & \iftfanswers\textbf{#2}\fi\\\hline}
% explicit row label (a set that keeps the Basics Worksheet letters, e.g. (e) omitted)
\newcommand{\tfrowl}[4]{\textbf{#1} & #2%
  \iftfanswers\if\relax\detokenize{#4}\relax\else\newline{\small\color{black!60}#4}\fi\fi
  & \iftfanswers\textbf{#3}\fi\\\hline}
% \tfcompact: allowed exception (Vectors I HW2/HW3) - compressed rows so the whole set
% stays on page 1 of the question sheet (extrarowheight 1.5pt, tabcolsep 4pt, 0.85/1.45 cm)
\newif\iftf@compact
\newcommand{\tfcompact}{\tf@compacttrue\setlength{\tf@lab}{0.85cm}\setlength{\tf@ans}{1.45cm}}
\renewenvironment{tftable}{\setcounter{tfrow}{0}\par\smallskip
  \iftf@compact\setlength{\extrarowheight}{1.5pt}\setlength{\tabcolsep}{4pt}\else\setlength{\extrarowheight}{6pt}\fi
  \renewcommand{\arraystretch}{1.18}%
  \noindent\begin{tabular}{|>{\centering\arraybackslash}p{\tf@lab}|p{\dimexpr\linewidth-\tf@lab-\tf@ans-6\tabcolsep-4\arrayrulewidth\relax}|>{\centering\arraybackslash}p{\tf@ans}|}
  \hline &\textbf{Statement}&\textbf{T/F}\\\hline}
 {\end{tabular}\par\smallskip}
% tflongtable: same table, may break across pages with the header repeated (solutions
% copies whose reasons make the set longer than a page). Needs longtable.
\RequirePackage{longtable}
\newenvironment{tflongtable}{\setcounter{tfrow}{0}\par\smallskip
  \setlength{\extrarowheight}{6pt}\renewcommand{\arraystretch}{1.18}%
  \setlength{\LTleft}{0pt}\setlength{\LTright}{0pt}%
  \begin{longtable}{|>{\centering\arraybackslash}p{\tf@lab}|p{\dimexpr\linewidth-\tf@lab-\tf@ans-6\tabcolsep-4\arrayrulewidth\relax}|>{\centering\arraybackslash}p{\tf@ans}|}
  \hline &\textbf{Statement}&\textbf{T/F}\\\hline\endhead}
 {\end{longtable}\par}
\makeatother

  \ifx\housexm\housexsol \def\prelimkeepcolour{1}% house_prelim.tex: HOUSE COPY of definite_integrals/common/prelim_style.tex; [HOUSE] titles from \housechapter
% =====================================================================
% prelim_style.tex: makes the Test look like a Singapore JC H2 Mathematics
% 9758 prelim paper (studied from ACJC/RI 2026 P1/P2 on the box:
% body Times New Roman + Times-style maths, cover in Arial).
% Fonts: prelim_fonts.tex (real Times New Roman/Arial + Times-style maths).
% Shared by Ch10 (di_style) and Ch11 (de_style; sets \prelimtitle/\prelimsubject). Loaded AFTER the chapter style (overrides its layout for the test only).
% =====================================================================
% [HOUSE] fonts loaded by house.tex
\providecommand{\prelimtitle}{\MakeUppercase{\housechap} TEST}
\providecommand{\prelimsubject}{\housechap}
\geometry{a4paper, top=22mm, bottom=24mm, left=25mm, right=25mm, headsep=6mm, footskip=12mm}
\setlength{\parindent}{0pt}
\setlength{\parskip}{0.9ex plus 0.3ex minus 0.2ex}
\linespread{1.3}
\setlength{\emergencystretch}{3em}
\clubpenalty=10000 \widowpenalty=10000 \interlinepenalty=10000
\AtBeginDocument{\fontsize{12}{14.4}\selectfont}
\renewcommand{\normalsize}{\fontsize{12}{14.4}\selectfont}

% ---- figures: plain black/grey like a printed paper ----------------------
\ifdefined\prelimkeepcolour\else % question paper: black/grey figures; solutions keep colour
\pgfplotsset{curve/.style={black, semithick, smooth, samples=150},
  line2/.style={black, semithick, samples=2},
  sketch/.append style={axis line style={thin, black}, every axis label/.append style={black}}}
\tikzset{region/.style={fill=black!15, draw=none}, regionB/.style={fill=black!30, draw=none},
  rect/.style={fill=black!8, draw=black, thin}, ax3/.style={-{Stealth[length=1.5mm]}, thin, black}}
\colorlet{blue}{black}\colorlet{red}{black}\colorlet{orange}{black}
\fi

% ---- header / footer: page number centred top (bold), [Turn over at bottom right
\makeatletter
\providecommand{\prelimfootlabel}{\prelimtitle\ 2026}
\newcommand{\pp@footer}{{\coverfont\scriptsize \prelimfootlabel\hspace{3em}H2 MATHEMATICS 9758}}
\fancypagestyle{prelim}{\fancyhf{}%
  \chead{\bfseries\thepage}%
  \lfoot{\pp@footer}%
  \rfoot{\ifnum\value{page}<\getpagerefnumber{LastPage}\relax\textbf{[Turn over}\fi}%
  \renewcommand{\headrulewidth}{0pt}\renewcommand{\footrulewidth}{0pt}}
\fancypagestyle{prelimcover}{\fancyhf{}%
  \rfoot{\textbf{[Turn over}}%
  \renewcommand{\headrulewidth}{0pt}\renewcommand{\footrulewidth}{0pt}}
\pagestyle{prelim}

% ---- marks: [n] regular weight, right-aligned at the end of the part ----------
\renewcommand{\mk}[1]{{\unskip\nobreak\hfil\penalty50\hskip1.5em\hbox{}\nobreak\hfill[#1]%
  \parfillskip=0pt\finalhyphendemerits=0\par}}

% ---- questions: bold number at the left margin, text indented; grey [tag] straight after the number (3B)
\newsavebox{\pp@qmeas}
\newcommand{\pp@qlist}[1]{%
  \begin{list}{}{\setlength\leftmargin{2.6em}\setlength\labelwidth{2.1em}\setlength\labelsep{0.5em}%
     \setlength\topsep{0pt}\setlength\itemsep{0pt}\setlength\parsep{0.9ex}%
     \renewcommand\makelabel[1]{\textbf{##1}\hfil}}%
  \item[\thediq]\@endparpenalty\@M\setbox0=\hbox{\di@tag{#1}}\ifdim\wd0=0pt\else{\small\textcolor{black!60}{[\di@tag{#1}]}}\par\nobreak\fi\csname diQbody@#1\endcsname\par\penalty\@M
  \end{list}}
% a question starts on a new page unless it fits whole on the current one
% (questions taller than a page may break, but only between parts)
\renewcommand{\putQ}[1]{\stepcounter{diq}%
  \expandafter\xdef\csname diLab@#1\endcsname{\thediq}%
  \par\setbox\pp@qmeas=\vbox{\pp@qlist{#1}}%
  \ifdim\dimexpr\ht\pp@qmeas+\dp\pp@qmeas\relax<0.95\textheight
    \needspace{\dimexpr\ht\pp@qmeas+\dp\pp@qmeas\relax}\else\needspace{14\baselineskip}\fi
  \zsavepos{q-s\thediq}\pp@qlist{#1}\par\zsavepos{q-e\thediq}\vspace{5ex plus 2ex minus 1ex}}
% part labels (a)/(i) bold, indented under the question text
\renewenvironment{qparts}[1][2.6em]{\begin{list}{}{\setlength\leftmargin{#1}%
  \setlength\labelwidth{\dimexpr#1-0.5em\relax}\setlength\labelsep{0.5em}%
  \setlength\itemsep{1.4ex}\setlength\topsep{1.2ex}\setlength\parsep{0.8ex}%
  \renewcommand\makelabel[1]{\textbf{##1}\hfil}}}{\end{list}}
\makeatother

% ---- cover page --------------------------------------------------------------
\newcommand{\prelimcover}[3]{% #1 total marks, #2 duration, #3 date line
\thispagestyle{prelimcover}%
{\coverfont\fontsize{10}{12}\selectfont\linespread{1}\selectfont
\noindent\begin{minipage}[t]{\linewidth}
  {\bfseries \prelimtitle}\\[0.2ex]
  {\bfseries Compiled by Maximus Pung Jun}\\[1.6ex]
  Higher 2
\end{minipage}\par\vspace{2ex}
\noindent\rule{\linewidth}{1.2pt}\par\vspace{0.4ex}
\noindent{\Large\bfseries MATHEMATICS}\hfill{\Large\bfseries 9758}\par\vspace{1.2ex}
\noindent \prelimsubject\hfill\textbf{#3}\\
QUESTION PAPER\hfill\textbf{#2}\par\vspace{2.5ex}
\noindent\begin{tabular}{@{}p{46mm}l@{}}Additional Materials: & Writing paper\\ & List of Formulae and Results (MF27)\end{tabular}\par\vspace{2ex}
\noindent\rule{\linewidth}{0.6pt}\par\vspace{0.6ex}
{\bfseries READ THESE INSTRUCTIONS FIRST}\par\vspace{2ex}
\setlength{\parskip}{0pt}
Write your name and class on all the work you hand in.\\
Write in dark blue or black pen.\\
You may use an HB pencil for any diagrams or graphs.\\
Do not use staples, paper clips, glue or correction fluid.\par\vspace{2ex}
Answer \textbf{all} the questions.\\
Give non-exact numerical answers correct to 3 significant figures, or 1 decimal place in the case of angles in degrees, unless a different level of accuracy is specified in the question.\\
You are expected to use an approved graphing calculator.\\
Unsupported answers from a graphing calculator are allowed unless a question specifically states otherwise.\\
Where unsupported answers from a graphing calculator are \textbf{not} allowed in a question, you are required to present the mathematical steps using mathematical notations and not calculator commands.\\
You are reminded of the need for clear presentation in your answers.\par\vspace{2ex}
The number of marks is given in brackets [\ ] at the end of each question or part question.\\
The total number of marks for this paper is #1.
\vfill
\noindent\rule{\linewidth}{0.6pt}\par\vspace{0.4ex}
\centerline{This document consists of \textbf{\pageref{LastPage}} printed pages.}\par\vspace{1ex}
}\clearpage}

% ---- solutions header block (page 1 of the worked solutions; same cover style) ----
\newcommand{\prelimsolhead}[1]{% #1 total marks
{\coverfont\fontsize{10}{12}\selectfont\linespread{1}\selectfont
\noindent{\bfseries \prelimtitle}\\[0.2ex]{\bfseries Compiled by Maximus Pung Jun}\\[1.6ex]Higher 2\par\vspace{2ex}
\noindent\rule{\linewidth}{1.2pt}\par\vspace{0.4ex}
\noindent{\Large\bfseries MATHEMATICS}\hfill{\Large\bfseries 9758}\par\vspace{1.2ex}
\noindent \prelimsubject\hfill\textbf{2026}\\
WORKED SOLUTIONS\hfill\textbf{Total marks: #1}\par\vspace{1ex}
\noindent\rule{\linewidth}{0.6pt}\par}\vspace{4ex}}
% house_sol.tex: HOUSE COPY of definite_integrals/common/ri_solstyle.tex (unchanged)
% =====================================================================
% ri_solstyle.tex: worked solutions in the layout of the RI 2026 Year 6 Prelim
% "Solutions with comments" (grail.moe 26612/26613), loaded AFTER prelim_style.
%   * page 1: bold school-style block + bold underlined title; later pages: running
%     title top-left with a rule under it; footer "Page N of M" bottom right (Calibri).
%   * each question: prelim-format question text (bold number, marks right, source tag
%     bottom right), then a bordered table:  Solution | Comments  (bold header row);
%     body rows  label + [marks] (bold, narrow left column) | working | comments,
%     a full-width rule between parts, working left-aligned (fleqn).
% Rows come from the existing sources: \pt in qparts (Ch10 HW), \slabel (Ch10 Test),
% spart (Ch11). Marks under each label are taken from the question's \mk.
% =====================================================================
\makeatletter
\usepackage{lastpage}
\usepackage{adjustbox}
\usepackage{zref-savepos,zref-abspage}
\zref@addprop{savepos}{abspage}
\geometry{a4paper, top=22mm, bottom=24mm, left=22mm, right=22mm, headsep=8mm, footskip=12mm, headheight=24pt}
\providecommand{\risoltitle}{\prelimtitle: Solutions with comments}
\providecommand{\risolschool}{Compiled by Maximus Pung Jun}
\fancypagestyle{risol}{\fancyhf{}%
  \lhead{\fontsize{11}{12}\selectfont 2026 \risoltitle\\[-1.2ex]\rule{\textwidth}{0.6pt}}%
  \rfoot{\calibrifont\fontsize{11}{13}\selectfont Page \textbf{\thepage} of \textbf{\pageref*{LastPage}}}%
  \renewcommand{\headrulewidth}{0pt}\renewcommand{\footrulewidth}{0pt}}
\fancypagestyle{risolfirst}{\fancyhf{}%
  \rfoot{\calibrifont\fontsize{11}{13}\selectfont Page \textbf{\thepage} of \textbf{\pageref*{LastPage}}}%
  \renewcommand{\headrulewidth}{0pt}\renewcommand{\footrulewidth}{0pt}}
\pagestyle{risol}
\newif\ifri@fresh
\newcommand{\risolfront}{\thispagestyle{risolfirst}\global\ri@freshtrue%
  {\noindent\bfseries\linespread{1}\selectfont \risolschool\\ H2 Mathematics (9758)\\ 2026\par}%
  \vspace{4ex}{\noindent\bfseries\underline{2026 \risoltitle}\par}\vspace{3ex}}
% ---- table geometry (RI: label 42pt | working 284pt | comments 142pt of 468pt) ----
\newlength{\ri@L}\setlength{\ri@L}{15mm}
\newlength{\ri@C}\setlength{\ri@C}{40mm}
\newlength{\ri@pad}\setlength{\ri@pad}{1.8mm}
\newif\ifri@sol \newif\ifri@first
\newcount\ri@depth
\def\ri@outer{}
\newtcolorbox{risol}{enhanced, breakable, sharp corners, frame hidden, colback=white,
  boxrule=0pt, boxsep=0pt, left=\dimexpr\ri@L+\ri@pad\relax, right=\dimexpr\ri@C+\ri@pad\relax,
  top=1.2mm, bottom=2.5mm, toptitle=0.6mm, bottomtitle=0.6mm, lefttitle=\ri@pad, righttitle=0pt,
  colbacktitle=white, coltitle=black, fonttitle=\bfseries,
  title={Solution\hfill\makebox[\ri@C][l]{\hspace{\ri@pad}Comments}},
  before skip=2ex, after skip=2ex,
  overlay unbroken and first={\draw[line width=0.5pt] (frame.south west) rectangle (frame.north east);
     \draw[line width=0.5pt] (title.south west) -- (title.south east);
     \draw[line width=0.5pt] ([xshift=\ri@L]interior.north west) -- ([xshift=\ri@L]frame.south west);
     \draw[line width=0.5pt] ([xshift=-\ri@C]frame.north east) -- ([xshift=-\ri@C]frame.south east);},
  overlay middle and last={\draw[line width=0.5pt] (frame.south west) rectangle (frame.north east);
     \draw[line width=0.5pt] ([xshift=\ri@L]frame.north west) -- ([xshift=\ri@L]frame.south west);
     \draw[line width=0.5pt] ([xshift=-\ri@C]frame.north east) -- ([xshift=-\ri@C]frame.south east);},
  after upper={\ri@endrow\protected@write\@auxout{}{\string\ri@settotal{\ri@qid}{\the\ri@sumH}}},
  before upper={\global\ri@sumH=14mm\ri@soltrue\ri@firsttrue\ri@depth=0 \parindent=0pt \parskip=0.6ex
     \predisplaypenalty=100 
     \linespread{1.12}\selectfont\raggedright\@fleqntrue\@mathmargin=0pt
     \abovedisplayskip=5pt plus 2pt\belowdisplayskip=5pt plus 2pt
     \abovedisplayshortskip=3pt plus 1pt\belowdisplayshortskip=4pt plus 2pt
     \ifcsname ri@hp@\ri@qid\endcsname\else\ri@row{}\fi}}
\def\ri@qid{}
% one row: full-width rule (glue leaders, so it vanishes at a page break where the frame
% edge takes its place; not before the first row), then bold label + bold [marks]
% in the left column, on the baseline of the first line of the working
\ExplSyntaxOn
\tl_new:N \g__ri_tmp_tl
\cs_new_protected:Npn \ri_setkey:Nn #1#2
  { \group_begin: \cs_set_eq:NN \chg \use_ii:nn \cs_set_eq:NN \chgm \use_ii:nn
    \cs_set_eq:NN \textbf \use:n \cs_set_eq:NN \mathrm \use:n
    \tl_gset:Nx \g__ri_tmp_tl {#2} \group_end:
    \tl_set_eq:NN #1 \g__ri_tmp_tl \tl_remove_all:Nn #1 {~} }
\cs_set_eq:NN \ri@setkey \ri_setkey:Nn
\ExplSyntaxOff
% ---- page-break rules (positions from the previous run, zref labels in the .aux) ----
% * a row (one part) never breaks across pages: if it does not fit on the rest of the
%   page it starts the next page (forced break between rows); only a row taller than a
%   whole page may break inside;
% * a whole solution that fits on one page is kept together (needspace before the table);
%   a longer one starts where at least its first row fits;
% * a row is at least as tall as its label + [marks].
\newcount\ri@n \newbox\ri@labbox \def\ri@open{0}
\newdimen\ri@ybot \newdimen\ri@ytop \newdimen\ri@h \newdimen\ri@sumH
\AtBeginDocument{\ri@ybot=\dimexpr\paperheight-1in-\voffset-\topmargin-\headheight-\headsep-\textheight\relax
  \ri@ytop=\dimexpr\ri@ybot+\textheight\relax}
\def\ri@pg#1{\zref@extractdefault{#1}{abspage}{-1}}
% \ri@rowh{n}: height of row n in the previous run -> \ri@h (0pt if unknown)
\newcommand{\ri@rowh}[1]{\ri@h=0pt
  \ifnum\ri@pg{ri-s#1}>0 \ifnum\ri@pg{ri-e#1}>0
    \ifnum\ri@pg{ri-s#1}=\ri@pg{ri-e#1}\relax \ri@h=\dimexpr\zposy{ri-s#1}sp-\zposy{ri-e#1}sp\relax
    \else \ri@h=\dimexpr\zposy{ri-s#1}sp-\ri@ybot+\ri@ytop-\zposy{ri-e#1}sp-6mm\relax\fi\fi\fi}
\newcommand{\ri@setfirst}[2]{\expandafter\gdef\csname ri@first@#1\endcsname{#2}}
\newcommand{\ri@setkept}[1]{\expandafter\gdef\csname ri@kept@#1\endcsname{}}
\newcommand{\ri@settotal}[2]{\expandafter\gdef\csname ri@total@#1\endcsname{#2}}
\newcommand{\ri@endrow}{\ifnum\ri@open=1 \par\zsavepos{ri-e\the\ri@n}%
  \ifnum\ri@pg{ri-s\the\ri@n}=\ri@pg{ri-e\the\ri@n}\relax
    \dimen@=\dimexpr\ri@curH sp+1.6mm-\zposy{ri-s\the\ri@n}sp+\zposy{ri-e\the\ri@n}sp\relax
    \ifdim\dimen@>0pt \nobreak\kern\dimen@\penalty\z@\fi\fi
  \ri@rowh{\the\ri@n}\global\advance\ri@sumH\dimexpr\ri@h+3.5mm\relax
  \gdef\ri@open{0}\fi}
\newcommand{\ri@row}[1]{\par\ri@endrow\global\advance\ri@n\@ne\gdef\ri@open{1}\@beginparpenalty\@M
  \ifri@first\global\ri@firstfalse\ri@firstfalse
    \protected@write\@auxout{}{\string\ri@setfirst{\ri@qid}{\the\ri@n}}%
  \else
    \zsavepos{ri-p\the\ri@n}%
    % push the row to the next page if it would not fit on the rest of this one
    \ri@rowh{\the\ri@n}%
    \ifnum\ri@pg{ri-p\the\ri@n}>0 \ifdim\ri@h>0pt
      \ifdim\dimexpr\ri@h+4mm\relax<\dimexpr\textheight-10mm\relax
        \ifdim\dimexpr\ri@h+7mm\relax>\dimexpr\zposy{ri-p\the\ri@n}sp-\ri@ybot\relax
          \penalty-\@M\fi\fi\fi\fi
    \vskip0.6ex\nointerlineskip
    \leaders\hbox{\hskip-\dimexpr\ri@L+\ri@pad\relax\vrule width\dimexpr\linewidth+\ri@L+\ri@C+2\ri@pad\relax height0.5pt depth0pt}\vskip0.5pt
    \nobreak\vskip1.4mm\fi
  \ri@setkey\ri@key{\ri@qid#1}\def\ri@lab{#1}%
  \setbox\ri@labbox=\vtop{\hsize=\dimexpr\ri@L-1.6mm\relax\bfseries\raggedright\parskip=0pt
      \strut\ri@lab\ifcsname ri@mk@\ri@key\endcsname\ifx\ri@lab\@empty\else\\\fi
      [\csname ri@mk@\ri@key\endcsname]\fi\strut\par}%
  \xdef\ri@curH{\number\dimexpr\ht\ri@labbox+\dp\ri@labbox\relax}%
  \nointerlineskip\zsavepos{ri-s\the\ri@n}\nointerlineskip\nobreak
  \vbox to 0pt{\hbox{\hskip-\dimexpr\ri@L+\ri@pad-1.2mm\relax\box\ri@labbox}\vss}\nobreak\prevdepth=\dp\strutbox
  \vskip-\parskip\ignorespaces}
% ---- qparts / \pt: question mode records labels (for the marks); solution mode makes rows
\NewCommandCopy\ri@qparts@begin\qparts \let\ri@qparts@end\endqparts
\renewenvironment{qparts}[1][2.6em]{\advance\ri@depth\@ne
  \ifri@sol\else\ri@qparts@begin[#1]\fi}{\ifri@sol\par\else\ri@qparts@end\fi}
\def\ri@qpartsname{qparts}
\renewcommand{\pt}[1]{\ifnum\ri@depth<2 \gdef\ri@outer{#1}\gdef\ri@cur{#1}\else\expandafter\gdef\expandafter\ri@cur\expandafter{\ri@outer#1}\fi
  \ifri@sol\expandafter\ri@solpt\else\expandafter\ri@qpt\fi{#1}}
\def\ri@qpt#1{\expandafter\gdef\csname ri@hp@\ri@qid\endcsname{}\item[#1]}
\def\ri@solpt#1{\@ifnextchar\begin\ri@ptpeek{\ri@row{\ri@cur}}}
\def\ri@ptpeek\begin#1{\def\ri@tmp{#1}\ifx\ri@tmp\ri@qpartsname
   \expandafter\@firstoftwo\else\expandafter\@secondoftwo\fi
  {\begin{qparts}}{\ri@row{\ri@cur}\begin{#1}}}
% marks: record [n] against the current question/part
\let\ri@mk\mk
\renewcommand{\mk}[1]{\ri@setkey\ri@k{\ri@qid\ri@cur}\ri@setkey\ri@mv{#1}\expandafter\xdef\csname ri@mk@\ri@k\endcsname{\ri@mv}\ri@mk{#1}}
% \slabel (Ch10 Test answers): letter labels set the outer part, roman ones are appended;
% consecutive \slabel calls merge into one row label
\def\ri@romans{(i)(ii)(iii)(iv)(v)(vi)(vii)}
\def\ri@firstpart(#1)#2\@nil{(#1)}
\def\slabel#1{\ifri@sol\expandafter\ri@solslabel\else\expandafter\ri@qslabel\fi{#1}}
\def\ri@qslabel#1{\par\smallskip\noindent\textbf{#1}\par\nobreak}
\def\ri@solslabel#1{\@expandtwoargs\in@{#1}{\ri@romans}%
  \ifin@\expandafter\gdef\expandafter\ri@cur\expandafter{\ri@outer#1}\else\gdef\ri@cur{#1}\xdef\ri@outer{\ri@firstpart#1\@nil}\fi
  \@ifnextchar\slabel{}{\ri@row{\ri@cur}}}
% Ch11 spart{label}: a row; soln: the table
\def\spart#1{\gdef\ri@cur{#1}\ri@row{#1}}\def\endspart{\par}
\def\soln{\@ifnextchar[{\ri@solnopt}{\risol}}\def\ri@solnopt[#1]{\risol}\def\endsoln{\endrisol}
% ---- working that started on the label line is left-aligned in the working column
\AtBeginDocument{\@ifundefined{topalign}{}{%
  \renewenvironment{topalign}{\par\noindent$\displaystyle\begin{aligned}[t]}{\end{aligned}$\par\smallskip}%
  \renewcommand{\topdisplay}[1]{\par\noindent$\displaystyle #1$\par\smallskip}}}
% ---- Ch11 side-by-side (text | sketch): the sketch shrinks to fit its half of the working column
\AtBeginDocument{\@ifundefined{sbsplain}{}{%
  \renewenvironment{sbsplain}[1]{%
    \sbs@start\begin{minipage}[t]{#1\linewidth}\vspace{0pt}\sbs@shift\raggedright%
    \def\sbsnext{\end{minipage}\hfill\begin{minipage}[t]{\dimexpr0.97\linewidth-#1\linewidth\relax}\vspace{0pt}\sbs@shift\centering
      \begin{adjustbox}{max width=\linewidth}}%
  }{\end{adjustbox}\end{minipage}\par}}}
% ---- figures in the working column shrink to fit it (only if wider), colours kept
\newcommand{\rifitfig}[1]{\expandafter\let\csname ri@orig@\string#1\endcsname#1%
  \def#1{\@ifnextchar\bgroup{\ri@fitarg#1}{\ri@fitnoarg#1}}}
\def\ri@fitarg#1#2{\ifri@sol\adjustbox{max width=\linewidth}{\csname ri@orig@\string#1\endcsname{#2}}\else\csname ri@orig@\string#1\endcsname{#2}\fi}
\def\ri@fitnoarg#1{\ifri@sol\adjustbox{max width=\linewidth}{\csname ri@orig@\string#1\endcsname}\else\csname ri@orig@\string#1\endcsname\fi}
\@ifundefined{solfig}{}{\renewcommand{\solfig}[2][]{\par\begin{center}\adjustbox{max width=\linewidth}{\begin{tikzpicture}\begin{axis}[qfig, width=6.6cm, height=5cm, #1]#2\end{axis}\end{tikzpicture}}\end{center}}}
% \ricomment{text}: text in the Comments column of the current row (call right after \pt,
% \slabel or \spart). The row is made at least as tall as the comment.
\newbox\ri@cmtbox
\newcommand{\ricomment}[1]{\par\setbox\ri@cmtbox=\vtop{\hsize=\dimexpr\ri@C-2\ri@pad\relax
    \parskip=0pt\raggedright\small{\color{black!60}\strut#1\strut\par}}%
  \ifdim\dimexpr\ht\ri@cmtbox+\dp\ri@cmtbox\relax>\ri@curH sp\xdef\ri@curH{\number\dimexpr\ht\ri@cmtbox+\dp\ri@cmtbox\relax}\fi
  \nointerlineskip\vbox to 0pt{\hbox to\linewidth{\hskip\dimexpr\linewidth+2\ri@pad\relax\box\ri@cmtbox\hss}\vss}\nobreak\prevdepth=\dp\strutbox
  \vskip-\parskip\ignorespaces}
% \putTFlong{n}: the True/False question n with its answer table allowed to break across
% pages (the body of \defQ{n} must be: intro text, then \begin{tftable}...\end{tftable});
% the solutions copy prints it with tflongtable.
\renewcommand{\putTFlong}[1]{\ri@newq\begingroup\let\tftable\tflongtable\let\endtftable\endtflongtable
  \stepcounter{diq}\expandafter\xdef\csname diLab@#1\endcsname{\thediq}%
  \par\noindent\hbox to 2.6em{\textbf{\thediq}\hfil}\setbox0=\hbox{\di@tag{#1}}\ifdim\wd0=0pt\else{\small\textcolor{black!60}{[\di@tag{#1}]}}\par\noindent\hskip2.6em\fi\ignorespaces\csname diQbody@#1\endcsname\par
  \endgroup\vspace{3ex}}
% ---- question + solution; each question starts on a new page (as RI)
\def\ri@qid{}
\let\ri@putQ\putQ
\renewcommand{\putQ}[1]{\gdef\ri@qid{#1}\global\ri@depth=0 \gdef\ri@cur{}\ri@putQ{#1}}
% \risolQ{n}{body macro name}: question n + its solution body in the table;
% \risolQraw: body already contains its own soln environment (Ch11 Test)
\newcommand{\ri@newq}{\ifri@fresh\global\ri@freshfalse\else\clearpage\fi}
% keep the whole table together if it fits on a page; otherwise make sure at least its
% header and first row fit before starting it
\newcommand{\ri@keep}[1]{\par
  \ifcsname ri@total@#1\endcsname
    % hysteresis: a table kept whole in the previous run stays whole up to textheight-4.5mm,
    % so a total near the limit cannot flip between runs
    \ifdim\csname ri@total@#1\endcsname<\dimexpr\textheight-\ifcsname ri@kept@#1\endcsname4.5\else8\fi mm\relax
      \protected@write\@auxout{}{\string\ri@setkept{#1}}%
      \needspace{\dimexpr\csname ri@total@#1\endcsname+4mm\relax}%
    \else\ifcsname ri@first@#1\endcsname
      \ri@rowh{\csname ri@first@#1\endcsname}%
      \ifdim\ri@h<\dimexpr\textheight-20mm\relax\needspace{\dimexpr\ri@h+18mm\relax}\fi\fi\fi\fi}
\newcommand{\risolQ}[2]{\ri@newq\putQ{#1}\vspace{-3ex}\ri@keep{#1}\begin{risol}\csname #2@#1\endcsname\par\end{risol}}
\newcommand{\risolQraw}[2]{\ri@newq\putQ{#1}\vspace{-3ex}\ri@keep{#1}\csname #2@#1\endcsname\par}
\renewcommand{\putQS}[1]{\risolQ{#1}{diS}}
% end marker only if it fits on the last page (no page holding just the marker)
\renewcommand{\diend}[1][END OF SOLUTIONS]{\par\ifdim\dimexpr\pagegoal-\pagetotal\relax<4\baselineskip\else\vspace{1.5ex}\begin{center}\textbf{#1}\end{center}\fi}
\makeatother
\tfanswerstrue\fi
  \ifx\housexm\housextest % house_prelim.tex: HOUSE COPY of definite_integrals/common/prelim_style.tex; [HOUSE] titles from \housechapter
% =====================================================================
% prelim_style.tex: makes the Test look like a Singapore JC H2 Mathematics
% 9758 prelim paper (studied from ACJC/RI 2026 P1/P2 on the box:
% body Times New Roman + Times-style maths, cover in Arial).
% Fonts: prelim_fonts.tex (real Times New Roman/Arial + Times-style maths).
% Shared by Ch10 (di_style) and Ch11 (de_style; sets \prelimtitle/\prelimsubject). Loaded AFTER the chapter style (overrides its layout for the test only).
% =====================================================================
% [HOUSE] fonts loaded by house.tex
\providecommand{\prelimtitle}{\MakeUppercase{\housechap} TEST}
\providecommand{\prelimsubject}{\housechap}
\geometry{a4paper, top=22mm, bottom=24mm, left=25mm, right=25mm, headsep=6mm, footskip=12mm}
\setlength{\parindent}{0pt}
\setlength{\parskip}{0.9ex plus 0.3ex minus 0.2ex}
\linespread{1.3}
\setlength{\emergencystretch}{3em}
\clubpenalty=10000 \widowpenalty=10000 \interlinepenalty=10000
\AtBeginDocument{\fontsize{12}{14.4}\selectfont}
\renewcommand{\normalsize}{\fontsize{12}{14.4}\selectfont}

% ---- figures: plain black/grey like a printed paper ----------------------
\ifdefined\prelimkeepcolour\else % question paper: black/grey figures; solutions keep colour
\pgfplotsset{curve/.style={black, semithick, smooth, samples=150},
  line2/.style={black, semithick, samples=2},
  sketch/.append style={axis line style={thin, black}, every axis label/.append style={black}}}
\tikzset{region/.style={fill=black!15, draw=none}, regionB/.style={fill=black!30, draw=none},
  rect/.style={fill=black!8, draw=black, thin}, ax3/.style={-{Stealth[length=1.5mm]}, thin, black}}
\colorlet{blue}{black}\colorlet{red}{black}\colorlet{orange}{black}
\fi

% ---- header / footer: page number centred top (bold), [Turn over at bottom right
\makeatletter
\providecommand{\prelimfootlabel}{\prelimtitle\ 2026}
\newcommand{\pp@footer}{{\coverfont\scriptsize \prelimfootlabel\hspace{3em}H2 MATHEMATICS 9758}}
\fancypagestyle{prelim}{\fancyhf{}%
  \chead{\bfseries\thepage}%
  \lfoot{\pp@footer}%
  \rfoot{\ifnum\value{page}<\getpagerefnumber{LastPage}\relax\textbf{[Turn over}\fi}%
  \renewcommand{\headrulewidth}{0pt}\renewcommand{\footrulewidth}{0pt}}
\fancypagestyle{prelimcover}{\fancyhf{}%
  \rfoot{\textbf{[Turn over}}%
  \renewcommand{\headrulewidth}{0pt}\renewcommand{\footrulewidth}{0pt}}
\pagestyle{prelim}

% ---- marks: [n] regular weight, right-aligned at the end of the part ----------
\renewcommand{\mk}[1]{{\unskip\nobreak\hfil\penalty50\hskip1.5em\hbox{}\nobreak\hfill[#1]%
  \parfillskip=0pt\finalhyphendemerits=0\par}}

% ---- questions: bold number at the left margin, text indented; grey [tag] straight after the number (3B)
\newsavebox{\pp@qmeas}
\newcommand{\pp@qlist}[1]{%
  \begin{list}{}{\setlength\leftmargin{2.6em}\setlength\labelwidth{2.1em}\setlength\labelsep{0.5em}%
     \setlength\topsep{0pt}\setlength\itemsep{0pt}\setlength\parsep{0.9ex}%
     \renewcommand\makelabel[1]{\textbf{##1}\hfil}}%
  \item[\thediq]\@endparpenalty\@M\setbox0=\hbox{\di@tag{#1}}\ifdim\wd0=0pt\else{\small\textcolor{black!60}{[\di@tag{#1}]}}\par\nobreak\fi\csname diQbody@#1\endcsname\par\penalty\@M
  \end{list}}
% a question starts on a new page unless it fits whole on the current one
% (questions taller than a page may break, but only between parts)
\renewcommand{\putQ}[1]{\stepcounter{diq}%
  \expandafter\xdef\csname diLab@#1\endcsname{\thediq}%
  \par\setbox\pp@qmeas=\vbox{\pp@qlist{#1}}%
  \ifdim\dimexpr\ht\pp@qmeas+\dp\pp@qmeas\relax<0.95\textheight
    \needspace{\dimexpr\ht\pp@qmeas+\dp\pp@qmeas\relax}\else\needspace{14\baselineskip}\fi
  \zsavepos{q-s\thediq}\pp@qlist{#1}\par\zsavepos{q-e\thediq}\vspace{5ex plus 2ex minus 1ex}}
% part labels (a)/(i) bold, indented under the question text
\renewenvironment{qparts}[1][2.6em]{\begin{list}{}{\setlength\leftmargin{#1}%
  \setlength\labelwidth{\dimexpr#1-0.5em\relax}\setlength\labelsep{0.5em}%
  \setlength\itemsep{1.4ex}\setlength\topsep{1.2ex}\setlength\parsep{0.8ex}%
  \renewcommand\makelabel[1]{\textbf{##1}\hfil}}}{\end{list}}
\makeatother

% ---- cover page --------------------------------------------------------------
\newcommand{\prelimcover}[3]{% #1 total marks, #2 duration, #3 date line
\thispagestyle{prelimcover}%
{\coverfont\fontsize{10}{12}\selectfont\linespread{1}\selectfont
\noindent\begin{minipage}[t]{\linewidth}
  {\bfseries \prelimtitle}\\[0.2ex]
  {\bfseries Compiled by Maximus Pung Jun}\\[1.6ex]
  Higher 2
\end{minipage}\par\vspace{2ex}
\noindent\rule{\linewidth}{1.2pt}\par\vspace{0.4ex}
\noindent{\Large\bfseries MATHEMATICS}\hfill{\Large\bfseries 9758}\par\vspace{1.2ex}
\noindent \prelimsubject\hfill\textbf{#3}\\
QUESTION PAPER\hfill\textbf{#2}\par\vspace{2.5ex}
\noindent\begin{tabular}{@{}p{46mm}l@{}}Additional Materials: & Writing paper\\ & List of Formulae and Results (MF27)\end{tabular}\par\vspace{2ex}
\noindent\rule{\linewidth}{0.6pt}\par\vspace{0.6ex}
{\bfseries READ THESE INSTRUCTIONS FIRST}\par\vspace{2ex}
\setlength{\parskip}{0pt}
Write your name and class on all the work you hand in.\\
Write in dark blue or black pen.\\
You may use an HB pencil for any diagrams or graphs.\\
Do not use staples, paper clips, glue or correction fluid.\par\vspace{2ex}
Answer \textbf{all} the questions.\\
Give non-exact numerical answers correct to 3 significant figures, or 1 decimal place in the case of angles in degrees, unless a different level of accuracy is specified in the question.\\
You are expected to use an approved graphing calculator.\\
Unsupported answers from a graphing calculator are allowed unless a question specifically states otherwise.\\
Where unsupported answers from a graphing calculator are \textbf{not} allowed in a question, you are required to present the mathematical steps using mathematical notations and not calculator commands.\\
You are reminded of the need for clear presentation in your answers.\par\vspace{2ex}
The number of marks is given in brackets [\ ] at the end of each question or part question.\\
The total number of marks for this paper is #1.
\vfill
\noindent\rule{\linewidth}{0.6pt}\par\vspace{0.4ex}
\centerline{This document consists of \textbf{\pageref{LastPage}} printed pages.}\par\vspace{1ex}
}\clearpage}

% ---- solutions header block (page 1 of the worked solutions; same cover style) ----
\newcommand{\prelimsolhead}[1]{% #1 total marks
{\coverfont\fontsize{10}{12}\selectfont\linespread{1}\selectfont
\noindent{\bfseries \prelimtitle}\\[0.2ex]{\bfseries Compiled by Maximus Pung Jun}\\[1.6ex]Higher 2\par\vspace{2ex}
\noindent\rule{\linewidth}{1.2pt}\par\vspace{0.4ex}
\noindent{\Large\bfseries MATHEMATICS}\hfill{\Large\bfseries 9758}\par\vspace{1.2ex}
\noindent \prelimsubject\hfill\textbf{2026}\\
WORKED SOLUTIONS\hfill\textbf{Total marks: #1}\par\vspace{1ex}
\noindent\rule{\linewidth}{0.6pt}\par}\vspace{4ex}}
\fi
  \ifx\housexm\housextestsol \def\prelimkeepcolour{1}% house_prelim.tex: HOUSE COPY of definite_integrals/common/prelim_style.tex; [HOUSE] titles from \housechapter
% =====================================================================
% prelim_style.tex: makes the Test look like a Singapore JC H2 Mathematics
% 9758 prelim paper (studied from ACJC/RI 2026 P1/P2 on the box:
% body Times New Roman + Times-style maths, cover in Arial).
% Fonts: prelim_fonts.tex (real Times New Roman/Arial + Times-style maths).
% Shared by Ch10 (di_style) and Ch11 (de_style; sets \prelimtitle/\prelimsubject). Loaded AFTER the chapter style (overrides its layout for the test only).
% =====================================================================
% [HOUSE] fonts loaded by house.tex
\providecommand{\prelimtitle}{\MakeUppercase{\housechap} TEST}
\providecommand{\prelimsubject}{\housechap}
\geometry{a4paper, top=22mm, bottom=24mm, left=25mm, right=25mm, headsep=6mm, footskip=12mm}
\setlength{\parindent}{0pt}
\setlength{\parskip}{0.9ex plus 0.3ex minus 0.2ex}
\linespread{1.3}
\setlength{\emergencystretch}{3em}
\clubpenalty=10000 \widowpenalty=10000 \interlinepenalty=10000
\AtBeginDocument{\fontsize{12}{14.4}\selectfont}
\renewcommand{\normalsize}{\fontsize{12}{14.4}\selectfont}

% ---- figures: plain black/grey like a printed paper ----------------------
\ifdefined\prelimkeepcolour\else % question paper: black/grey figures; solutions keep colour
\pgfplotsset{curve/.style={black, semithick, smooth, samples=150},
  line2/.style={black, semithick, samples=2},
  sketch/.append style={axis line style={thin, black}, every axis label/.append style={black}}}
\tikzset{region/.style={fill=black!15, draw=none}, regionB/.style={fill=black!30, draw=none},
  rect/.style={fill=black!8, draw=black, thin}, ax3/.style={-{Stealth[length=1.5mm]}, thin, black}}
\colorlet{blue}{black}\colorlet{red}{black}\colorlet{orange}{black}
\fi

% ---- header / footer: page number centred top (bold), [Turn over at bottom right
\makeatletter
\providecommand{\prelimfootlabel}{\prelimtitle\ 2026}
\newcommand{\pp@footer}{{\coverfont\scriptsize \prelimfootlabel\hspace{3em}H2 MATHEMATICS 9758}}
\fancypagestyle{prelim}{\fancyhf{}%
  \chead{\bfseries\thepage}%
  \lfoot{\pp@footer}%
  \rfoot{\ifnum\value{page}<\getpagerefnumber{LastPage}\relax\textbf{[Turn over}\fi}%
  \renewcommand{\headrulewidth}{0pt}\renewcommand{\footrulewidth}{0pt}}
\fancypagestyle{prelimcover}{\fancyhf{}%
  \rfoot{\textbf{[Turn over}}%
  \renewcommand{\headrulewidth}{0pt}\renewcommand{\footrulewidth}{0pt}}
\pagestyle{prelim}

% ---- marks: [n] regular weight, right-aligned at the end of the part ----------
\renewcommand{\mk}[1]{{\unskip\nobreak\hfil\penalty50\hskip1.5em\hbox{}\nobreak\hfill[#1]%
  \parfillskip=0pt\finalhyphendemerits=0\par}}

% ---- questions: bold number at the left margin, text indented; grey [tag] straight after the number (3B)
\newsavebox{\pp@qmeas}
\newcommand{\pp@qlist}[1]{%
  \begin{list}{}{\setlength\leftmargin{2.6em}\setlength\labelwidth{2.1em}\setlength\labelsep{0.5em}%
     \setlength\topsep{0pt}\setlength\itemsep{0pt}\setlength\parsep{0.9ex}%
     \renewcommand\makelabel[1]{\textbf{##1}\hfil}}%
  \item[\thediq]\@endparpenalty\@M\setbox0=\hbox{\di@tag{#1}}\ifdim\wd0=0pt\else{\small\textcolor{black!60}{[\di@tag{#1}]}}\par\nobreak\fi\csname diQbody@#1\endcsname\par\penalty\@M
  \end{list}}
% a question starts on a new page unless it fits whole on the current one
% (questions taller than a page may break, but only between parts)
\renewcommand{\putQ}[1]{\stepcounter{diq}%
  \expandafter\xdef\csname diLab@#1\endcsname{\thediq}%
  \par\setbox\pp@qmeas=\vbox{\pp@qlist{#1}}%
  \ifdim\dimexpr\ht\pp@qmeas+\dp\pp@qmeas\relax<0.95\textheight
    \needspace{\dimexpr\ht\pp@qmeas+\dp\pp@qmeas\relax}\else\needspace{14\baselineskip}\fi
  \zsavepos{q-s\thediq}\pp@qlist{#1}\par\zsavepos{q-e\thediq}\vspace{5ex plus 2ex minus 1ex}}
% part labels (a)/(i) bold, indented under the question text
\renewenvironment{qparts}[1][2.6em]{\begin{list}{}{\setlength\leftmargin{#1}%
  \setlength\labelwidth{\dimexpr#1-0.5em\relax}\setlength\labelsep{0.5em}%
  \setlength\itemsep{1.4ex}\setlength\topsep{1.2ex}\setlength\parsep{0.8ex}%
  \renewcommand\makelabel[1]{\textbf{##1}\hfil}}}{\end{list}}
\makeatother

% ---- cover page --------------------------------------------------------------
\newcommand{\prelimcover}[3]{% #1 total marks, #2 duration, #3 date line
\thispagestyle{prelimcover}%
{\coverfont\fontsize{10}{12}\selectfont\linespread{1}\selectfont
\noindent\begin{minipage}[t]{\linewidth}
  {\bfseries \prelimtitle}\\[0.2ex]
  {\bfseries Compiled by Maximus Pung Jun}\\[1.6ex]
  Higher 2
\end{minipage}\par\vspace{2ex}
\noindent\rule{\linewidth}{1.2pt}\par\vspace{0.4ex}
\noindent{\Large\bfseries MATHEMATICS}\hfill{\Large\bfseries 9758}\par\vspace{1.2ex}
\noindent \prelimsubject\hfill\textbf{#3}\\
QUESTION PAPER\hfill\textbf{#2}\par\vspace{2.5ex}
\noindent\begin{tabular}{@{}p{46mm}l@{}}Additional Materials: & Writing paper\\ & List of Formulae and Results (MF27)\end{tabular}\par\vspace{2ex}
\noindent\rule{\linewidth}{0.6pt}\par\vspace{0.6ex}
{\bfseries READ THESE INSTRUCTIONS FIRST}\par\vspace{2ex}
\setlength{\parskip}{0pt}
Write your name and class on all the work you hand in.\\
Write in dark blue or black pen.\\
You may use an HB pencil for any diagrams or graphs.\\
Do not use staples, paper clips, glue or correction fluid.\par\vspace{2ex}
Answer \textbf{all} the questions.\\
Give non-exact numerical answers correct to 3 significant figures, or 1 decimal place in the case of angles in degrees, unless a different level of accuracy is specified in the question.\\
You are expected to use an approved graphing calculator.\\
Unsupported answers from a graphing calculator are allowed unless a question specifically states otherwise.\\
Where unsupported answers from a graphing calculator are \textbf{not} allowed in a question, you are required to present the mathematical steps using mathematical notations and not calculator commands.\\
You are reminded of the need for clear presentation in your answers.\par\vspace{2ex}
The number of marks is given in brackets [\ ] at the end of each question or part question.\\
The total number of marks for this paper is #1.
\vfill
\noindent\rule{\linewidth}{0.6pt}\par\vspace{0.4ex}
\centerline{This document consists of \textbf{\pageref{LastPage}} printed pages.}\par\vspace{1ex}
}\clearpage}

% ---- solutions header block (page 1 of the worked solutions; same cover style) ----
\newcommand{\prelimsolhead}[1]{% #1 total marks
{\coverfont\fontsize{10}{12}\selectfont\linespread{1}\selectfont
\noindent{\bfseries \prelimtitle}\\[0.2ex]{\bfseries Compiled by Maximus Pung Jun}\\[1.6ex]Higher 2\par\vspace{2ex}
\noindent\rule{\linewidth}{1.2pt}\par\vspace{0.4ex}
\noindent{\Large\bfseries MATHEMATICS}\hfill{\Large\bfseries 9758}\par\vspace{1.2ex}
\noindent \prelimsubject\hfill\textbf{2026}\\
WORKED SOLUTIONS\hfill\textbf{Total marks: #1}\par\vspace{1ex}
\noindent\rule{\linewidth}{0.6pt}\par}\vspace{4ex}}
% house_sol.tex: HOUSE COPY of definite_integrals/common/ri_solstyle.tex (unchanged)
% =====================================================================
% ri_solstyle.tex: worked solutions in the layout of the RI 2026 Year 6 Prelim
% "Solutions with comments" (grail.moe 26612/26613), loaded AFTER prelim_style.
%   * page 1: bold school-style block + bold underlined title; later pages: running
%     title top-left with a rule under it; footer "Page N of M" bottom right (Calibri).
%   * each question: prelim-format question text (bold number, marks right, source tag
%     bottom right), then a bordered table:  Solution | Comments  (bold header row);
%     body rows  label + [marks] (bold, narrow left column) | working | comments,
%     a full-width rule between parts, working left-aligned (fleqn).
% Rows come from the existing sources: \pt in qparts (Ch10 HW), \slabel (Ch10 Test),
% spart (Ch11). Marks under each label are taken from the question's \mk.
% =====================================================================
\makeatletter
\usepackage{lastpage}
\usepackage{adjustbox}
\usepackage{zref-savepos,zref-abspage}
\zref@addprop{savepos}{abspage}
\geometry{a4paper, top=22mm, bottom=24mm, left=22mm, right=22mm, headsep=8mm, footskip=12mm, headheight=24pt}
\providecommand{\risoltitle}{\prelimtitle: Solutions with comments}
\providecommand{\risolschool}{Compiled by Maximus Pung Jun}
\fancypagestyle{risol}{\fancyhf{}%
  \lhead{\fontsize{11}{12}\selectfont 2026 \risoltitle\\[-1.2ex]\rule{\textwidth}{0.6pt}}%
  \rfoot{\calibrifont\fontsize{11}{13}\selectfont Page \textbf{\thepage} of \textbf{\pageref*{LastPage}}}%
  \renewcommand{\headrulewidth}{0pt}\renewcommand{\footrulewidth}{0pt}}
\fancypagestyle{risolfirst}{\fancyhf{}%
  \rfoot{\calibrifont\fontsize{11}{13}\selectfont Page \textbf{\thepage} of \textbf{\pageref*{LastPage}}}%
  \renewcommand{\headrulewidth}{0pt}\renewcommand{\footrulewidth}{0pt}}
\pagestyle{risol}
\newif\ifri@fresh
\newcommand{\risolfront}{\thispagestyle{risolfirst}\global\ri@freshtrue%
  {\noindent\bfseries\linespread{1}\selectfont \risolschool\\ H2 Mathematics (9758)\\ 2026\par}%
  \vspace{4ex}{\noindent\bfseries\underline{2026 \risoltitle}\par}\vspace{3ex}}
% ---- table geometry (RI: label 42pt | working 284pt | comments 142pt of 468pt) ----
\newlength{\ri@L}\setlength{\ri@L}{15mm}
\newlength{\ri@C}\setlength{\ri@C}{40mm}
\newlength{\ri@pad}\setlength{\ri@pad}{1.8mm}
\newif\ifri@sol \newif\ifri@first
\newcount\ri@depth
\def\ri@outer{}
\newtcolorbox{risol}{enhanced, breakable, sharp corners, frame hidden, colback=white,
  boxrule=0pt, boxsep=0pt, left=\dimexpr\ri@L+\ri@pad\relax, right=\dimexpr\ri@C+\ri@pad\relax,
  top=1.2mm, bottom=2.5mm, toptitle=0.6mm, bottomtitle=0.6mm, lefttitle=\ri@pad, righttitle=0pt,
  colbacktitle=white, coltitle=black, fonttitle=\bfseries,
  title={Solution\hfill\makebox[\ri@C][l]{\hspace{\ri@pad}Comments}},
  before skip=2ex, after skip=2ex,
  overlay unbroken and first={\draw[line width=0.5pt] (frame.south west) rectangle (frame.north east);
     \draw[line width=0.5pt] (title.south west) -- (title.south east);
     \draw[line width=0.5pt] ([xshift=\ri@L]interior.north west) -- ([xshift=\ri@L]frame.south west);
     \draw[line width=0.5pt] ([xshift=-\ri@C]frame.north east) -- ([xshift=-\ri@C]frame.south east);},
  overlay middle and last={\draw[line width=0.5pt] (frame.south west) rectangle (frame.north east);
     \draw[line width=0.5pt] ([xshift=\ri@L]frame.north west) -- ([xshift=\ri@L]frame.south west);
     \draw[line width=0.5pt] ([xshift=-\ri@C]frame.north east) -- ([xshift=-\ri@C]frame.south east);},
  after upper={\ri@endrow\protected@write\@auxout{}{\string\ri@settotal{\ri@qid}{\the\ri@sumH}}},
  before upper={\global\ri@sumH=14mm\ri@soltrue\ri@firsttrue\ri@depth=0 \parindent=0pt \parskip=0.6ex
     \predisplaypenalty=100 
     \linespread{1.12}\selectfont\raggedright\@fleqntrue\@mathmargin=0pt
     \abovedisplayskip=5pt plus 2pt\belowdisplayskip=5pt plus 2pt
     \abovedisplayshortskip=3pt plus 1pt\belowdisplayshortskip=4pt plus 2pt
     \ifcsname ri@hp@\ri@qid\endcsname\else\ri@row{}\fi}}
\def\ri@qid{}
% one row: full-width rule (glue leaders, so it vanishes at a page break where the frame
% edge takes its place; not before the first row), then bold label + bold [marks]
% in the left column, on the baseline of the first line of the working
\ExplSyntaxOn
\tl_new:N \g__ri_tmp_tl
\cs_new_protected:Npn \ri_setkey:Nn #1#2
  { \group_begin: \cs_set_eq:NN \chg \use_ii:nn \cs_set_eq:NN \chgm \use_ii:nn
    \cs_set_eq:NN \textbf \use:n \cs_set_eq:NN \mathrm \use:n
    \tl_gset:Nx \g__ri_tmp_tl {#2} \group_end:
    \tl_set_eq:NN #1 \g__ri_tmp_tl \tl_remove_all:Nn #1 {~} }
\cs_set_eq:NN \ri@setkey \ri_setkey:Nn
\ExplSyntaxOff
% ---- page-break rules (positions from the previous run, zref labels in the .aux) ----
% * a row (one part) never breaks across pages: if it does not fit on the rest of the
%   page it starts the next page (forced break between rows); only a row taller than a
%   whole page may break inside;
% * a whole solution that fits on one page is kept together (needspace before the table);
%   a longer one starts where at least its first row fits;
% * a row is at least as tall as its label + [marks].
\newcount\ri@n \newbox\ri@labbox \def\ri@open{0}
\newdimen\ri@ybot \newdimen\ri@ytop \newdimen\ri@h \newdimen\ri@sumH
\AtBeginDocument{\ri@ybot=\dimexpr\paperheight-1in-\voffset-\topmargin-\headheight-\headsep-\textheight\relax
  \ri@ytop=\dimexpr\ri@ybot+\textheight\relax}
\def\ri@pg#1{\zref@extractdefault{#1}{abspage}{-1}}
% \ri@rowh{n}: height of row n in the previous run -> \ri@h (0pt if unknown)
\newcommand{\ri@rowh}[1]{\ri@h=0pt
  \ifnum\ri@pg{ri-s#1}>0 \ifnum\ri@pg{ri-e#1}>0
    \ifnum\ri@pg{ri-s#1}=\ri@pg{ri-e#1}\relax \ri@h=\dimexpr\zposy{ri-s#1}sp-\zposy{ri-e#1}sp\relax
    \else \ri@h=\dimexpr\zposy{ri-s#1}sp-\ri@ybot+\ri@ytop-\zposy{ri-e#1}sp-6mm\relax\fi\fi\fi}
\newcommand{\ri@setfirst}[2]{\expandafter\gdef\csname ri@first@#1\endcsname{#2}}
\newcommand{\ri@setkept}[1]{\expandafter\gdef\csname ri@kept@#1\endcsname{}}
\newcommand{\ri@settotal}[2]{\expandafter\gdef\csname ri@total@#1\endcsname{#2}}
\newcommand{\ri@endrow}{\ifnum\ri@open=1 \par\zsavepos{ri-e\the\ri@n}%
  \ifnum\ri@pg{ri-s\the\ri@n}=\ri@pg{ri-e\the\ri@n}\relax
    \dimen@=\dimexpr\ri@curH sp+1.6mm-\zposy{ri-s\the\ri@n}sp+\zposy{ri-e\the\ri@n}sp\relax
    \ifdim\dimen@>0pt \nobreak\kern\dimen@\penalty\z@\fi\fi
  \ri@rowh{\the\ri@n}\global\advance\ri@sumH\dimexpr\ri@h+3.5mm\relax
  \gdef\ri@open{0}\fi}
\newcommand{\ri@row}[1]{\par\ri@endrow\global\advance\ri@n\@ne\gdef\ri@open{1}\@beginparpenalty\@M
  \ifri@first\global\ri@firstfalse\ri@firstfalse
    \protected@write\@auxout{}{\string\ri@setfirst{\ri@qid}{\the\ri@n}}%
  \else
    \zsavepos{ri-p\the\ri@n}%
    % push the row to the next page if it would not fit on the rest of this one
    \ri@rowh{\the\ri@n}%
    \ifnum\ri@pg{ri-p\the\ri@n}>0 \ifdim\ri@h>0pt
      \ifdim\dimexpr\ri@h+4mm\relax<\dimexpr\textheight-10mm\relax
        \ifdim\dimexpr\ri@h+7mm\relax>\dimexpr\zposy{ri-p\the\ri@n}sp-\ri@ybot\relax
          \penalty-\@M\fi\fi\fi\fi
    \vskip0.6ex\nointerlineskip
    \leaders\hbox{\hskip-\dimexpr\ri@L+\ri@pad\relax\vrule width\dimexpr\linewidth+\ri@L+\ri@C+2\ri@pad\relax height0.5pt depth0pt}\vskip0.5pt
    \nobreak\vskip1.4mm\fi
  \ri@setkey\ri@key{\ri@qid#1}\def\ri@lab{#1}%
  \setbox\ri@labbox=\vtop{\hsize=\dimexpr\ri@L-1.6mm\relax\bfseries\raggedright\parskip=0pt
      \strut\ri@lab\ifcsname ri@mk@\ri@key\endcsname\ifx\ri@lab\@empty\else\\\fi
      [\csname ri@mk@\ri@key\endcsname]\fi\strut\par}%
  \xdef\ri@curH{\number\dimexpr\ht\ri@labbox+\dp\ri@labbox\relax}%
  \nointerlineskip\zsavepos{ri-s\the\ri@n}\nointerlineskip\nobreak
  \vbox to 0pt{\hbox{\hskip-\dimexpr\ri@L+\ri@pad-1.2mm\relax\box\ri@labbox}\vss}\nobreak\prevdepth=\dp\strutbox
  \vskip-\parskip\ignorespaces}
% ---- qparts / \pt: question mode records labels (for the marks); solution mode makes rows
\NewCommandCopy\ri@qparts@begin\qparts \let\ri@qparts@end\endqparts
\renewenvironment{qparts}[1][2.6em]{\advance\ri@depth\@ne
  \ifri@sol\else\ri@qparts@begin[#1]\fi}{\ifri@sol\par\else\ri@qparts@end\fi}
\def\ri@qpartsname{qparts}
\renewcommand{\pt}[1]{\ifnum\ri@depth<2 \gdef\ri@outer{#1}\gdef\ri@cur{#1}\else\expandafter\gdef\expandafter\ri@cur\expandafter{\ri@outer#1}\fi
  \ifri@sol\expandafter\ri@solpt\else\expandafter\ri@qpt\fi{#1}}
\def\ri@qpt#1{\expandafter\gdef\csname ri@hp@\ri@qid\endcsname{}\item[#1]}
\def\ri@solpt#1{\@ifnextchar\begin\ri@ptpeek{\ri@row{\ri@cur}}}
\def\ri@ptpeek\begin#1{\def\ri@tmp{#1}\ifx\ri@tmp\ri@qpartsname
   \expandafter\@firstoftwo\else\expandafter\@secondoftwo\fi
  {\begin{qparts}}{\ri@row{\ri@cur}\begin{#1}}}
% marks: record [n] against the current question/part
\let\ri@mk\mk
\renewcommand{\mk}[1]{\ri@setkey\ri@k{\ri@qid\ri@cur}\ri@setkey\ri@mv{#1}\expandafter\xdef\csname ri@mk@\ri@k\endcsname{\ri@mv}\ri@mk{#1}}
% \slabel (Ch10 Test answers): letter labels set the outer part, roman ones are appended;
% consecutive \slabel calls merge into one row label
\def\ri@romans{(i)(ii)(iii)(iv)(v)(vi)(vii)}
\def\ri@firstpart(#1)#2\@nil{(#1)}
\def\slabel#1{\ifri@sol\expandafter\ri@solslabel\else\expandafter\ri@qslabel\fi{#1}}
\def\ri@qslabel#1{\par\smallskip\noindent\textbf{#1}\par\nobreak}
\def\ri@solslabel#1{\@expandtwoargs\in@{#1}{\ri@romans}%
  \ifin@\expandafter\gdef\expandafter\ri@cur\expandafter{\ri@outer#1}\else\gdef\ri@cur{#1}\xdef\ri@outer{\ri@firstpart#1\@nil}\fi
  \@ifnextchar\slabel{}{\ri@row{\ri@cur}}}
% Ch11 spart{label}: a row; soln: the table
\def\spart#1{\gdef\ri@cur{#1}\ri@row{#1}}\def\endspart{\par}
\def\soln{\@ifnextchar[{\ri@solnopt}{\risol}}\def\ri@solnopt[#1]{\risol}\def\endsoln{\endrisol}
% ---- working that started on the label line is left-aligned in the working column
\AtBeginDocument{\@ifundefined{topalign}{}{%
  \renewenvironment{topalign}{\par\noindent$\displaystyle\begin{aligned}[t]}{\end{aligned}$\par\smallskip}%
  \renewcommand{\topdisplay}[1]{\par\noindent$\displaystyle #1$\par\smallskip}}}
% ---- Ch11 side-by-side (text | sketch): the sketch shrinks to fit its half of the working column
\AtBeginDocument{\@ifundefined{sbsplain}{}{%
  \renewenvironment{sbsplain}[1]{%
    \sbs@start\begin{minipage}[t]{#1\linewidth}\vspace{0pt}\sbs@shift\raggedright%
    \def\sbsnext{\end{minipage}\hfill\begin{minipage}[t]{\dimexpr0.97\linewidth-#1\linewidth\relax}\vspace{0pt}\sbs@shift\centering
      \begin{adjustbox}{max width=\linewidth}}%
  }{\end{adjustbox}\end{minipage}\par}}}
% ---- figures in the working column shrink to fit it (only if wider), colours kept
\newcommand{\rifitfig}[1]{\expandafter\let\csname ri@orig@\string#1\endcsname#1%
  \def#1{\@ifnextchar\bgroup{\ri@fitarg#1}{\ri@fitnoarg#1}}}
\def\ri@fitarg#1#2{\ifri@sol\adjustbox{max width=\linewidth}{\csname ri@orig@\string#1\endcsname{#2}}\else\csname ri@orig@\string#1\endcsname{#2}\fi}
\def\ri@fitnoarg#1{\ifri@sol\adjustbox{max width=\linewidth}{\csname ri@orig@\string#1\endcsname}\else\csname ri@orig@\string#1\endcsname\fi}
\@ifundefined{solfig}{}{\renewcommand{\solfig}[2][]{\par\begin{center}\adjustbox{max width=\linewidth}{\begin{tikzpicture}\begin{axis}[qfig, width=6.6cm, height=5cm, #1]#2\end{axis}\end{tikzpicture}}\end{center}}}
% \ricomment{text}: text in the Comments column of the current row (call right after \pt,
% \slabel or \spart). The row is made at least as tall as the comment.
\newbox\ri@cmtbox
\newcommand{\ricomment}[1]{\par\setbox\ri@cmtbox=\vtop{\hsize=\dimexpr\ri@C-2\ri@pad\relax
    \parskip=0pt\raggedright\small{\color{black!60}\strut#1\strut\par}}%
  \ifdim\dimexpr\ht\ri@cmtbox+\dp\ri@cmtbox\relax>\ri@curH sp\xdef\ri@curH{\number\dimexpr\ht\ri@cmtbox+\dp\ri@cmtbox\relax}\fi
  \nointerlineskip\vbox to 0pt{\hbox to\linewidth{\hskip\dimexpr\linewidth+2\ri@pad\relax\box\ri@cmtbox\hss}\vss}\nobreak\prevdepth=\dp\strutbox
  \vskip-\parskip\ignorespaces}
% \putTFlong{n}: the True/False question n with its answer table allowed to break across
% pages (the body of \defQ{n} must be: intro text, then \begin{tftable}...\end{tftable});
% the solutions copy prints it with tflongtable.
\renewcommand{\putTFlong}[1]{\ri@newq\begingroup\let\tftable\tflongtable\let\endtftable\endtflongtable
  \stepcounter{diq}\expandafter\xdef\csname diLab@#1\endcsname{\thediq}%
  \par\noindent\hbox to 2.6em{\textbf{\thediq}\hfil}\setbox0=\hbox{\di@tag{#1}}\ifdim\wd0=0pt\else{\small\textcolor{black!60}{[\di@tag{#1}]}}\par\noindent\hskip2.6em\fi\ignorespaces\csname diQbody@#1\endcsname\par
  \endgroup\vspace{3ex}}
% ---- question + solution; each question starts on a new page (as RI)
\def\ri@qid{}
\let\ri@putQ\putQ
\renewcommand{\putQ}[1]{\gdef\ri@qid{#1}\global\ri@depth=0 \gdef\ri@cur{}\ri@putQ{#1}}
% \risolQ{n}{body macro name}: question n + its solution body in the table;
% \risolQraw: body already contains its own soln environment (Ch11 Test)
\newcommand{\ri@newq}{\ifri@fresh\global\ri@freshfalse\else\clearpage\fi}
% keep the whole table together if it fits on a page; otherwise make sure at least its
% header and first row fit before starting it
\newcommand{\ri@keep}[1]{\par
  \ifcsname ri@total@#1\endcsname
    % hysteresis: a table kept whole in the previous run stays whole up to textheight-4.5mm,
    % so a total near the limit cannot flip between runs
    \ifdim\csname ri@total@#1\endcsname<\dimexpr\textheight-\ifcsname ri@kept@#1\endcsname4.5\else8\fi mm\relax
      \protected@write\@auxout{}{\string\ri@setkept{#1}}%
      \needspace{\dimexpr\csname ri@total@#1\endcsname+4mm\relax}%
    \else\ifcsname ri@first@#1\endcsname
      \ri@rowh{\csname ri@first@#1\endcsname}%
      \ifdim\ri@h<\dimexpr\textheight-20mm\relax\needspace{\dimexpr\ri@h+18mm\relax}\fi\fi\fi\fi}
\newcommand{\risolQ}[2]{\ri@newq\putQ{#1}\vspace{-3ex}\ri@keep{#1}\begin{risol}\csname #2@#1\endcsname\par\end{risol}}
\newcommand{\risolQraw}[2]{\ri@newq\putQ{#1}\vspace{-3ex}\ri@keep{#1}\csname #2@#1\endcsname\par}
\renewcommand{\putQS}[1]{\risolQ{#1}{diS}}
% end marker only if it fits on the last page (no page holding just the marker)
\renewcommand{\diend}[1][END OF SOLUTIONS]{\par\ifdim\dimexpr\pagegoal-\pagetotal\relax<4\baselineskip\else\vspace{1.5ex}\begin{center}\textbf{#1}\end{center}\fi}
\makeatother
\tfanswerstrue\fi
\fi

\usetikzlibrary{calc}
\providecommand{\cv}[3]{\begin{pmatrix}#1\\#2\\#3\end{pmatrix}}
\providecommand{\vv}[1]{\overrightarrow{#1}}
\housechapter{Vectors I}
\ditle{Homework 1 (Introduction to Vectors in Three Dimensions)}
\defQ{1}{}{%
The vectors $\mathbf a$ and $\mathbf b$ are given by $\mathbf a=2\mathbf i-\mathbf j+3\mathbf k$ and $\mathbf b=-\mathbf i+4\mathbf j$.
\begin{qparts}
\pt{(a)} Write $\mathbf a$ and $\mathbf b$ as column vectors and find $\mathbf a+2\mathbf b$. \mk{2}
\pt{(b)} Find $\mathbf c$ such that $2\mathbf a-3\mathbf b+\mathbf c=\mathbf0$. \mk{2}
\pt{(c)} Points $P$ and $Q$ have position vectors $\mathbf a+2\mathbf b$ and $\mathbf c$ respectively. Find $\vv{PQ}$. \mk{1}
\end{qparts}}

\defS{1}{%
\begin{qparts}
\pt{(a)} $\displaystyle\mathbf a=\cv{2}{-1}{3},\qquad\mathbf b=\cv{-1}{4}{0},$
\begin{align*}
\mathbf a+2\mathbf b&=\cv{2}{-1}{3}+2\cv{-1}{4}{0}\\
&=\cv{0}{7}{3}.
\end{align*}
\pt{(b)} $\displaystyle\mathbf c=3\mathbf b-2\mathbf a=\cv{-3}{12}{0}-\cv{4}{-2}{6}=\cv{-7}{14}{-6}$.
\pt{(c)} $\displaystyle\vv{PQ}=\vv{OQ}-\vv{OP}=\cv{-7}{14}{-6}-\cv{0}{7}{3}=\cv{-7}{7}{-9}$.
\end{qparts}}

\defQ{2}{}{%
The points $A$ and $B$ have coordinates $(1,-2,2)$ and $(3,1,-4)$ respectively.
\begin{qparts}
\pt{(a)} Find a unit vector that is parallel to $AB$. \mk{2}
\pt{(b)} Find a vector of length 21 parallel to $AB$. \mk{2}
\pt{(c)} A point $P$ is such that $\vv{AP}$ has magnitude $14$ and the same direction as $\vv{AB}$. Find the coordinates of $P$ and state the order of the three points along their common straight path. \mk{2}
\end{qparts}}

\defS{2}{%
\begin{qparts}
\pt{(a)} \[\vv{AB}=\cv{3}{1}{-4}-\cv{1}{-2}{2}=\cv{2}{3}{-6},\qquad |\vv{AB}|=\sqrt{4+9+36}=7.\]
One such unit vector is $\displaystyle\frac17\cv{2}{3}{-6}$. Its negative is also a valid answer.
\pt{(b)} One such vector is $\displaystyle21\left(\frac17\cv{2}{3}{-6}\right)=\boxed{\cv{6}{9}{-18}}$. Its negative is also a valid answer.
\pt{(c)} \[\vv{AP}=14\left(\frac17\cv{2}{3}{-6}\right)=\cv{4}{6}{-12},\qquad \vv{OP}=\vv{OA}+\vv{AP}=\cv{5}{4}{-10}.\]
Thus $\boxed{P(5,4,-10)}$. Since $\vv{AP}=2\vv{AB}$, the order is $A,B,P$.
\end{qparts}}

\defQ{3}{}{%
The vectors $\mathbf p=\cv{6}{k}{-3}$ and $\mathbf q=\cv{2}{-4}{m}$ are parallel, where $k$ and $m$ are real constants.
\begin{qparts}
\pt{(a)} Find $k$ and $m$. \mk{3}
\pt{(b)} State whether $\mathbf p$ and $\mathbf q$ have the same or opposite directions, and find $|\mathbf p|:|\mathbf q|$. \mk{1}
\pt{(c)} Find all vectors of magnitude $\sqrt{21}$ that are parallel to $\mathbf p$. \mk{2}
\end{qparts}}

\defS{3}{%
\begin{qparts}
\pt{(a)} Since the vectors are parallel, $\mathbf p=h\mathbf q$ for some non-zero real $h$. Comparing first components gives $6=2h$, so $h=3$.
\[k=3(-4)=-12,\qquad -3=3m\quad\Longrightarrow\quad\boxed{k=-12,\ m=-1}.\]
\pt{(b)} Since $\mathbf p=3\mathbf q$, they have the same direction and $\boxed{|\mathbf p|:|\mathbf q|=3:1}$.
\pt{(c)} $\displaystyle|\mathbf q|=\sqrt{2^2+(-4)^2+(-1)^2}=\sqrt{21}$. Both directions must be included, so all such vectors are
\[\boxed{\pm\cv{2}{-4}{-1}}.\]
\end{qparts}}

\defQ{4}{}{%
The points $A$, $B$ and $C$ have coordinates $(1,2,-1)$, $(3,1,2)$ and $(7,-1,8)$ respectively.
\begin{qparts}
\pt{(a)} Show that $A$, $B$ and $C$ are collinear, and find $AB:BC$. \mk{3}
\pt{(b)} The point $R(u,-2,v)$ lies on the same straight path as $A$, $B$ and $C$. Find $u$ and $v$. \mk{3}
\end{qparts}}

\defS{4}{%
\begin{qparts}
\pt{(a)} \[\vv{AB}=\cv{2}{-1}{3},\qquad \vv{BC}=\cv{4}{-2}{6}=2\vv{AB}.\]
Thus $\vv{AB}$ and $\vv{BC}$ are parallel and have the common point $B$, so $A,B,C$ are collinear. Taking magnitudes gives $\boxed{AB:BC=1:2}$.
\pt{(b)} Collinearity gives $\vv{AR}=s\vv{AB}$ for some real $s$, so
\[\cv{u-1}{-4}{v+1}=s\cv{2}{-1}{3}.\]
The second components give $s=4$. Hence $u-1=8$ and $v+1=12$, giving $\boxed{u=9,\ v=11}$.
\end{qparts}}

\defQ{5}{}{%
The points $A$ and $B$ have coordinates $(2,-1,3)$ and $(-4,5,0)$ respectively. A point $M$ divides $AB$ internally in the ratio $AM:MB=2:1$. A point $N$ lies on $AB$ produced beyond $B$ and satisfies $AN:NB=5:2$.
\begin{qparts}
\pt{(a)} Find the position vector of $M$. \mk{2}
\pt{(b)} Find the position vector of $N$. \mk{3}
\pt{(c)} State the order of the four points $A$, $B$, $M$ and $N$, and find $MB:BN$. \mk{3}
\end{qparts}}

\defS{5}{%
\begin{qparts}
\pt{(a)} By the ratio theorem,
\[\vv{OM}=\frac{\vv{OA}+2\vv{OB}}3=\frac13\left[\cv{2}{-1}{3}+2\cv{-4}{5}{0}\right]=\boxed{\cv{-2}{3}{1}}.\]
\pt{(b)} Since $AN=AB+BN$ and $AN:BN=5:2$, $AB:BN=3:2$. Hence $\vv{BN}=\frac23\vv{AB}$, and
\[\vv{ON}=\vv{OB}+\frac23(\vv{OB}-\vv{OA})=\frac{5\vv{OB}-2\vv{OA}}3=\boxed{\cv{-8}{9}{-2}}.\]
\pt{(c)} The order is $A,M,B,N$. Since $MB=\frac13 AB$ and $BN=\frac23 AB$,
\[\boxed{MB:BN=1:2}.\]
\end{qparts}}

\defQ{6}{}{%
Relative to the origin $O$, the points $A$ and $B$ have non-zero, non-parallel position vectors $\mathbf a$ and $\mathbf b$ respectively. The point $X$ lies on $OA$ with $OX:XA=2:1$, and the point $Y$ lies on $OB$ with $OY:YB=1:2$. The line segments $AY$ and $BX$ intersect at $P$.
\begin{center}\begin{tikzpicture}[scale=0.55,every node/.style={font=\small}]
\coordinate(O)at(0,0);\coordinate(A)at(0,4.2);\coordinate(B)at(7,0);
\coordinate(X)at(0,2.8);\coordinate(Y)at(2.3333,0);\coordinate(P)at(1,2.4);
\draw[thick](O)--(A)--(B)--cycle;\draw(A)--(Y);\draw(B)--(X);
\foreach \n in {O,A,B,X,Y,P}{\fill (\n) circle(1.4pt);}
\node[left]at(O){$O$};\node[above]at(A){$A$};\node[right]at(B){$B$};
\node[left]at(X){$X$};\node[below]at(Y){$Y$};\node[above right]at(P){$P$};
\end{tikzpicture}\end{center}
\begin{qparts}
\pt{(a)} By expressing $\vv{OP}$ in two ways using the ratio theorem, find $\vv{OP}$ in terms of $\mathbf a$ and $\mathbf b$. \mk{6}
\pt{(b)} Find $AP:PY$ and $BP:PX$. \mk{2}
\pt{(c)} The point $Z$ divides $AB$ internally with $AZ:ZB=1:4$. Show that $O,P,Z$ are collinear and find $OP:PZ$. \mk{2}
\end{qparts}}

\defS{6}{%
\begin{qparts}
\pt{(a)} First, $\vv{OX}=\frac23\mathbf a$ and $\vv{OY}=\frac13\mathbf b$.
Let $AP:PY=s:(1-s)$ and $BP:PX=t:(1-t)$, where $0<s,t<1$. By the ratio theorem,
\begin{align*}
\vv{OP}&=(1-s)\vv{OA}+s\vv{OY}=(1-s)\mathbf a+\frac{s}{3}\mathbf b,\\
\vv{OP}&=(1-t)\vv{OB}+t\vv{OX}=\frac{2t}{3}\mathbf a+(1-t)\mathbf b.
\end{align*}
Since $\mathbf a$ and $\mathbf b$ are non-parallel, comparing coefficients gives
\[1-s=\frac{2t}{3},\qquad \frac{s}{3}=1-t.\]
The second equation gives $s=3-3t$. Substituting into the first,
\[-2+3t=\frac{2t}{3}\quad\Longrightarrow\quad 7t=6\quad\Longrightarrow\quad t=\frac67,\quad s=\frac37.\]
Consequently,
\[\boxed{\vv{OP}=\frac47\mathbf a+\frac17\mathbf b}.\]
\pt{(b)} $\displaystyle\boxed{AP:PY=3:4,\qquad BP:PX=6:1}$.
\pt{(c)} By the ratio theorem,
\[\vv{OZ}=\frac{4\mathbf a+\mathbf b}{5},\qquad \vv{OP}=\frac{4\mathbf a+\mathbf b}{7}=\frac57\vv{OZ}.\]
The two vectors are parallel and have the common point $O$, so $O,P,Z$ are collinear. Since $0<\frac57<1$, $P$ is between $O$ and $Z$; hence $\boxed{OP:PZ=5:2}$.
\end{qparts}}

\defQ{7}{}{%
A solid house consists of a cuboid $OBCDEFGH$ and a triangular prism roof $EFPHGQ$, as shown. The base $OBCD$ is horizontal. The vertices $E$, $F$, $G$ and $H$ are directly above $O$, $B$, $C$ and $D$ respectively.

The front-left corner $O$ is the origin. The positive $\mathbf i$, $\mathbf j$ and $\mathbf k$ directions are along $OB$, $OD$ and $OE$ respectively, and each unit represents $1$ metre. It is given that $OB=6\,\mathrm m$, $OD=4\,\mathrm m$ and $OE=3\,\mathrm m$.

The ridge $PQ$ is parallel to $OD$. The front gable $EFP$ and back gable $HGQ$ are congruent isosceles triangles, with $P$ vertically above the midpoint of $EF$ and $Q$ vertically above the midpoint of $HG$. Each sloping roof face makes an angle $\alpha$ with the horizontal, where $\tan\alpha=\frac23$; the vertical roof rise is $2\,\mathrm m$.
\begin{center}
\begin{tikzpicture}[
  x={(0.85cm,0cm)},
  y={(0.29cm,0.32cm)},
  z={(0cm,0.85cm)},
  line join=round, line cap=round, >=stealth,
  every node/.style={font=\small}
]
\coordinate(O)at(0,0,0);
\coordinate(B)at(6,0,0);
\coordinate(C)at(6,4,0);
\coordinate(D)at(0,4,0);
\coordinate(E)at(0,0,3);
\coordinate(F)at(6,0,3);
\coordinate(G)at(6,4,3);
\coordinate(H)at(0,4,3);
\coordinate(P)at(3,0,5);
\coordinate(Q)at(3,4,5);
\coordinate(M)at(3,0,3);
% hidden edges
\draw[dashed,thick](O)--(D)--(C);
\draw[dashed,thick](D)--(H)--(G);
% visible edges
\draw[thick](O)--(B)--(F)--(E)--cycle;
\draw[thick](B)--(C)--(G)--(F);
\draw[thick](E)--(P)--(F);
\draw[thick](P)--(Q)--(G);
\draw[thick](E)--(H)--(Q);
% roof rise and angle
\draw[dashed](M)--(P);
\draw ($(M)+(-0.25,0,0)$)--++(0,0,0.25)--++(0.25,0,0);
\node[right=1pt] at ($(M)!0.5!(P)$) {$2\,\mathrm{m}$};
\draw ($(F)+(-0.85,0,0)$) arc[start angle=180, end angle=146.31, radius=0.85*0.85cm];
\node at ($(F)+(-1.30,0,0.30)$) {$\alpha$};
% lengths
\node[below=3pt] at ($(O)!0.62!(B)$) {$6\,\mathrm{m}$};
\node[below right=1pt] at ($(O)!0.62!(D)$) {$4\,\mathrm{m}$};
\node[left=2pt] at ($(O)!0.68!(E)$) {$3\,\mathrm{m}$};
% axes
\draw[->,very thick](O)--(1.3,0,0);
\node[below=2pt] at (0.95,0,0) {$\mathbf{i}$};
\draw[->,very thick](O)--(0,1.5,0);
\node[above left=-1pt] at (0,1.15,0) {$\mathbf{j}$};
\draw[->,very thick](O)--(0,0,1.2);
\node[left=2pt] at (0,0,0.9) {$\mathbf{k}$};
\foreach \n in {O,B,C,D,E,F,G,H,P,Q} {\fill(\n) circle (1.4pt);}
\node[below left] at (O) {$O$};
\node[below] at (B) {$B$};
\node[right] at (C) {$C$};
\node[below right=-1pt] at (D) {$D$};
\node[left] at (E) {$E$};
\node[below right=-1pt] at (F) {$F$};
\node[right] at (G) {$G$};
\node[left] at (H) {$H$};
\node[above left=-1pt] at (P) {$P$};
\node[above right=-1pt] at (Q) {$Q$};
\end{tikzpicture}\par
{\small All lengths are in metres. The diagram is not drawn to scale.}
\end{center}Find the coordinates of all the vertices of the house $(O,B,C,D,\ldots)$. \mk{6}}

\defS{7}{%
The positive coordinate directions are along $OB$, $OD$ and $OE$. Hence the four base vertices have $z=0$, the front vertices have $y=0$, and the back vertices have $y=4$.

The walls are $3\,\mathrm m$ high, so $E,F,G,H$ all have $z=3$. The roof is symmetric across the middle of its $6\,\mathrm m$ width, so the ridge has $x=\frac62=3$.

Writing $h$ for the vertical roof rise,
\[\tan\alpha=\frac{h}{6/2}=\frac23\quad\Longrightarrow\quad h=2\,\mathrm m.\]
This agrees with the labelled rise. The ridge is at height $3+2=5\,\mathrm m$. Therefore all ten vertices are
\[\boxed{\begin{aligned}
O&=(0,0,0),& B&=(6,0,0),\\
C&=(6,4,0),& D&=(0,4,0),\\
E&=(0,0,3),& F&=(6,0,3),\\
G&=(6,4,3),& H&=(0,4,3),\\
P&=(3,0,5),& Q&=(3,4,5).
\end{aligned}}\]}

\defQ{8}{}{%
The diagram shows a container with a horizontal rectangular base $OABC$, where $OA=16\,\mathrm{cm}$ and $AB=11\,\mathrm{cm}$. The top of the container $DEFG$ is also a horizontal rectangle, where $DE=10\,\mathrm{cm}$ and $EF=5\,\mathrm{cm}$. The four sloping faces (for example $OAED$) of the container, each a trapezium, are inclined at the same angle to the horizontal, and the vertical distance between the base and the top is $6\,\mathrm{cm}$.

The point $O$ is taken as the origin $(0,0,0)$, and perpendicular unit vectors $\mathbf i$, $\mathbf j$, $\mathbf k$ are such that $\mathbf i$ is parallel to $OA$, $\mathbf j$ is parallel to $OC$, and $\mathbf k$ is vertical.
\begin{center}
\begin{tikzpicture}[
  x={(0.48cm,-0.03cm)},
  y={(0.32cm,0.22cm)},
  z={(0cm,0.48cm)},
  line join=round, line cap=round, >=stealth,
  every node/.style={font=\bfseries\small}
]
\coordinate(O)at(0,0,0);
\coordinate(A)at(16,0,0);
\coordinate(B)at(16,11,0);
\coordinate(C)at(0,11,0);
\coordinate(D)at(3,3,6);
\coordinate(E)at(13,3,6);
\coordinate(F)at(13,8,6);
\coordinate(G)at(3,8,6);
\draw[dashed,thick](O)--(C)--(B);
\draw[dashed,thick](C)--(G);
\draw[thick](O)--(A)--(B);
\draw[thick](O)--(A)--(E)--(D)--cycle;
\draw[thick](A)--(B)--(F)--(E)--cycle;
\draw[thick](D)--(E)--(F)--(G)--cycle;
\draw[thick](O)--(D);
\draw[thick](G)--(D);
\draw[->,thick](O)--(2.8,0,0) node[below] {$\mathbf{i}$};
\draw[->,thick](O)--(0,2.4,0) node[above right=-1pt] {$\mathbf{j}$};
\draw[->,thick](O)--(0,0,2.6) node[left] {$\mathbf{k}$};
\foreach \n in {O,A,B,C,D,E,F,G} {\fill(\n) circle (1.45pt);}
\node[below left] at (O) {$O$};
\node[below] at (A) {$A$};
\node[right] at (B) {$B$};
\node[left] at (C) {$C$};
\node[above left] at (D) {$D$};
\node[above right] at (E) {$E$};
\node[right] at (F) {$F$};
\node[above] at (G) {$G$};
\node[below=2pt] at ($(O)!0.5!(A)$) {$16$};
\node[below right=-1pt] at ($(A)!0.5!(B)$) {$11$};
\node[above=1pt] at ($(D)!0.5!(E)$) {$10$};
\node[right=1pt] at ($(E)!0.5!(F)$) {$5$};
\end{tikzpicture}\par
{\small All lengths are in centimetres. The diagram is not drawn to scale.}
\end{center}Find the coordinates of $D$. \mk{3}}

\defS{8}{%
Since the four sloping faces are inclined at the same angle to the horizontal, the top rectangle $DEFG$ is centred above the base rectangle $OABC$.

The horizontal inset from each base edge to the corresponding top edge is
\[\frac{OA-DE}{2}=\frac{16-10}{2}=3,\qquad \frac{AB-EF}{2}=\frac{11-5}{2}=3.\]
The top lies $6\,\mathrm{cm}$ above the base, so the $z$-coordinate of every top vertex is $6$. Taking $D$ as the front-left vertex of the top,
\[\boxed{D=(3,3,6)}.\]}

\ifdefined\diquestionsonly\expandafter\endinput\fi
\begin{document}
\dititleblock[Original practice questions. Total: 50 marks. Give exact answers unless otherwise stated.]
\disection{Questions}
\putQ{1}\putQ{2}\putQ{3}\putQ{4}\putQ{5}\putQ{6}\putQ{7}\putQ{8}
\diend[END OF HOMEWORK]
\end{document}
